\documentclass[11pt,a4paper]{article}
\usepackage{amsmath}
\usepackage{fontspec}
\usepackage[mathrm=sym]{unicode-math}
\setmainfont{NotoSans-Regular.ttf}[Path=../fonts/,BoldFont=NotoSans-Bold.ttf,ItalicFont=NotoSans-Italic.ttf,BoldItalicFont=NotoSans-BoldItalic.ttf]
\setsansfont{NotoSans-Regular.ttf}[Path=../fonts/,BoldFont=NotoSans-Bold.ttf,ItalicFont=NotoSans-Italic.ttf,BoldItalicFont=NotoSans-BoldItalic.ttf]
\setmonofont{NotoSansMono-Regular.ttf}[Path=../fonts/,BoldFont=NotoSansMono-Bold.ttf]
\setmathfont{FiraMath-Regular.otf}[Path=../fonts/]
\usepackage[margin=26.5mm,includefoot]{geometry}
\usepackage{microtype}
\usepackage{graphicx}
\usepackage{xurl}
\usepackage[numbers,sort&compress]{natbib}
\usepackage[colorlinks=true,urlcolor=blue,linkcolor=blue,citecolor=blue]{hyperref}
\setlength{\parindent}{0pt}
\setlength{\parskip}{6pt}
\emergencystretch=3em
\hypersetup{pdftitle={Brain2Body: The Next Generation In Body-Brain Interfacing.},pdfauthor={@eliminatenoise},pdfsubject={Conceptual working draft; peer review pending; biological teaching untested; current model gates fail}}
% R9 retains the sans-serif master and the historical R6/R8 Source provenance.
\usepackage{ragged2e}
\usepackage[justification=raggedright,singlelinecheck=false]{caption}
\setlength{\RaggedRightRightskip}{0pt plus 1fil}
\AtBeginDocument{\RaggedRight}
% Match the already qualified reader paragraph-flow controls without changing type size.
\makeatletter
\AtBeginDocument{\clubpenalty=10000\@clubpenalty=10000\widowpenalty=10000\displaywidowpenalty=10000}
\makeatother
\renewcommand{\bibfont}{\RaggedRight}
\AddToHook{cmd/LARGE/after}{\hyphenpenalty=10000\exhyphenpenalty=10000}
\begin{document}
{\sffamily\LARGE\bfseries Brain2Body: The Next Generation In Body-Brain Interfacing.\par}\vspace{8pt}
{\sffamily @eliminatenoise\par}\vspace{8pt}
{\sffamily Working Draft — Review Pending\par}\vspace{8pt}


\begin{abstract}
An assisted hand can perform a movement without its nervous system having
learned to perform it. Brain2Body asks what remains when teaching commands end,
with any compensatory movement bridge held constant. The proposal connects stated nervous-system and body
observations to independently measured outcomes, then separates assistance,
simulated control, retained learning and biological repair. Published
motor-control research motivates the experiments; the companion computational
studies provide bounded comparisons rather than biological teaching evidence.
In the current LINK study, 18 runs and 720 saved epochs fail both model-admission
gates and all five benefit claims. The proposed next experiment pairs EEG, EMG
and calibrated movement with a musculoskeletal physics engine and a learning
model. Forward movement prediction and reverse recorded-feature reconstruction
are separate tests; neither identifies a unique brain state. Sensory delivery
requires a separately qualified writer. The comparison must separate what the
model predicts, what the controller executes and what the trainee retains.
It reports no new biological
experiment or retained skill acquisition.
A separate self-forecast diagnostic used one original TRAIN acquisition with 64 surface-EMG RMS features and 18 glove angles. Its 40-update pilot and once-completed continuation to 990 cumulative updates shared one GRU core across separate same-stream forecasting tasks. Final 990 improved over initialization and 40 updates but lost to past-value persistence in both pooled streams. No DEV/TEST generalization, neural-to-body mapping or biological learning follows from that within-acquisition comparison.
\end{abstract}
% Derived from the owned full Markdown narrative; source-stage only.


\section{Assistance and Retained Learning}

Let a computer guide a hand through a tune, then stop the teaching commands.
What can the trainee play now? The first performance shows what the combined
system could do; the second begins to test what the nervous system retained.
Any movement bridge needed to express the trainee's own control must remain
available in both comparisons. Otherwise we would be testing a disconnected
hand and calling the result a test of memory. Brain2Body connects
nervous-system measurements to physical outcomes to make this distinction
testable. The numeral 2 means ``to''; body-to-brain feedback matters as much
as the outward command.\par

Rapid acquisition of difficult physical skills is the ambition. Sewing, playing an instrument, manipulating a puzzle and coordinating a whole body put that ambition under pressure: each demands more than one correct movement. The programme must explain how useful guidance could become a capability the trainee retains. These are proposed tests, not demonstrated abilities. Neither arbitrary skill upload nor complete nervous-system reproduction is established, and the current models have not passed their original performance gates.\par

A decoder turns an observation into a prediction. Before calling it correct, we must decide what answer it is meant to give. NeuroTranslation proposes relating neural signals and their representations between explicitly described sources and targets; Brain2Body follows that prediction into a body, where movement changes the feedback available for the next action and learning becomes a further question. The argument therefore needs an independently measured action, a model whose contribution can be isolated, and a retention test that can fail. Published research motivates those requirements. The companion computational work supplies bounded comparisons, including negative results; it has not supplied evidence of biological teaching.\par

\section{Compare Predicted and Measured Movement}

Perception is subjective. Movement is not. A person experiences a sensation; an independent instrument can measure a finger's position, a contact, a force or the timing of a release. That contrast does not make every movement measurement exact or reveal private experience. It gives the programme a tractable starting point: relate an uncertain neural observation to an outcome that can be checked.\par

Consider a key press. Choosing the right key, reaching it, pressing at the right time, releasing it and avoiding neighbouring keys are separate outcomes. A video can show position; a contact sensor can show a strike; a force sensor can measure force. A task score records success under declared rules. Neural activity near an electrode is another observation. None silently supplies the others.\par

In the proposed loop, sensory input guides action. The body and environment
change; measured outcomes supply feedback for the next action. Feedback control
adjusts action using the difference between the desired and observed outcome.
Todorov and Jordan provide an established account of motor coordination under
uncertainty. Applying it here is a programme rationale, not an empirical result
of their paper. Every claimed loop must disclose its delays and whether its
feedback came from a living subject, a simulation or a recorded replay.
\href{https://www.nature.com/articles/nn963}{Todorov and Jordan, 2002}.\par

Accuracy means agreement with an independent reference. Here we assess precision through repeatability under stated conditions. A consistently wrong model can be precise. Ground truth means an independently acquired measured reference with its own calibration, sampling limits and uncertainty; it is not error-free access to biology. Finer timesteps, more decimal places and a confident animation cannot establish these properties. \href{https://www.nist.gov/pml/nist-technical-note-1297/nist-tn-1297-appendix-d1-terminology}{NIST measurement terminology}.\par

\section{Structure, Activity and Function}

A movement score tells us what happened. Explaining why requires more than a map of the routes a signal could take. Timing, the cells' responses to chemical signals, the body's current state, learning history and task conditions affect what follows. An intervention deliberately changes something to test its effect; a connection map alone cannot tell us that response. Structure narrows the explanation, while measured function must test it. That is why a body-grounded prediction needs both.\par

The Human Genome Project offers a bounded analogy. Obtaining a reference sequence did not finish the investigation of regulation and function; subsequent projects investigated functional elements. That is a distinction between a map and its interpretation, not a claim that genome sequencing failed. \href{https://www.genome.gov/11009066/2003-release-scientists-venture-deeper-into-the-human-genome}{NHGRI, Beyond Genes}.\par

The MICrONS Consortium combined recorded visual responses with detailed connectivity in a volume of mouse visual cortex. Its strength is the combination of different measurements. It does not provide a whole human nervous system, peripheral or spinal signals, or an unrestricted account of learned skills. Co-registration means relating measurements to corresponding locations; such correspondence must be demonstrated for the actual source, not borrowed from an attractive template. \href{https://www.nature.com/articles/s41586-025-08790-w}{MICrONS Consortium, 2025}.\par

Measurement itself needs explanation. Functional magnetic resonance imaging, or fMRI, measures an indirect blood-related signal over time. In the experiments of Logothetis and colleagues, that signal related more closely to local field potentials, aggregate electrical activity near a recording site, than to spiking output. It is not a photograph of thought. Tang and colleagues demonstrated constrained semantic reconstruction with substantial dependence on the subject and cooperation; this is not universal one-to-one decoding of subjective experience. \href{https://pubmed.ncbi.nlm.nih.gov/11449264/}{Logothetis et al., 2001}, \href{https://www.nature.com/articles/s41593-023-01304-9}{Tang et al., 2023}.\par

Quantum measurement can illustrate that observing a system may involve interaction. It does not establish quantum backaction, the influence of a quantum measurement on a system, as the cause of these large-scale neural limits. Scanner constraints, indirect signals, person-specific learning and ambiguous representations supply concrete explanations that must be tested on their own terms. The analogy stops before it becomes an unsupported brain mechanism. \href{https://www.nobelprize.org/uploads/2018/06/advanced-physicsprize2012-1.pdf}{Royal Swedish Academy, quantum measurement background}.\par

\section{Existing Brain--Body Interfaces}

Bouton and colleagues linked recordings from motor cortex, movement-related tissue in the brain's outer layer, to electrical activation of forearm muscles, enabling useful hand and wrist movements in one participant with quadriplegia, paralysis affecting four limbs. This demonstrates a participant-specific route from brain signals to bodily action. It does not demonstrate an unfamiliar skill transplant or a complete map of the nervous system. \href{https://www.nature.com/articles/nature17435}{Bouton et al., 2016}.\par

Lorach and colleagues linked cortical recording to spinal stimulation and evaluated standing and walking in a participant with chronic spinal injury. They also reported rehabilitation-associated improvement with the interface switched off. Movement in a real participant provides functional evidence beyond a rendered demonstration. The study remains a particular participant, intervention and rehabilitation history; assistance, retained improvement and tissue regeneration are distinct claims. \href{https://www.nature.com/articles/s41586-023-06094-5}{Lorach et al., 2023}.\par

Flesher and colleagues used stimulation of sensory brain tissue to provide touch-related feedback during robotic arm control. The work supports testing the return path of an interface. It does not establish transfer of an expert's repertoire. These studies motivate bounded experiments; they cannot be stacked into proof of a new integrated system whose relevant links have never been measured together. \href{https://pubmed.ncbi.nlm.nih.gov/34016775/}{Flesher et al., 2021}.\par

\section{Digital-Fly Models}

Digital flies sharpen the question of attribution: when a composite system
moves well, which component supplied the ability? The answer determines
whether the result tests neural wiring, a trained controller, body mechanics
or their engineered combination. Those contributions must remain visible
before a movement is credited to a reproduced nervous system.\par

FlyWire reconstructed 139,255 neurons, nerve cells, in one adult female fly brain. This is an exceptional structural resource with reconstruction and coverage limits. A map of chemical connections does not supply every connection's physiological effect, the complete nervous system, the body or the environment. Pospisil and colleagues propose combining intervention measurements with structural constraints to estimate causal effects; their estimator tests in simulation are not a completed empirical map of all fly-brain dynamics. \href{https://www.nature.com/articles/s41586-024-07558-y}{Dorkenwald et al., 2024}, \href{https://www.nature.com/articles/s41586-024-07982-0}{Pospisil et al., 2024}.\par

Shiu and colleagues show the positive route: specific connectome-derived predictions were tested through neural stimulation and behaviour. A connectome is a map of neural connections; it supplied constraints, not an automatic functional certificate. Useful predictions coexist with missed experimentally observed Phantom and Usnea effects and incorrect G2N-1/Usnea water-silencing predictions. The authors report 91\% agreement across 164 tested predictions, or 84\% after excluding one mostly negative screen. Neither figure measures whole-fly fidelity. Missing receptor dynamics, electrical connections and chemical modulation limit the account, including absolute simulated firing rates. \href{https://pmc.ncbi.nlm.nih.gov/articles/PMC11446845/}{Shiu et al., 2024}.\par

Body mechanics introduce another layer. NeuroMechFly v2 documents simulated leg tips slightly below the ground because of numerical errors or physics instability, with a 0.05 mm margin to prevent false overstretch detection. Its adhesion model uses engineered forces and stepping phases while biological liftoff mechanisms remained uncertain. These are specific disclosed modeling limitations, not evidence of dramatic tunnelling or an ability no living fly could ever acquire. FlyBody uses trained artificial controllers to imitate measured movement; its time-averaged aerodynamic approximation and omitted detailed muscle and sensory physiology remain separate limits. \href{https://gizemozd.github.io/assets/pdf/2024_neuromechflyv2.pdf}{Wang-Chen et al., 2024}, \href{https://pmc.ncbi.nlm.nih.gov/articles/PMC12310536/}{Vaxenburg et al., 2025}.\par

Eon's March 2026 technical account describes hand-chosen brain--body mappings, substantial reliance on existing body controllers, largely missing state and learning mechanisms, and internal dynamics not yet validated against relevant biological signatures. That dated developer disclosure supports a conclusion about an engineered integration, not a complete biological upload. We have not independently reproduced its dynamics. \href{https://eon.systems/updates/embodied-brain-emulation}{Eon Systems, 10 March 2026}.\par

The same attribution problem appears in games. Pinned fly-chess code supplies symbolic board information, selected mathematical dynamics, a trained readout and external legal-move search. Pinned DoomFly code uses the distinct male central-nervous-system map and injects game-engine enemy geometry into model neurons after the intended visual route reportedly failed. These observations come from inspected code and author accounts, not reproduced gameplay or evidence that a living fly learned either game. \href{https://github.com/mncrftfrcnm/making-fly-play-chess/blob/b97dfefbf9b8850f3cdf6b412c2ed840be94e260/scripts/fly_chess_trainer.py}{Fly-chess trainer}, \href{https://github.com/shreyash-sharma/doomFly/blob/5f6c2881be5f5f40b0592a75e6ee231b50333b40/src/doom/fearless/danger.py}{DoomFly threat input}, \href{https://github.com/shreyash-sharma/doomFly/blob/5f6c2881be5f5f40b0592a75e6ee231b50333b40/docs/how-it-works.md}{pinned methods and limitations}.\par

\section{Simulation Failures}

Attribution also depends on the world in which the controller succeeds. A
controller can improve against a score while moving farther from the intended
task. To use simulated practice as evidence, we need to know both what the
score rewards and whether the simulated movement obeys the relevant
physical constraints.\par

An optimizer, the procedure adjusting a model to improve its objective, can exploit an incorrect simulator or a poorly chosen score. First-hand historical accounts describe virtual creatures gaining propulsion from numerical errors and soft robots receiving artificial energy from ground-penetration corrections. CoastRunners illustrates a different failure: circling reward targets improved the score while leaving the race unfinished. Physics violation and objective misspecification need different diagnoses. Neither example establishes that every fly model shares those failures. \href{https://arxiv.org/html/1803.03453v4}{Lehman et al., first-hand accounts}, \href{https://deepmind.google/blog/specification-gaming-the-flip-side-of-ai-ingenuity/}{Google DeepMind, specification gaming}.\par

Protein prediction offers another carefully bounded comparison. AlphaFold3's original paper describes invented structure in unstructured regions and overlapping predicted chains. Its antibody-interface benchmark reports 62.9\% cluster-averaged success (reported 95\% confidence interval, 51.4--73.5\%) at DockQ > 0.23, a score threshold for agreement with experimentally determined interfaces, hence 37.1\% failing that threshold. It covers 65 interface clusters from 71 complexes, with confidence-based selection from 1,000 seeds and five samples per seed for each target. This is not a universal percentage of proteins that fail to function. Crystallization is one measurement bottleneck; other methods include cryo-electron microscopy, which images samples using electrons, and nuclear magnetic resonance, which uses atomic responses to magnetic conditions. Synthetic examples can help training, but they cannot create independent biological validation of unmeasured function. \href{https://www.nature.com/articles/s41586-024-07487-w}{Abramson et al., 2024}, \href{https://www.nature.com/articles/s41586-021-03819-2}{Jumper et al., 2021}, \href{https://pdb101.rcsb.org/learn/guide-to-understanding-pdb-data/methods-for-determining-structure}{RCSB structure methods}.\par

Biological virtualization means reproducing specified biological behaviour and causal dynamics in a digital model within independently tested conditions and error bounds. Invented biology presented as faithful reproduction is hallucination. A disclosed approximation can still be a useful scientific model. Missing validation does not prove a prediction false; one correct guess does not validate its mechanism or general reliability. The dice analogy concerns unknown reliability, not equal odds or an assertion that learned predictions are random.\par

\section{Recording Access, Context and Stable Signals}

The proposed measurement chain has to survive three questions. Access it: can we obtain the nervous-system observations relevant to a defined question, and what burden does access impose? Ground it: can we relate those observations to independently measured body behaviour on known clocks, in stated quantities and with stated uncertainty? Keep it: does the recording and its mapping remain useful as tissue, movement and signal quality change? Keeping the recording usable is a different endpoint from a trainee retaining a skill.\par

Sources need organism, subject, region, task, modality, units, sampling and acquisition conditions. Cortex, spinal cord, peripheral nerve, muscle and movement are different observables. Electroencephalography (EEG) measures electrical activity using scalp electrodes; electrospinography (ESG) records signals over the spine; electroneurography (ENG) records nerve electrical activity; electromyography (EMG) records muscle electrical activity. EMG is not direct peripheral nerve recording. Separate studies of these signals do not become one synchronized voluntary cortex--spine--periphery chain by placing their charts together.\par

The resources associated with HHMI/Janelia, Allen, Harvard and other groups contribute particular observations to particular questions. Their value depends on that match. A broad catalogue cannot supply a paired signal that was never recorded; a model can estimate the missing value, but the estimate remains a prediction. An independent measurement can check it, but must remain a separately identified observation. When data are combined across organisms, source identities and physiological differences must remain traceable. Whole-body coverage is the objective to investigate, not a property conferred by the catalogue.\par

\section{Physical Constraints}

MyoSuite supplies muscle-based simulations with contact and task objects, making it possible to examine how a command becomes movement. The proposed chain keeps the neural predictor, goal, muscle controller, body mechanics and outcome score distinct. If the predictor selects the wrong key, a competent controller can execute that mistaken goal perfectly. Setting the joint directly to the desired pose would bypass the muscle-control question and must not be reported as a learned muscle action. These controlled computational studies still need their own tests of participant-specific biological accuracy. \href{https://proceedings.mlr.press/v168/caggiano22a.html}{Caggiano et al., 2022}.\par

NeuroTranslator names the implemented software joining particular recordings, models, replay and computer-body checks. NeuroTranslation Network names a proposed design for joining further components while keeping their origins and meanings clear. Neither name establishes full nervous-system coverage. Model size, learning from neural recordings and learning through simulated rewards require matched comparisons on the actual task and available hardware. A larger model or an impressive reward does not replace missing measurements.\par

Reaching, contact and release offer possible task relationships to compare across organisms. The mechanics still belong to each organism: a hand, paw and whole-body quadruped have different proportions, joints and control requirements. An averaged anatomy cannot establish a mapping that works for any of them. Physical scale, computational detail and statistical aggregation must also remain separate; a small organism can have a detailed model, while a pooled error can hide a failed recording. Biological generalization requires a test in the organism and conditions named by the claim.\par

\section{Current Study: Neither Model Passed}

The current hand study puts that demand for comparison into practice. Does
retaining recent neural history improve the stated prediction without
violating the other required conditions? Finishing a training run cannot
answer it; the original conjunction of tests must decide.\par

The current development study comprises six sessions from one macaque: four training dates and two development-validation dates. Ninety-six recording contacts in the source's HandM1 motor-cortex region contribute threshold-crossing-rate (TCR) and spiking-band-power (SBP) measurements, counts of detected signal crossings and power in an archived signal band. They give 192 features, not 192 independently located electrodes. Its targets are two source-normalized finger-sensor groups: index and combined middle/ring/small fingers. They are not measured joint angles, forces or a full-body state.\par

Eighteen runs saved 720 epochs, passes through the training material. Both gated recurrent unit (GRU) candidates, sequence models with a mechanism for retaining earlier information, fail their original model gates. Both fail the required no-harm and positive neural-information conditions: the predeclared comparisons must show the required absence of harm and contribution of the neural information. All five separate history, order-use and auxiliary-benefit claims also FAIL. Reduced aggregate error relative to a selected ridge baseline, a simpler predictor trained with a penalty on large fitted coefficients, does not override a failed required combination of criteria. The negative report records these failures; it does not establish that either model is effective.\par

The failed comparison also leaves a reproducibility problem. Both candidates use Torch automatic differentiation to calculate derivatives and Adam to update the model from those derivatives and its stored state. Agreement between candidates using that common route cannot independently certify the route. Full independent NumPy calculation using a distinct numerical software library remains unverified, as do complete model/update state, random-number-generator (RNG) state, reload and training-trajectory evidence. The final-checkpoint/update gate therefore remains unresolved.\par

The LINK study leaves its complete numerical update trajectory independently unverified, and its four reserved dates remain closed. Its recorded 18 fits and 720 epochs supply the failed development comparison, not a qualified decoder or GPU training route. The separate same-stream diagnostic in Section~\ref{sec:self-forecast} neither continues nor verifies that neural study. \href{research-evidence.json}{Study evidence and limits}.\par

\section{Earlier Studies}

An earlier human finger-attempt classifier obtained 116 correct predictions among 150 later trials after calibration using the first 50 labelled examples of each recording session. In its separate keyboard replay, 41 of 60 source-labelled right-finger trials produced the intended contact, while all 54 supported commands actually issued were executed; 13 presses were unintended or mistimed. The controller's execution score does not fix a prediction error. The model waited for a completed window; its earliest decision was one second after the cue plus computation. Inspected trials become development material if later tuning uses them.\par

A separate goal-driven simulated policy completed 32 of 32 scored phrase cases: eight reserved simulated orders crossed with two tempos and two poses, not 32 independent organisms. It uses a constrained generic four-key hand whose wrist, thumb and sideways finger-motion joints remain braced. It receives requested keys, phase, duration and smooth fingertip goals; displayed neural clips do not drive the simulated hand. The record preserves an interrupted original evaluator and a completed independent replay check; no test was reopened. Soft simulated key limits are not guaranteed physical travel bounds. A current saved development replay contains 336 frames at a 20 ms control interval, with 22 of 22 correct strikes, six of six chords, 22 of 22 timely releases, 34 ms worst timing error and no extra strikes. These are bounded simulation/readout results, not biological teaching, neural policy learning or unexpected real-world musical mastery.\par

The human recordings, macaque hand mapping and separate human EEG/ESG/ENG/EMG clips have distinct acquisition conditions and clocks. They do not demonstrate a voluntary central--spinal--peripheral pathway or drive a combined macaque/human arm simulation. Actual joint and force ground truth and physical unit calibration remain gaps. A task label, a normalized sensor value and a simulated joint angle must retain their different meanings. \href{research-evidence.json}{Programme computational evidence summary}.\par

\section{Failure Analysis}

The failure narrows the next question: where did the extra error come from? In the reported index-finger comparison, the auxiliary decoder's excess mean squared error, the average squared prediction difference, separates mainly into a mean offset and a smaller difference in residual variation. That decomposition locates a mismatch without identifying its cause. Trial composition, recording changes and state-dependent offsets could each contribute. Re-centering the validation predictions after inspecting that miss would use the answers to repair the prediction, so the adjusted result would no longer preserve the independent test.\par

Training-only context summaries show variation within trials, between trials and between dates; correlated overlapping frames are not independent observations. Named scalar calculations have bounded independent checks, including a preserved rejected skewness calculation, a measure of distribution asymmetry, and corrected arithmetic. Those checks do not demonstrate physiological state, neural drift or a useful new decoder. The proposed small context branch remains unexecuted. Any next fit requires independently checked training-only selection, normalization, controls, initialization, uncertainty and full update/state evidence, while preserving every original threshold. \href{research-evidence.json}{Current evidence and limits}.\par

\section{More optimization still lost to persistence}
\label{sec:self-forecast}

A model can improve on its earlier weights and still lose to carrying the last observation forward. A separate development diagnostic tested that possibility on one original TRAIN acquisition, sub-03/Cyl. It concerns self-forecasting of muscle and glove features, not the LINK neural decoder or the simulated hand. Its completed continuation supplies a new negative comparison rather than a successful four-source movement model.\cite{movementTrain990Recorded2026}

\subsection{Two tasks around one recurrent core}
The source is Reach\&Grasp version 1.0, released through IIT Dataverse with a declared CC BY 4.0 licence.\cite{reachGraspV1} The selected sub-03/Cyl pair is the cyberglove-motion file, provider ID 14486 (999,735 bytes), and the sessantaquattro-EMG file, ID 14393 (77,809,601 bytes). The retained acquisition record reports provider-MD5 verification for both. This attribution joins existing metadata; it supplies no fresh payload checksum, calibration or synchronization test, and does not settle every participant or combined-use right. The numerical result still concerns one acquisition.

The acquisition retained 3,123 rows from 3,127 source rows after excluding four invalid rows. It supplies 64 derived surface-EMG root-mean-square features with trailing 40 ms support and 18 glove-angle coordinates. The source declares millivolt-derived RMS and degree angles; physical calibration, synchronization, feature-availability times and online latency remain unqualified. A trailing feature summarizes an interval, so its row label alone cannot establish when all of its underlying samples would have been available.

Separate modality adapters and prediction heads surround one gated recurrent unit (GRU) core. Each task receives 64 rows of its own stream and forecasts its own later recorded features at positive row offsets 1 and 10. Sharing a core does not create a cross-modal target: no EMG-to-glove, glove-to-neural, force or afferent mapping is demonstrated. Recorded order is the temporal support being tested; no physiological delay or causal online acquisition is inferred from a row offset.

The first pilot made 40 joint optimizer updates with four ordered windows per task per update, exposing the first 160 windows of each task. Its normalized full-TRAIN errors exceeded both initialization and persistence. The once-completed continuation restored those parent 40 weights, AdamW moments and step 40, then the saved Torch random-generator state last, after model construction and optimizer loading. It made 950 additional joint updates: 949 batches of 32 windows per task and a final batch of 17. Thus 990 cumulative updates correspond to 30,545 window exposures per task, or five passes through the 6,109-window catalog. The catalog contains 3,059 offset 1 and 3,050 offset 10 windows per task. Repeated exposure changes optimization; it adds no independent acquisition or person.

\subsection{A matched TRAIN comparison}
Initial (0 updates), frozen parent 40, continued 990 and past-featurewise persistence were compared on the same catalog, targets, masks and saved original-TRAIN scalers. The scalers were fitted once on original TRAIN frames; continuation and evaluation fitted none. Persistence carries forward each feature's last supported past value. Evaluation took no optimizer step and kept the initial and parent comparison states frozen. The full catalog and the original first 4-window subset were reported separately; the subset is contained in the catalog, not an independent test.

Let $z$ denote a target on its saved normalized scale, $\widehat z$ a prediction and $m$ the observed-target mask. The reported score is
\[
 \operatorname{MSE}_{\rm norm}=\frac{\sum m(\widehat z-z)^2}{\sum m}.
\]
Predictions and residual coordinates use the recorded float32 route; squared-error sums accumulate in float64 and divide by the exact observed-feature count. Pooling sums errors and counts, rather than averaging offset scores without regard to support. The full pooled counts are 390,976 RMS features and 109,962 glove features; offset-specific counts and all literal sums are retained in the accompanying scientific summary. Normalized error is dimensionless, not squared millivolts or degrees, and its decimal places do not establish biological accuracy. No independent-acquisition uncertainty interval or statistical $p$-value is supplied by these correlated windows.

% Literal recorded metrics; no recalculation or numerical/model execution.
\begin{table}[htbp]
\centering
\setlength{\tabcolsep}{3pt}
\begin{tabular}{p{.29\linewidth}p{.32\linewidth}p{.32\linewidth}}
Condition & Surface EMG RMS & Glove angles\\\hline
Initial (0 updates) & 1.0195127528126222 & 0.8750431745338253\\
Frozen (40 updates) & 1.065703117562344 & 1.1992171367822229\\
Continued (990 updates) & 0.48371488886038944 & 0.37219431399606173\\
Persistence & 0.4391345726797162 & 0.024541819466343505\\
\hline
\end{tabular}
\caption{Pooled normalized mean squared error on all 6,109 original TRAIN windows per stream, across positive row offsets 1 and 10. Each condition uses the same targets, observed-feature masks and saved scalers. Scores are dimensionless and lower is better. The pooled counts are 390,976 observed RMS target features and 109,962 glove target features. Initial (0 updates) and frozen 40 are fixed comparison states. Continued (990 updates) improves on both but exceeds persistence in both streams; it remains unpromoted. These are correlated within-acquisition development observations.\cite{movementTrain990Recorded2026}}
\label{tab:self-forecast-pooled}
\end{table}

\begin{table}[htbp]
\centering
\setlength{\tabcolsep}{3pt}
\begin{tabular}{p{.29\linewidth}p{.32\linewidth}p{.32\linewidth}}
Stream and row offset & Continued (990 updates) & Persistence\\\hline
RMS, offset 1 & 0.41254365424001344 & 0.14465972419162815\\
RMS, offset 10 & 0.5550961369599728 & 0.7344783633436713\\
Glove, offset 1 & 0.34938925569254925 & 0.0004908467144512168\\
Glove, offset 10 & 0.3950666659142399 & 0.048663762301765986\\
\hline
\end{tabular}
\caption{The same full-TRAIN comparison separated by row offset. Offset 1 has 3,059 windows and offset 10 has 3,050 per stream. Only RMS at offset 10 beats persistence; a favourable offset does not replace the pooled failure or the other offset's result. A row offset alone establishes no physical forecasting delay.\cite{movementTrain990Recorded2026}}
\label{tab:self-forecast-offsets}
\end{table}

\begin{table}[htbp]
\centering
\setlength{\tabcolsep}{3pt}
\begin{tabular}{p{.29\linewidth}p{.32\linewidth}p{.32\linewidth}}
Condition & Surface EMG RMS & Glove angles\\\hline
Initial (0 updates) & 0.7221646495859888 & 1.4936251670580845\\
Frozen (40 updates) & 0.13923739944560742 & 0.781681849149845\\
Continued (990 updates) & 0.12858087160530363 & 1.5822619236435527\\
Persistence & 0.20836011081823158 & 0.0004564336184945723\\
\hline
\end{tabular}
\caption{The originally selected first 4 windows per stream, two at each offset: 256 observed RMS target features and 72 glove target features. This subset is included in the full TRAIN catalog and in the initial optimization prefix. Its RMS improvement cannot establish independent generalization; its glove score worsens after continuation. No subset is substituted for the full-catalog comparison.\cite{movementTrain990Recorded2026}}
\label{tab:self-forecast-prefix}
\end{table}


\subsection{Improvement is not the comparator}
Final 990 lowers full-TRAIN RMS error from 1.0195127528126222 at initialization and 1.065703117562344 at 40 updates to 0.48371488886038944; persistence remains lower at 0.4391345726797162. Glove error falls from 0.8750431745338253 and 1.1992171367822229 to 0.37219431399606173, while persistence is 0.024541819466343505. Both streams therefore fail the retained persistence comparison. The earlier 40-update failure remains visible rather than being replaced by the later improvement.

The pooled result has structure. RMS at offset 10 beats persistence, while offset 1 does not; glove loses at both offsets. In the original first 4 windows, RMS improves, but glove worsens relative to both initialization and parent 40. The prefix cannot be chosen to rescue the complete result. Software/runtime and recorded arithmetic checks support the specified calculation; they do not promote the private checkpoint or demonstrate model quality on new data.

Batch size changed from 4 to 32 while the update budget and repeated catalog exposure increased. Those changes form one combined engineering intervention. They do not isolate coverage as the cause of improvement, test added dataset diversity or establish a quantity/diversity power law. This remains one TRAIN acquisition, with no DEV or TEST generalization result. Nor is there a neural observation, cross-modal translation, simulated force target, biological feedback or retained-skill outcome in this diagnostic.

That leaves a concrete requirement for the broader proposal: a representation must beat the appropriate simple predictor on its stated source before shared weights can support a larger translation claim. Broad neural observations, myography, measured movement and physics-generated examples remain four distinct proposed training groups. Their source rights, physical reference series, timing, calibration, sensory context and independent held-out tests remain necessary; a shared GRU and repeated training on one acquisition cannot supply them. Brain2Body's further question still begins when trainer commands stop: what does the trainee retain under matched controls?


\section{Viewer Capabilities and Missing Measurements}

A graphical user interface (GUI) should let the reader tell a measurement from a prediction, a simulation from a missing value. Saved replay must preserve the source-time clock and complete-run metrics, even when the reader pauses at a persuasive frame. Cache-only playback displays an identified saved result without silently starting a new experiment; an uncached selection must show absence, and a failed source must clear stale geometry and results. Natural-language goals should become explicit task specifications. Receiving files in a chat does not itself establish inference or body control, so the interface must show what operation the system actually performs.\par

A saved replay and a live research view therefore need separate checks. The reader should be able to see which one is available, which source it uses and which measurements are missing.\par

\subsection{What the online and offline versions require}

The online candidate needs its whole directory. A disconnected reader instead
needs the separate static export compiled from the original source; research
views that use the application programming interface (API), the interface to
the software service, are unavailable offline. Neither route has current
whole-interface verification. Earlier bounded replay checks cannot establish
that the present interface or the separate export works.\par

The static package and original-source offline export have different dependencies. Their file sizes, missing browser comparison and unavailable API-dependent views are described in Appendix~\ref{app:reproducibility-availability}. Neither establishes complete current interface function or an interaction-speed improvement.\par

\section{Anatomical and Functional Validation}

A viewer can expose an error only if the body it shows is itself accountable to references. Begin with proportions, muscles, bones and joints, then trace the neural routes and check contact and viewing geometry wherever those features are claimed. Each component and the assembled scene need lifelike three-dimensional review against a series of real-world references. Biological identification colours, useful close-ups and tilted orthogonal views can make those comparisons legible. Transparent layers, cross-sections and paired views can reveal relationships; illustrative signal paths must remain distinct from physical parts. This review tests the anatomy we show. Hidden physiological function needs its own evidence, however plausible the image. A conceptual schematic remains an explanation, not a biological acceptance result.\par

A template surface or atlas provides a reference. Registration has to establish its correspondence to the recorded individual's anatomy before a location in the template can be attributed to that individual. The named source-binary reconstruction matches template geometry and atlas values, and its optional region boundary retains disclosed surface-mesh defects. Those bounded correspondences neither localize the recorded neurons nor establish their function. The retained framework illustration is historical and conceptual; its labels describe earlier computational work, not the current model-gate outcome.\par

The proposed left and right tripod views of the hands and arms need calibration, occlusion checks and independently verified timing. The motion-capture relay-composite would align or stitch those views. A separate perpendicular LiDAR view, using light to measure distance, may make particular surfaces easier to observe; it is not guaranteed to reveal every surface or solve neural registration. The acquisition method must earn its reconstruction claim against real measurements.\par

\section{Proposed Skill Tests}

Once the movement and its measurement can be trusted, the next task should expose something the present test cannot answer. Key presses separate selection, contact and timing errors. Threading adds flexible-object geometry, friction and corrective contact; puzzle manipulation adds object state and coordinated grasping. Archery or breakdancing needs forces, balance and wider body coverage absent from the current hand study. This gives the task ladder a reason to grow without pretending that a classifier or successful replay has already climbed it.\par

Reserve new task orders, object properties, starting positions and recordings before selecting the method. Then give simpler controllers and copied-example methods the same information and practice, and score the complete task: success, unintended actions, delay and failure. A selected clip cannot answer that comparison. Right-angle and non-right-angle mazes could test orientation and memory demands, but their difficulty needs measurement; 10\ensuremath{\times} or 100\ensuremath{\times} labels do not provide it. Computer training duration likewise cannot establish biological learning duration.\par

\section{Development, Twins and Disease}

The body also changes while a mapping is being learned. The native-development branch asks how recording access, maturation and changing function affect that mapping. Cook and colleagues provide species-specific skeletal staging in the fat-tailed dunnart. It does not establish corresponding timing in antechinus, a safe neonatal neural interface or an artificial rearing system. Armstrong and Ferguson studied age-related skin-wound scarring in \emph{Monodelphis domestica}; skin healing does not certify scar-free neural access, spinal repair or stable electrodes. \href{https://www.nature.com/articles/s42003-021-02506-2}{Cook et al., 2021}, \href{https://pubmed.ncbi.nlm.nih.gov/7750642/}{Armstrong and Ferguson, 1995}.\par

Identical-twin, non-identical-twin and unrelated comparisons would ask which sources of variation affect a mapping. Similar genetics leaves differences in experience, physiology and recording location to be measured. Each comparison needs matched tasks, developmental stages and acquisition conditions before a channel or skill-transfer claim can be interpreted. The twins model must undergo component and assembled-scene reference review as well; kinship and a convincing linked image supply neither functional evidence nor visual acceptance.\par

Human-cell and graft branches ask defined questions about tissue contribution and host interaction. Revah and colleagues found maturation and circuit integration of transplanted human cortical organoids, laboratory-grown tissue representing aspects of the brain's outer layer, in rats. This does not make the whole host human, transplant expertise or demonstrate injury repair. The later image-panel correction is retained in the citation record. Native development, altered development, mixed-species tissue, grafts and disease conditions are separate proposed comparisons, each requiring its own evidence. \href{https://www.nature.com/articles/s41586-022-05277-w}{Revah et al., 2022}, \href{https://www.nature.com/articles/s41586-026-10943-4}{author correction, 2026}.\par

A disease experiment must identify the impaired function and the change produced by the interface. Compensation for interrupted motor output, improved task performance, retained learning and tissue repair require different endpoints. The cause of one condition may leave a route usable that another condition disrupts, so a benefit cannot be carried across them without its own test. These branches remain proposed; the present computational record demonstrates none of them.\par

\section{Next Experiments}

The failed comparison leaves us short of a qualified predictor, before we even reach the teaching question. The measured targets are still two normalized finger-sensor groups, without calibrated joints, forces or a complete body state, and neither recurrent candidate is admitted. A constrained controller could be studied by declaring its observations, permitted actions and feedback. If it receives explicit goals, its result must remain a goal-driven result; it cannot stand in for a neural prediction or evidence of biological learning. Only with those contributions separated can a trainer comparison ask whether useful movement assistance leaves lasting capability.\par

The next comparison can start with the grasp itself. The proposed two-person movement-only study asks whether pretraining on a qualified source person's movement improves a target person's movement forecast. Compare target-person pretraining with source-person pretraining under deliberately incorrect correspondence and under the proposed correspondence, matching architecture, initializations, update budgets, target-training batches and final-checkpoint selection. The proposed shared representation contains 13 movement features and 13 missing-feature masks: wrist-relative index, thumb and object sensor positions, digit aperture and causal wrist velocity. Proper wrist frames and a characteristic sensor-distance scale fitted only on each person's training observations separate coordinate orientation and size; a common 0.1 second reference interval also makes the velocity features dimensionless. The scale is not a measured bone length. Frames, units, clocks, gaps and sensor calibration still need qualification before this representation can be used.\par

Both homologous and mechanically analogous features remain eligible for the wider programme. This particular comparison tests a declared human sensor-landmark correspondence, not cross-species neural equivalence. Motion from another recording can supply a learned movement constraint without supplying a neural observation paired to the target person's action. Scans may constrain a specified anatomical gap, but an external scan or simulated movement cannot supply that missing measured neural pairing. Keep acquisition parents, overlapping segments and derivatives in one role before windowing or fitting statistics; development observations must not influence training scales or become independent final confirmation. This comparison remains unexecuted. Any later neural adaptation must use actual paired observations and retain its distinct input sets, causal timing, controls and protected partitions.\par

A trainer experiment needs supported interfaces and new evidence. Compare experienced, inexperienced and artificial trainers with matched task information, practice, assistance budgets and relevant burden. Include sham, delayed and unrelated-feedback controls. Improvement while guidance is present does not establish retained learning.\par

Withdraw trainer information while keeping any necessary movement bridge, practice, feedback and task access constant. Then measure delayed performance and unfamiliar variations against matched controls. Record the help that remains. If a gain fails to persist, that failure counts against learning. A retained gain still cannot establish tissue repair, restored pathways or regeneration: those claims require separate biological evidence.\par

\section{Biological and Artificial Trainers}

Brain2Brain names specified information links between nervous systems. Pais-Vieira and colleagues linked recorded activity and stimulation to share task-relevant information between trained rats. Prior training and dependence on stimulation limit what the result means: it is not instant acquisition of another animal's general skill repertoire. \href{https://www.nature.com/articles/srep01319}{Pais-Vieira et al., 2013}.\par

The proposed biological-trainer experiment compares experienced and inexperienced partners with matched practice. Genuine guidance, sham guidance without the intended information, delayed guidance and matched feedback unrelated to current need test different explanations. A simple artificial controller tests whether the neural trainer adds benefit beyond a less elaborate route. Random assignment, recording of actual help and independently scored outcomes are needed to distinguish expertise, information, timing and additional exposure.\par

An artificial trainer would supply guidance from a model trained on supported observations. Its comparison begins with what the model receives and controls. Desired goals, expert trajectories, privileged labels and sensory feedback give it different resources; supplying the correct action removes the prediction problem that a neural decoder is supposed to solve. Trainers therefore need matched task information, practice and relevant burden. A useful gain while guidance remains present can support an assistance claim even if no lasting learning follows.\par

\section{Trainer-Withdrawal Tests}

Test retention by withdrawing trainer commands while holding any necessary
compensatory movement bridge constant. A bridge lets a trainee express action
despite an interrupted pathway. Remove that bridge together with teaching and
movement may become impossible regardless of learning: the test would confound
body access with skill. Keep movement access, feedback and task conditions
comparable between the baseline and later test.\par

Measure immediate performance, delayed retention and unfamiliar variations separately. Record how much help remained, whether task performance improved without trainer information, and whether the benefit exceeds matched controls. If the gain disappears when trainer commands stop, report assistance. If a gain persists, test its domain and uncertainty. Tissue repair still requires separate biological evidence; a retained task gain or successful compensatory interface cannot establish regeneration.\par

BrainNet is a network metaphor for specified links. It is not an operating market for biological mastery. Its value would have to appear in what a trainee can do through those links. Broad enablement and possible skill services motivate the programme, but they do not establish deployability, ownership of biological information or superiority over every other method. The proposal must earn support through its outcomes.\par

The failed comparisons set a condition on the broader ambition: a more elaborate learner must earn its advantage over the appropriate baseline. Grounding its outputs in body mechanics gives the next mechanism to test; compatible source observations constrain what can train it. More examples and greater diversity then form a separate scaling hypothesis, not an explanation already established by repeated optimization. The task sequence makes those proposed relationships observable. A paired retention test finally asks whether any benefit survives after trainer commands end, with necessary movement support held constant.

% R152 source proposal; provenance and exact prior text remain in the proposal manifest.
\section{A grounded conditional learner with distinct output interfaces}
The retention question now needs a mechanism: what information would guide
the movement, and what independent observation could show that the guidance
was useful? Follow one action through the proposed measurements and learner
before adding a writer or attributing a gain to biological teaching.

Start with a movement we can measure well. For the proposed snap showcase,
the arm begins extended with the palm down. The elbow and forearm swing and
turn; the wrist and hand turn upward during the snap, then move down to rest.
The model must account for that whole movement, including the preparation,
contact and return. Finger pairing, organism-specific anatomy and timing still
need a real-reference series; a human-looking animation cannot establish that
a particular monkey can perform the same gesture.

The cup reach supplies a simpler precision comparison: relate recorded signals
to measured movement before attempting the coordinated snap. Each task still
needs its own anatomy, contact, timing and independent comparison.

That comparison connects three interacting groups: task-linked brain
observations, named-muscle observations and a musculoskeletal simulation.
A learner would relate the observations to the mechanics; independent movement
and contact measurements would test its predictions. EEG records scalp
electrical activity and surface EMG records muscle electrical activity.
A separately recorded nerve or spinal signal needs its own declared channel.
The proposed apparatus uses two arm sleeves, scalp EEG and synchronized,
calibrated left and right tripod views. Vendor, channel, bandwidth and clock
qualification remain open.

Intent to action is the target. By understanding and measuring the intermediary,
we aim to make that mapping as close to one-to-one as possible within a declared
task. Specify what counts as the intended action, the participant and session,
state and context, permitted history, tolerated error and independent reference.
The test then asks whether observing the middle reduces ambiguity and improves
prediction or qualified control. It does not assume a unique hidden intention
or claim that the best possible mapping has already been achieved.

The learner's first test would compare its predictions with recorded brain
and muscle activity and with measurements of where the hand moved and when.
Sensory feedback and physical action would add a further question: what did
each separate output interface actually deliver? A successful prediction
cannot answer that question, and EEG and EMG recorders alone cannot write
sensation. For this proposal, no paired trial, trained model, delivered sensory
output or physical actuation is established. Sources and permissions, physical
and computational behaviour, clinical interfaces, reference comparisons and
held-out tests still require independent assessment.

The retained LINK comparison belongs to a different experiment. Its 18 runs
and 720 saved epochs are not paired EEG/EMG/movement trials; both model-admission
gates and all five benefit claims failed. The prospective methods below would
have to supply the missing paired evidence and separately test recording,
prediction, simulated movement, sensory delivery and physical action. A result
for one of those jobs would not establish the others.

\subsection{Form, function and action need a measured middle}
A structural map constrains what a system could do. A NeuroConnectome maps
specified neural connections; its organism, coverage, resolution and source
would have to be declared. The signals passing through those connections
still have to be observed or inferred, with the difference made explicit.
To predict the next action, we need the processes that turn structural
constraints into an outcome. The proposed programme follows form to function
to action, while intent to action names the task-level translation we want
to improve.

DNA to RNA and RNA to protein put a process between a starting constraint and
its downstream product. Regulatory context, RNA processing and cell state
change what that process produces; protein structure alone does not determine
the behaviour of a whole cell. AlphaGenome, RNA or cell-state models and
AlphaFold illustrate these different modelling jobs. AlphaTranscriptome names
a desired predictive capability here. Connecting the predictors would still
leave system dynamics and the environment to account for, with experimental
outcomes needed to test the combined prediction.
% Source-role pointers: alphagenome2026, scgpt, alphafold3; bounded owner checks.
\par\noindent \href{https://www.nature.com/articles/s41586-025-10014-0}{AlphaGenome regulatory prediction}~\cite{alphagenomeRetained};
\href{https://www.nature.com/articles/s41592-024-02201-0}{scGPT cell-state modelling}~\cite{scgptRetained};
\href{https://www.nature.com/articles/s41586-024-07487-w}{AlphaFold 3 structure prediction}~\cite{abramson2024}.
These provide modelling examples, not evidence of an integrated cell programmer.
Our account of AlphaGenome and scGPT is limited to their reported modelling roles;
their complete methods have not been independently assessed here.

In the biological pathway, Brain2Spine and Spine2Body name proposed links to
test; their composition is the Brain2Body aim. The spinal cord integrates
descending and sensory signals and contributes local coordination, reflex
regulation and adaptation. This is a different mechanism from RNA processing,
but it serves the same explanatory purpose: account for the context-dependent
process between an initial constraint and its outcome. Supported spinal
observations would test that account. Where they are absent, the internal
states remain inferred or simulated. No programme spinal intervention or qualified spinal-writing interface is demonstrated here.
\par\noindent \href{https://pubmed.ncbi.nlm.nih.gov/38603479/}{Spinal sensorimotor adaptation}~\cite{lavaud2024}
is a constrained example, not evidence of unrestricted skill transfer or
validation of our proposed pathway. Our comparison draws on the abstract and
relevant author-manuscript passages; equivalence to the final published edition
and the complete methods and statistics remain unverified here.

Thought2Action relates a declared intention to an independently measured action.
Action2World follows the action to its consequence in the task environment;
Thought2World names the proposed composition. Body2World drew attention to the
consequence outside the body; Body2Action focuses on body mechanics and a
specified task outcome, while Brain2Motion focuses on the movement pathway.
These working labels identify related scopes for a test. They do not establish
interchangeable or working controllers.

The engineering route connects neural activity to musculoskeletal simulation
and then to a programmable model, toward NeuroTranslation. The proposed
movement-focused learner is a large movement model. System 3 thinking names
the proposed human--AI supervisory partnership around it. The supervisory
model would compare task strategies and diagnose where neural or mechanical
predictions miss independent measurements, then propose training or calibration
comparisons for the signal mapping and virtual body. The task learner would
perform the declared prediction or control; the supervisor would help choose
what to test. Neither role replaces the biological spinal intermediary or a
physical interface. Final held-out outcomes must remain outside rewards and
model selection.
\par\noindent Shaw and Nave's SSRN preprint,
\emph{Thinking---Fast, Slow, and Artificial}, and its
\href{https://executiveeducation.wharton.upenn.edu/thought-leadership/wharton-at-work/2026/05/thinking-fast-slow-and-artificially/}{institutional discussion}
provide conceptual context for this proposed human--AI role.~\cite{shawNaveSystem3Retained}
The comparison is limited to that framing; the complete preprint and institutional
article have not been reviewed here. Neither establishes a third anatomical brain
system or validates this programme's proposed learner or supervisor. Expert footage may supply
movement and tactical examples, with rights and camera uncertainty accounted
for; it supplies no measured expert EEG, EMG, muscle force or stimulation plan.

Action changes the outcome. Its sensory consequences inform perception and
the next intention. Independently measured movement would anchor the action
part of that loop; it would not by itself decode a unique sensation or
perception. Known stimuli, measured feedback and appropriate recordings or
reports would have to test those further mappings. The return path must
therefore be measured rather than drawn into existence. The full
Interoceptomics aim includes internal-body sensing and connectivity; external
touch and proprioceptive tasks contribute evidence but do not by themselves
solve that wider problem. A prediction, assisted execution and retained skill
without assistance answer different questions. Each wider claim must earn
its own independent comparison.

\subsection{Context connects. Grounding checks}
Within a reach, context has a concrete job: relate brain activity to muscle
activity in this task and at this time. Attention is one proposed way to
represent those relationships. The model compares a query with keys and uses the resulting weights to combine their values. Masks keep unavailable observations and future
samples out. The calculation weights information; it does not establish that
attention understands all context.
% Source registry IDs: context-native-01a12f62, context-attention2017-v7, context-mult2019, context-policy4
\par\noindent \href{https://arxiv.org/html/1706.03762v7}{Vaswani et al. (2017)}~\cite{vaswani2017}; \href{https://aclanthology.org/P19-1656/}{Tsai et al. (2019)}~\cite{tsai2019}. These methods motivate the computation; validation of the proposed paired system still requires the independent observations below.

A pattern can fit inside the model and still describe the wrong movement.
Grounding therefore means testing the declared prediction against qualified
observations under the same conditions. Independently measured movement would
check the outcome, while muscle activity would add an observation that
constrains the brain model. Even a matching reach cannot identify a unique
private experience. No acquired paired data or trained model is established
for this proposed loop; those comparisons still have to be made.

\subsection{Muscle activity tests the brain model}
Record brain and named-muscle activity during the same reach. In the proposed
paired trial, EEG would sample scalp electrical activity and surface EMG would
sample activity associated with the named muscles. Paired training could use
the muscle observation to constrain their relationship. Then withhold the
muscle target and movement reference from the EEG-only predictor: can the
permitted brain observations predict a specified muscle-activity observable
or movement?

Pairing and calibration remain proposed. Shared cues and movement artefacts
could give that predictor a shortcut, so the comparison needs controls for
both. The muscle observations would cover only the named muscles, and EMG
supplies neither afferent nerve activity nor direct force. Sensory claims
still need known stimuli and measured feedback. Several neural states could
remain compatible with even a well-matched pair; the pair could narrow a
task-specific explanation without revealing the complete hidden brain.
Checking the movement itself requires the next independent observation.
% Source registry IDs: context-native-01a12f62, context-policy4, overhaul-brantley2018
\par\noindent Brantley and colleagues reported simultaneous scalp EEG, lower-limb EMG and locomotion observations from 17 body-worn inertial measurement units (IMUs) in ten able-bodied people. This supplies a paired gait precedent; it does not establish the proposed hand-reach apparatus, afferent recordings or optical three-dimensional motion capture. Selected-file rights, clocks, calibration and movement-artifact handling still need qualification. \href{https://www.nature.com/articles/sdata2018133}{Brantley et al., 2018}.

\subsection{Measured motion tests the virtual body}
Measure the same hand and arm from the aligned left and right tripod views.
The proposed motion-capture relay-composite would align the two point maps
in a common clock and coordinate frame, then join them into one composite
while retaining each point's origin. Before that joined map can check a
prediction, its calibration, alignment and capture precision would have to
be independently verified.

The movement map would constrain the external movement model: surface motion
and pose in the task's physical frame. To explain that motion mechanically,
the internal musculoskeletal model would describe specified bones, joints
and soft tissues. Organism, developmental stage and task would determine the
appropriate anatomical sources. X-ray projections could inform bone and
joint relationships; appropriate MRI references could inform specified soft
tissues. Registration, rights and source-specific limits would need independent
checks. No scan set is selected for this proposed comparison. Anatomy still
leaves activation, forces and neural function to test.

With a qualified composite, the musculoskeletal-simulation physics-engine's
joints, muscles, contacts and physical limits would face a measured comparison.
Fit the model on training data, then compare its predicted reach with
held-out measured movement. A force claim would require an additional force
measurement. The body's form limits possible motion; measured function would
test the mechanics and show where they need refinement.
% Source registry IDs: context-native-01a12f62, context-component-names-01a12f81, context-policy4, context-opensim2018, context-myosuite-docs, context-precision-01a12f77
\par\noindent \href{https://journals.plos.org/ploscompbiol/article?id=10.1371/journal.pcbi.1006223}{Seth et al. (2018)}~\cite{opensim2018}; \href{https://myosuite.readthedocs.io/en/latest/}{MyoSuite documentation}~\cite{myosuiteDocsUndated}. The simulation tools supply methods for that comparison; the proposed capture and body models still need its independent validation.

A lifelike hand makes that comparison visible. It cannot establish biological
fidelity. Capture alignment, precision, mechanics and biological fidelity have
not been demonstrated by this proposal, and captured geometry cannot establish
hidden neural state, internal forces or unique subjective experience. The
measured reach would check a movement prediction within those limits.

\subsection{Two pairs. One conditional learner}
The two comparisons would meet in the proposed large movement model: EEG
observations paired with muscle activity, and measured movement compared
with predicted mechanics. Large movement model names the proposed family;
deep neural network describes a possible implementation. Large refers to
the intended scope, rather than a demonstrated parameter count. Encoders
would receive only the observations permitted for the named test. Two
branches would summarize the relationships and meet in a shared task
representation; separate output parts would predict each direction's targets.
% Source registry IDs: context-native-01a12f62, context-policy4, context-attention2017-v7, context-mult2019
\par\noindent \href{https://arxiv.org/html/1706.03762v7}{Vaswani et al. (2017)}~\cite{vaswani2017}; \href{https://aclanthology.org/P19-1656/}{Tsai et al. (2019)}~\cite{tsai2019}. These computational methods offer precedents for a design, leaving the proposed paired learner and its performance to be tested.

Signal features belong to branches, and branches belong to the learner.
That hierarchy describes the computation, not two anatomical thinking systems.
A common clock, task and physical frame would connect the branches; independent
measurements would check their errors. Form constrains the prediction, and
checked function would guide refinement. Several hidden states could still
fit the same reach, so neither a shared representation nor conditional
predictions would establish an exact inverse nervous system. The architecture
has not been selected, implemented, trained or independently accepted.
To choose it, we would have to compare it with simpler baselines on
appropriately reserved observations.

\subsection{Learn the movement, then test what transfers}

A reach can provide a training target even when no EEG accompanies it. Earlier movement supplies the input; a later observation supplies the target. A separate proposed comparison would ask whether learning that relationship in one person helps forecast another person's grasp-and-lift movement. The WAY-EEG-GAL descriptor provides a relevant continuous-recording candidate, alongside its distinct windowed trials and derived lift summaries.\citep{luciw2014way} This movement-only question needs no AllLifts Run/trial join or invented cross-person neural pairing. The original condition-specific five-weight study remains a separate design.

The mapping has to explain what two bodies share. Homologous structures share evolutionary anatomical origin; mechanically analogous structures can serve comparable roles without that origin. A human digit aperture and a robot gripper opening could support a task-space comparison after separate mechanical qualification; this human-only proposal would not test that robotic extension. Joint degrees of freedom, axes and signs, proportions, muscle--tendon roles, loads, inertia, contact and time scales remain body-specific. Similar geometry cannot establish equivalent dynamics or neural coding. Retargeted motion remains derived; tracking, estimation, prediction and simulation keep their different evidential roles.

For the proposed two-person comparison, the representation contains index-, thumb- and object-sensor positions relative to the wrist, digit aperture and causal wrist velocity: thirteen features, with thirteen availability masks. Each participant needs qualified position units, a proper wrist-frame rotation, orientation convention and sensor-attachment calibration. A positive median wrist-to-index-sensor distance, expressed in qualified metres, supplies the characteristic length from assigned training support only; it is not a bone measurement. Multiplying normalized velocity by the same fixed 0.1-second reference interval makes all thirteen features dimensionless. Physical velocity errors invert that scaling with length divided by 0.1 seconds. Missing features receive no target loss; an unavailable observation is not a measured zero. Full finger joints, individual muscle forces and subject-specific inertia are not supplied by these sensor features.

The continuous target would be the same observed feature set 100 milliseconds beyond the latest actually available filtered input time. The proposal retains fifty frames at 100 Hz and a causal 51-tap filter before decimating qualified 500-Hz kinematics by five. Its 50-millisecond linear-phase delay must remain in the physical timing account; shifting information backward would change the prediction claim. Coefficients, passband, timestamp convention, gaps and boundary handling still require qualification. Only named kinematic columns may enter: the documented kinematics matrix also contains force-plate signals.

Identity and independence answer different questions. Trustworthy versioned participant-release attribution, an immutable selected archive entry and a qualified capture can identify the participant record and input object without requiring one particular internal ID-field spelling. They cannot establish biological identity or independent original acquisitions. Shared samples, trims, derivatives and overlapping recordings must form indivisible acquisition groups before filtering, windows, normalization or role assignment. Different names, checksums or series numbers do not prove separation. The complete public non-ST candidate roster and these relations still need establishing; all 294 rows of the prior AllLifts table remain excluded from training, evaluation and assigning trials or acquisition groups, and ST, competition storage and the reserved LINK recordings remain unavailable.

Qualified P1 groups would supply source pretraining; P2 would supply target training and development. With $G\geq2$ independently established P2 acquisition groups, the fixed identifier-hash ranking assigns the first $\lfloor3G/5\rfloor$ to training and the remainder to development. Canonical original-parent identifiers, their serialization and the roster must be fixed before ranking or examining outcomes. This is a nominal 60/40 split, not chronological ordering or a promise of condition balance. If unresolved ancestry merges P2 into one group, the comparison stops; splitting its windows cannot manufacture an independent development group. No group or segment has been assigned by the proposed identity method.

Three matched arms would test the transfer. B0 pretrains on correctly mapped P2 training movement. B1 and B2 use identical P1 windows and order: B1 sets the index-role input to the thumb sensor, the thumb-role input to the object sensor and the object-role input to the index sensor, while B2 retains the correct correspondence. Both share the correct source calibration and scale before that corruption, keep the wrist anchor, and use the same correct target adapter and ordered target batches. Parent-balanced sampling must record unique windows and repetitions. The false mapping can be relearned or can alter feature availability; matched update counts do not match source diversity. These are limits on attributing any difference to anatomical correspondence, not reasons to replace the control after seeing its result.

The proposed learner is one GRU layer of hidden width 32, with 26 inputs and thirteen outputs. Each arm and fixed seed---17, 31 and 73---receives 200 pretraining and 100 target updates, batch size 16, identical per-seed initialization and the fixed last checkpoint. Adam uses learning rate 0.001, betas 0.9 and 0.999, epsilon $10^{-8}$ and zero weight decay, without clipping, scheduling or early stopping. Training minimizes mean squared residuals over observed target features; a batch with no observed targets is rejected. The resulting 2,700 updates are planned, not performed. No external checkpoint is selected; numerical updates, filtering, precision and actual workload limits remain to be checked independently.

For each development parent and seed, the proposed score takes the square root of summed squared errors divided by the count of observed target features, then averages parent scores without weighting by parent duration. A zero denominator rejects the comparison. Report feature counts, missingness, physical-unit errors, uncertainty and persistence at the same information endpoint, including every worse parent or feature. B2 would have to lower normalized error against both B0 and B1 in every fixed seed, with no applicable original criterion failing, to justify another investigation. This remains a two-person development comparison, not independent final confirmation, population transfer or biological teaching.

Broader motion collections could supply additional movement examples only after their exact assets, rights, task coverage and original-recording overlap are qualified. In \href{https://www.addbiomechanics.org/download_data.html}{AddBiomechanics}, measured markers and ground reactions must remain distinct from fitted kinematics and estimated inverse-dynamics torques. \href{https://amass.is.tue.mpg.de/license.html}{AMASS} requires checking the selected asset's research-use, non-commercial and redistribution restrictions. Copies across collections stay with their original acquisition. LaBraM's WAY source-family exposure is flagged, while exact checkpoint/participant/trial exposure remains unknown; reserving a recording later cannot erase earlier pretraining exposure. No such checkpoint is selected for this comparison.

Before adding X-ray/CT or MRI, inspect the anatomical sources already used by the selected musculoskeletal model. A scan is useful when it constrains a stated gap or supplies an independent check beyond that model. Its person or species, scale, registration, rights and uncertainty still matter. Another person's scan remains an external anatomical reference; a generic model does not implicitly measure this participant's geometry, activation or forces. No scan set is selected here.

Actual source eligibility, acquisition groups, units, frames, clock and calibration remain unresolved for this continuous comparison, so no new fit or successful transfer is reported. A useful movement prior could later be tested with actual paired neural/body observations and the permitted input masks. Its transferred targets would still not become simultaneous measurements. The next result must identify where the mechanical forecast improves, where it worsens and which measurement or assumption accounts for the difference; neural equivalence, causal physiology, sensory delivery and retained skill would require further experiments.

\subsection{Pretrain. Pair. Verify. Refine}
Those comparisons start with the observations we may use. Data rights,
calibration and held-out subjects, sessions and tasks determine which can
enter the study; resolve them before processing data or choosing a model.
Open-source code does not grant blanket permission to use datasets,
checkpoints or clinical interfaces. Qualified data could first train
representations for each recording type. LaBraM is one EEG encoder candidate,
not a complete brain-body writer. Synchronized paired observations could
then train the declared mapping from permitted inputs.

Paired prediction would give later refinement a baseline to beat. Before
optional reinforcement learning, post-training or task-specific fine-tuning,
the physical predictions and reward need checking against independent
outcomes. Offline reinforcement learning needs qualified coverage of actions,
transitions and rewards: labels alone cannot tell us what an untried action
would do. Backpropagation can update differentiable neural components, but
it cannot automatically pass through every physics engine; universal
differentiability is not established. Simpler baselines and frozen final
tests remain through every phase.
% Source registry IDs: context-native-01a12f62, context-policy4, context-labram2405, context-cql2020-abstract, context-opensim2018, context-myosuite-docs
\par\noindent \href{https://arxiv.org/html/2405.18765v1}{Jiang, Zhao and Lu (2024)}~\cite{labram2024}; \href{https://arxiv.org/abs/2006.04779v3}{Kumar et al. (2020)}~\cite{cql2020}; \href{https://journals.plos.org/ploscompbiol/article?id=10.1371/journal.pcbi.1006223}{Seth et al. (2018)}~\cite{opensim2018}; \href{https://myosuite.readthedocs.io/en/latest/}{MyoSuite documentation}~\cite{myosuiteDocsUndated}. These sources offer candidate methods for the sequence; validation of the proposed paired system still requires the experiment below.

No data acquisition, training, fine-tuning, RL rollout or accepted checkpoint
is established for this proposed loop. The next experiment must test what
the prediction and reward represent against independent outcomes. A reward
that the computer can calculate supplies no biological proof on its own.

\subsection{Respond quickly. Improve carefully}
\label{b2-frozen-response}
The proposed response path would use a frozen accepted version to estimate
the current state, predict declared targets and prepare encoded outputs for
a separately qualified interface. Permitted causal time windows would limit
it to the allowed past and present observations. Estimating the state of this
moment would not change learned weights. A separate adaptation path would
inspect errors, calibrate or recalibrate, and propose a candidate update;
an independent comparison and a versioned saved state whose restoration reproduces the specified computation would precede any replacement. Examining errors would not
silently change the running model, and no unchecked live weight promotion
would occur.

The fast-and-slow comparison is a Kahneman-like analogy for this proposed
division of work. It identifies neither two biological or anatomical systems
nor inherent thinking speeds for EEG, spatial data or other sensor modalities.
Real-time performance would have to be measured across acquisition, filtering,
alignment, inference, output preparation, transduction or actuation, feedback
and rendering. Report delay tails, jitter, missed deadlines, uncertainty and
performance under specified concurrent workloads; component inference time alone cannot
establish end-to-end performance. No accepted model, running system that adapts
under these rules or real-time demonstration is established here. The loop has
no measured end-to-end latency or accuracy guarantee.
% Source registry IDs: context-native-01a12f62, context-policy4

\subsection{Write sensation. Enact action}
\label{b2-output-interfaces}
The prediction ends at the computer. A separate interface would have to turn
it into sensation or action. Input and output describe direction relative to
that computer; afferent and efferent describe physiological direction relative
to the central nervous system. A computer output could drive a sensory
interface supplying input toward the nervous system, while a computer input
could record activity related to the body's motor output. Both boundaries
must remain explicit.

The sensory route would use encoders and a separately qualified visual, haptic
or other sensory transducer to deliver specified feedback. Direct neural
writing remains a separate later research interface, with its own device,
measurement, safety and clinical evidence. Direct cortical writing could have
a sensory role without being a peripheral afferent signal. Predicting an
EEG-like pattern would not deliver sensation: the sensory content, interface,
authority and actual measured response would need their own tests, including
conditions, error and delay.

The action route would send separately qualified encoded controls through a
controller to a virtual body, a robotic actuator or a later qualified biological
interface. Action encoders, controllers and actuators have distinct jobs in
that route. A robot command is an action-role output, not an anatomical efferent
nerve. Simulated action, physical actuation, retained learning and biological
causal claims each require their corresponding evidence. EEG and EMG recorders
neither supply sensation nor actuate a body. These sensory and action
subinterfaces belong within the unchanged five programme pillars.

The retained references supply no assessed methods comparison that makes EEG or EMG recorders equivalent to an implanted sensory writer or action device. The missing clinical-methods source is identified as an unresolved provenance gap in Appendix~\ref{app:reproducibility-availability}. No method or capability is inferred from that missing support. Physical writers, transducers and actuators remain unqualified and unpriced, outside the recording-and-virtual-modelling cost scope.

This recording-and-virtual-modelling proposal cannot test physical sensory
delivery or biological actuation. No such physical interface, stimulation,
acquisition or actuation is established here, and no loop, training or
physical-writing result is claimed.

A physical study would need its own qualified interface and procedure. This
proposal supplies no stimulation dose, implantation, electrode-placement,
acquisition or actuation procedure. No such participant study, stimulation, biological actuation or qualified physical training interface has been implemented or tested. Participant, device, safety and clinical evidence would be separate requirements.
% Source registry IDs: context-write-native-01a12f66, context-native-01a12f62, context-policy4

\subsection{Close the loop with evidence}
Return to the cup and test each prediction against the observation it was
asked to explain. The brain-to-body prediction would face withheld observations;
the body-to-recorded-neural estimate would face its own recorded target.
A sensory or action interface would need a further test of what it delivered,
under which conditions, with what error and delay. Two models can agree and
still be wrong; a high reward can have the same problem. Neither agreement nor
a successful prediction would, by itself, identify a unique inverse brain or
establish someone's private experience.

The programme's distinction remains that approximation gets us towards
plausibility and virtualization gets us to precision. Independent measurements,
repeatability, causal checks and stated error bounds must earn that precision.
The original held-out and real-reference controls would test the prediction;
retained-learning controls would test whether a gain remains after useful
guidance ends. Precision is not an automatic result, and the proposal establishes
no perfect biology.

The next experiment must report where each declared prediction works and where
it fails against its independent measured target. Its answer would determine
which part of the proposed loop needs refinement.
% Source registry IDs: context-native-01a12f62, context-write-native-01a12f66, context-precision-01a12f77, context-policy4, context-spec15, context-glossary13, b2b-circular-checklist-original

\subsection{Technical representation and mask}
The information boundary must survive the calculation. A model given tomorrow's
sample or today's withheld target could make a persuasive prediction of an
answer it was never meant to receive. The representation and mask specify
which comparison the proposed attention computation would actually perform.

Learned query, key and value representations give the proposed attention
calculation a precise form. An admitted mask would restrict which observations
may contribute; biological grounding still requires the independent comparison.
\[\operatorname{Attention}(Q,K,V)=\operatorname{softmax}\!\left(QK^{\mathsf{T}}/\sqrt{d_k}+M\right)V.\]
Here Q, K and V are queries, keys and values; M expresses the admitted mask. The dimension $d_k$ is the common feature width of each query and key vector. The equation is methodological context from the supplied technical proposal. Modality, channel, task and time are part of the observation contract. No weights, EEG encoder or gradient route are selected or executed.

\subsection{Permitted observations and causal modality routes}
Declare B=EEG, E=EMG, Y=independently measured movement, C=task/stimulus/context, and P=the source-bound physics model. Training may relate qualified paired observations. Inference uses only its explicitly admitted subset:

\begin{itemize}
\item EEG-only: permitted past/present B and declared context predict a specified   EMG/movement target. Withheld E, Y and future samples cannot enter the predictor.
\item Body-only: permitted E and/or Y predict specified recorded neural features as   a conditional distribution with uncertainty. This is not a unique inverse brain.
\item Hybrid: additional modalities require a separately named route and evaluation;   a hybrid result cannot be reported as EEG-only.
\item Simulation-derived inputs stay simulated. They do not replace independent   observations or inherit the biological validity of a rendered body.
\end{itemize}

Separate training signals from evaluation references. Cycle consistency can make two models agree while both are wrong. Shared cues, movement artefacts, subject identity, session overlap and preprocessing leakage need the original controls. Do not expose final held-out targets to model, physics or reward selection.

\subsection{Representation, supervised learning and policy refinement}
\begin{enumerate}
\setcounter{enumi}{-1}
\item Establish source rights, subjects/tasks, measurement units, calibration,    uncertainty, clock alignment and frozen subject/session/task splits. Audit    pretraining overlap before preprocessing and model selection.
\item Pretrain modality representations on the qualified inventory. LaBraM is one    candidate EEG encoder: neural-tokenizer training, masked neural-code prediction    and downstream fine-tuning. It is not a complete brain-body writer. Code,    checkpoint, dataset and individual-use rights remain separate qualifications.
\item Learn paired supervised mappings with named targets and explicit modality    masks. Compare appropriate simpler, cue-only and artefact baselines. Keep raw,    processed and model-derived values distinct.
\item Qualify physical predictions and the reward verifier against independent    measurements. A successful task, smaller calibrated movement error and valid    physical states are prospective reward ingredients, not completed results.    Declare loss/reward weights on training/development evidence. A computable    objective or high reward does not confer biological ground truth.
\item Where admitted action/transition/reward coverage supports it, evaluate optional    RL and post-training. Offline trajectory labels alone do not establish the    consequences of actions never supported by the data. Conservative Q-Learning    is a methodological candidate, not a safety or performance guarantee here.
\item Fine-tune for the declared task and qualify both directions on frozen final    tests, uncertainty, baselines and every original threshold. Keep failures.
\item Evaluate a version with its weights fixed. Any later update must be a separate, versioned comparison with independently checked results and reproducible saved state.
\end{enumerate}

An illustrative supervised objective would combine declared errors for
specified EMG, movement and recorded-neural targets with applicable physical
constraints and optional consistency terms. Each term needs its units,
measurement uncertainty and target definition before a weight is chosen.
Consistency can check whether the outputs fit together; an independent
observation checks whether they fit the measured outcome.

Backpropagation changes differentiable components of the computer model;
biological feedback is a different process. A physics engine in the path does
not automatically give the learner a gradient. Qualify the actual
OpenSim/MyoSuite rollout and optimization route, or a separately validated
surrogate, rather than assuming gradients through every route. No paired four-source training,
RL rollout or task checkpoint for this proposed loop is reported here.


% Proposed conceptual-article insertion only; existing manuscript untouched.
% Inherit Noto Sans/Fira Math from the selected master.
\section{A larger model still owes us a learning experiment}

A model could predict a hand's path more accurately while teaching its owner nothing. That distinction matters as the programme expands from recording-specific mappings toward a proposed Large Movement Model. The new hypothesis is that greater amounts and greater diversity of qualified neural, muscle and movement data could improve generalization. Brain2Body asks what any resulting prediction would let a trainee do, what information or assistance supplied the gain, and what remains when teaching ends. A scaling result would answer the first modelling question only within its tested domain.

\subsection{Four training sources, five component roles}

The expanded training proposal has at least four source groups: neural observations; muscle and other myographic observations; measured motion and displacement; and data generated by physics-based motion or musculoskeletal simulation. Suitable neural sources may span brain, spinal, peripheral and relevant visceral recordings. Their actual modalities, sites and task roles remain distinct. The present emphasis is motor--sensory. Myoception is the programme's term for the sensory side of musculoskeletal contraction and movement, including proprioception, tactile feedback and other task-relevant sensations linked to action. It marks proposed scope, without establishing a new external scientific standard. That scope supplies neither acquired coverage nor body-wide simultaneous observations, and the four groups are not an exhaustive taxonomy.

Known tactile stimuli, proprioceptive feedback, contact, force and task context contribute different evidence. Stimulus delivery, reported sensation, recorded sensory activity, calibrated force and tracked configuration cannot be substituted without a tested observation relationship. EMG records muscle electrical activity, not direct peripheral nerve or sensory activity; kinematics alone does not measure proprioceptive experience. Each measured, reported, estimated or simulated quantity keeps its provenance.

The simulation group supplies potential training examples as well as constraints: trajectories, state transitions, muscle--tendon or contact-force outputs and declared movement priors. Record the simulator, body-model version, parameters, solver, control representation and underlying sources. Joint-target control must not be called muscle-driven control. Retargeted recordings retain original lineage and transformations; generated states or forces remain simulated. The generator requires calibration against appropriate measured support with uncertainty and training-only role boundaries.

Test its contribution by comparing a fixed-real-data reference with a simulation-augmented arm under matched permitted information, compute and original criteria. Record real-example exposure, simulated exposure and updates in each arm; a fixed real corpus alone does not match training exposure. Count simulated amount and diversity separately from real acquisitions. Held-out real observations must check the result without entering generator calibration, tuning, reward choice or model selection. Wrong dynamics, contacts or source correspondence can make the synthetic contribution fail. Generated data is never independent biological ground truth for the mechanism that produced it. Generator methods, implementation and use conditions remain to be checked; no frontier implementation or dataset is selected here.

Five component roles describe another axis: brain signals, muscle signals, measured movement, a virtual body and a learning model. A virtual body can supply training data and physical constraints, while a learning model can draw from several source groups. Neither count excludes another useful, properly characterized observation. These resources give the later teaching study more hypotheses to test, rather than supplying the missing retained-learning result.

\subsection{Breadth needs a controlled comparison}

Amount and diversity must be varied separately. Repeated windows from one acquisition increase presentations without adding independent subjects, days or tasks. A proposed amount comparison would increase support from a stated source distribution. A diversity comparison would retain a declared amount while changing named aspects of that distribution, with a crossed comparison testing their interaction. Organism, subject, session, movement, body region, modality, apparatus, environment and measurement quality require separate descriptions. Dataset count alone cannot measure them.

The architecture, permitted input information, targets, preprocessing and selection rules would remain fixed. Source rosters, repetitions, updates and measured computing work must accompany the results; equal update counts alone do not establish matched compute. If diversity also changes recording quality or task difficulty, the comparison cannot attribute its outcome to diversity alone. All exact budgets, controls, seeds, uncertainty methods and decisions remain to be registered before a new experiment. This proposal changes neither the adopted continuous WAY-EEG-GAL development design nor any original protected partition.

A candidate amount--error curve, for a fixed diversity condition and evaluation domain, is
\[
E(N)=E_{\infty}+A(N/N_0)^{-\alpha}.
\]
$N>0$ and its fixed reference $N_0>0$ use the same declared training-amount units. $E$ is a chosen nonnegative error; $E_{\infty}$ and $A$ use its units, while $\alpha$ is dimensionless. This decreasing-error candidate requires $E_{\infty}\geq0$, $A\geq0$ and $\alpha>0$: $A>0$ gives strict decrease, while $A=0$ gives constant error. Positive $\alpha$ alone is insufficient. Predictions must also respect any upper bound of the metric throughout the declared fitting and prediction range. The floor is a prospective curve parameter, not an established irreducible biological limit. No exponent is estimated here and no universal law follows. Constant, saturating, increasing, non-monotonic and domain-specific alternatives need prospective comparison and predictive checks at withheld scaling levels. Observed negative transfer remains a result even when this constrained candidate cannot fit it. Related windows cannot supply independent uncertainty replicates. Neural-target, EMG-target, pixel and physical-motion errors retain their own units and denominators.

Generalization would be evaluated separately across excluded subjects, sessions, tasks, datasets and, where mechanically qualified, species. The target-only/no-transfer reference, applicable incorrect correspondence, simpler predictors and original null, shuffle, cue/confound and causal-time controls remain. Negative transfer, worse domains and saturation stay in the record. A better pooled score cannot discharge an original per-output no-harm criterion or erase known pretraining exposure.

\subsection{Build a composite without inventing a shared observation}

FALCON M1-A and AJILE12 illustrate complementary source candidates. The fixed FALCON M1-A release describes simultaneous M1 activity and processed intramuscular EMG from one rhesus macaque; released synchronized motion capture is not established. Its acausally processed EMG targets limit causal real-time interpretation. The small minival subset is drawn from calibration data for submission-format checks and is not independent validation. AJILE12 describes intracranial recordings and two-dimensional, video-estimated upper-body pose from twelve clinical participants. Paired EMG and instrumented three-dimensional mocap are not established there. Its pose-label construction is participant-adapted; particular held-out-day estimator exposure and numerical alignment uncertainty remain unqualified. \href{https://doi.org/10.48324/dandi.000941/0.241029.1405}{FALCON M1-A release}; \href{https://papers.nips.cc/paper_files/paper/2024/file/8c2e6bb15be1894b8fb4e0f9bcad1739-Paper-Datasets_and_Benchmarks_Track.pdf}{FALCON methods}; \href{https://doi.org/10.1038/s41597-022-01280-y}{Peterson et al., 2022}; \href{https://doi.org/10.48324/dandi.000055/0.220127.0436}{AJILE12 release}.

Qualified source-specific pretraining could exploit each relationship. A transfer model would then have to demonstrate a correspondence on actual target measurements. The separate recordings never become one simultaneous neural--muscle--movement event. The release descriptions do not establish downloaded recordings or a trained model. Their exact use conditions, versioned files, integrity, original acquisition relationships, timing, units, earlier exposure and independent partitions still need checking before training. Retargeting between homologous or mechanically analogous structures also requires declared geometry, frames, scale, load, contact, timing and uncertainty. Similar movement cannot establish identical neural coding.

\subsection{Model the inside and the outside of one action}

When a hand closes on an object, external displacement and contact put demands on internal mechanics. The proposed world model would represent specified musculoskeletal states alongside the supported body, objects and environment. NeuroTranslate would connect named observations and targets; MoCap ReComp would provide its separately qualified source-preserving movement composite; MyoSim Replay would simulate declared mechanics. Their proposed combination is the Large Movement Model. Neither that name nor a neural composite implies a measured whole nervous system, equivalent cells across organisms, complete internal organs or unique access to private experience.

The model must distinguish measured neural activity, muscle electrical activity and tracked motion from estimated activation, internal forces, latent neural state and simulated motion. EMG is not direct peripheral nerve recording or automatically calibrated force. A proposed transition model could forecast a distribution over joint configuration, generalized velocity, muscle-model activation and environmental state from permitted past observations, context and declared controls. Separate observation models would connect each sensor to the state it can constrain. An unavailable observation stays unavailable; a conditional reverse prediction cannot identify a unique hidden brain.

Physical tests would compare aligned predicted movement with independent measured motion, and any internal-force claim with appropriate independent force evidence. Joint limits, contacts, loads, energy accounting, friction, gravity and numerical solver behaviour need their own checks. OpenSim supplies established methods for declared musculoskeletal models; it cannot confer biological validity on this proposed assembly. \href{https://journals.plos.org/ploscompbiol/article?id=10.1371/journal.pcbi.1006223}{Seth et al., 2018}. A rendering can show where a model says a muscle pulls. Only the applicable observations can test whether that pull explains the claimed result.

% R4 conceptual-article insertion candidate; no manuscript/source promotion.
% Inherit the selected master's Noto Sans/Noto Sans Mono/Fira Math setup.
\subsection{One muscle pulse, two simulated trajectories}

A moving virtual hand has answered a mechanical question before it has answered a teaching question. In one numerical diagnostic, we compared two trajectories of the installed \texttt{myohand\_r} human right-hand mechanism on a fixed passive torso. A muscle excitation was added to one branch. The other supplied its matched control. This asks whether the specified software input changes the computed state. It provides no neural decoding, measured biological response or retained-learning outcome. The mechanism used \texttt{myo\_sim} 0.2.3, commit \texttt{93b0ca8f4ec90c9899ee7f05fee561e9911da91b}, MuJoCo 3.14.0 and NumPy 2.5.3.\cite{myohandMyoSimPinned2026,myohandMuJoCo3140,mujoco2012,myohandNumericalRecord20261105}

\subsubsection{Mechanism and matched comparison}

The compiled model has 23 scalar hinge degrees of freedom, 39 muscle--tendon actuators, 39 activation states and zero sensors. Both branches began with matching full recorded initial states, including positions, velocities and activation states. The control branch received no additional excitation pulse; the pulse branch received dimensionless excitation 0.2 at \texttt{FDS2\_r}, with all other controls zero. Excitation is a muscle-model input, distinct from a joint-position request, measured EMG or neural activity. Gravity and passive dynamics remain present, so the unpulsed trajectory need not remain stationary.

The native Euler integrator advanced each branch through 100 explicit steps of $0.002\,\mathrm{s}$, giving nominal $T=0.200\,\mathrm{s}$. Pulse steps 10--59 inclusive correspond to the pre-step interval $0.020\leq t<0.120\,\mathrm{s}$. Each branch received explicit initial and final forward evaluations to synchronize state-dependent quantities. Thus the record contains 200 explicit integration steps and four explicit forward evaluations. The executed source path also implies two compilation calls and one construction-time internal forward evaluation. Those counts do not enumerate every internal native solver operation.

The construction path changed physical parameters even with optional \texttt{inertia\_floor=None}: each diagonal inertia component of four fixed scaffold bodies changed from $1$ to $10^{-4}\,\mathrm{kg\,m^2}$. This native transformation is part of the tested mechanism, rather than evidence that all parameters remained untouched. No additional physical retuning or \texttt{implicitfast} fallback was used. Position, orientation, scale, contact behaviour and biological fidelity have not been independently calibrated by this diagnostic.

\subsubsection{Recorded response}

For pulse branch $P$, control $C$ and the synchronized final state, the reported quantities are
\[
\Delta q_{\infty}=\max_{1\leq j\leq23}|q_{P,j}(T)-q_{C,j}(T)|,
\qquad
\Delta a=a_{P,\mathrm{FDS2}}(T)-a_{C,\mathrm{FDS2}}(T),
\]
\[
\Delta F=F_{P,\mathrm{FDS2}}(T)-F_{C,\mathrm{FDS2}}(T).
\]
The first quantity is the maximum absolute separation across the 23 hinge coordinates at the final time, in radians. The other quantities are signed pulse-minus-control contrasts in the selected dimensionless activation and tendon-actuator force in newtons. Force follows the model's convention and uses synchronized endpoint forward evaluations, rather than derived values left over from the last integration step.

% Shared by all four proposed R3 inserts; no added packages or font defaults.
\begin{table}[htbp]
\centering
\caption{Final contrasts for one simulated right-hand mechanism. All values compare the pulse branch with its matched control at nominal $T=0.200\,\mathrm{s}$. The angle value is the largest absolute separation across 23 hinges at that endpoint; activation and force retain the pulse-minus-control sign. Activation and force concern the selected \texttt{FDS2\_r} actuator. Force uses the synchronized final forward evaluations and the model's sign convention. These are simulated contrasts, not errors against biological measurements. The displayed digits preserve the recorded numerical values, without establishing biological accuracy.\cite{myohandNumericalRecord20261105}}
\label{tab:myohand-numerical-response}
\begin{tabular}{p{0.53\linewidth}p{0.26\linewidth}p{0.10\linewidth}}
\hline
Quantity & Contrast & Unit \\
\hline
Largest absolute final joint-coordinate separation, $\Delta q_{\infty}$ & $0.5027347840656727$ & rad \\
Selected final activation difference, $\Delta a$ & $0.05205729904534248$ & none \\
Selected final tendon-actuator force difference, $\Delta F$ & $-6.407722136774943$ & N \\
\hline
\end{tabular}
\end{table}


Independent auditing of the saved evidence completed 132 standard-library checks without a failed check or model rerun. All required recorded numerical channels were finite and recorded native warning counts were zero. Four stored endpoint mass-matrix spectra were positive. Raw mass matrices and residual operands were not independently exported; their stored diagnostics were checked without independently resolving those matrices. No contact was observed. At the final control endpoint, the recorded soft-limit excess was $0.0004121829780273552\,\mathrm{rad}$, so exact hard-bound enforcement is not established. Neither conservation, numerical convergence nor a physiological response threshold follows.

These contrasts establish a response to the added excitation within this simulated configuration. They are not errors against measured human movement, EMG or force. Two dependent trajectories with 200 time steps supply neither 200 independent experiments nor biological replication. No statistical $p$-value is inferred, and numerical precision in the record does not establish biological accuracy or repeatability across cases.

\subsubsection{From a response to a teaching experiment}

The diagnostic establishes the recorded input-to-state response, not a performance benchmark. Its complete computing-resource record and unresolved memory/timing attribution are retained in Appendix~\ref{app:reproducibility-availability}.

The hand's zero sensors leave Myoception---the programme's proposed sensory side of muscle-driven action---without a measurement in this test. Simulated activation, force and motion do not supply afferent recordings, biological feedback, a learned controller, a trained Large Movement Model or generalization. The four proposed source groups remain neural observations, myography, measured motion and physics-generated examples. Mechanics, controls, calibration, uncertainty, source rights and data lineage require qualification before this generator could enter training, and its contribution would need independent held-out real tests. The amount/diversity hypothesis remains untested.

For Brain2Body, even a later qualified controller would leave the teaching question: what does a trainee retain when trainer commands end? That comparison still needs the original matched information, practice and assistance controls, sham/delayed/unrelated guidance, independent outcomes and a constant compensatory movement bridge where required. This diagnostic contains no trainee, teaching, withdrawal or physical intervention. It selects no training generator and supplies no biological ground truth. LINK's eighteen fits and 720 saved epochs retain both failed admission gates and all five failed benefit claims; protected tests and quarantined sources remain closed.

\subsection{Carry the test through teaching and withdrawal}

A better movement prior might improve assistance, but the teaching hypothesis requires another comparison. A prospective trainer--trainee study must identify the model's inputs, expert information, feedback, commands, practice and assistance burden. Match those resources against the appropriate simpler controller and original sham, delayed and unrelated-guidance controls. Otherwise an apparent neural benefit could come from supplying the target action, more practice or a less demanding task.

Retention is assessed after trainer commands stop, with any necessary compensatory movement bridge held constant. Immediate performance, delayed retention and unfamiliar task variations remain separate outcomes. A gain that depends on continued guidance may support assistance. A controlled gain that persists would need its own domain and uncertainty; it would still establish neither tissue repair nor arbitrary skill transfer. Simulation, offline replay, measured biological feedback and physical intervention retain distinct evidence and authority.

No scaling result, trained body-wide world model or biological teaching result is reported here. LINK retains eighteen fits and 720 saved epochs, with both model-admission gates and all five benefit claims failed. A broader programme gives those failures somewhere useful to lead: a new, source-qualified comparison that can show exactly which prediction improves, what supplies the improvement and whether a trainee retains anything from it.

% Prospective assistance/retention protocol; no biological experiment reported.
\section{From a guided action to a retained skill}

Catch a ball with outside guidance, then try again after the guidance ends. The
second catch asks what changed in the trainee. The comparison is useful only if
the task and access to movement remain comparable; removing a necessary movement
bridge along with the trainer would confound the test. This distinction carries
the proposed sequence from touch and object handling to NeuroBoxer and NeuroArcher.

The near-term apparatus would record scalp EEG and surface EMG from sleeves on
both arms, alongside synchronized, calibrated left and right tripod views.
OpenBCI or an equivalent configurable system remains a candidate requiring
channel, site, reference, bandwidth, unit, delay and clock qualification. Each
event needs its independently observed movement and, where claimed, an
appropriate contact or force reference. Surface EMG supplies neither direct
force nor afferent nerve activity. A virtual body can predict a response, but
its predicted contact or sensory state cannot stand in for those observations.

This left/right hand-task apparatus and the earlier opposing top/bottom dot-map
apparatus are different proposed configurations of the motion-capture
relay-composite. Their calibration, frames, visibility, occlusion and timing
must be qualified separately. Preserve the source of each measured point and
the uncertainty of the joined map. The musculoskeletal-simulation physics-engine
also needs task- and organism-specific anatomy and mechanics; lifelike rendering
cannot establish its function.

\subsection{Build coordination through four declared tasks}

\paragraph{Touch and movement.}
Rapid hand rubbing brings contact and movement together. Record the requested
knuckle movement and interlaced-hand stretch, with palms outward and arms
extended, as separate conditions. Matched observations must distinguish
contact-related effects from movement-related effects; the cold-hands metaphor
does not establish a cold-exposure result or identify a unique sensation.

\paragraph{Ball transfer and juggling.}
The proposed hot-potato task passes one juggling ball between the hands, followed
by the preferred two-ball juggling task. One limb's action changes the next
limb's problem. Independently observed release and catch events, trajectories
and timing would test that coordination. Preserve uncertainty when an event or
body part is not observed, and compare with cue- and history-based controls.

\paragraph{Punch discrimination and NeuroBoxer.}
Independently labelled punch classes would test discrimination and movement
prediction before becoming the NeuroBoxer virtual game. Boxing includes hand
and wrist stabilization as well as arm motion. A prime-era Mike Tyson opponent
is a proposed virtual game character. A virtual opponent supplies no physical
opponent or evidence of safety. Rights-qualified expert footage can inform
movement and tactics, but supplies no expert neural, EMG or force recordings.

\paragraph{Toy archery and NeuroArcher.}
Toy archery would add explicit grasp, asymmetric two-hand roles, draw under
opposing loads, release and return. Its proposed arrows have small, cup-shaped
suction tips and are carried in an over-arm pouch with multiple arrows. NeuroArcher would use that coordination to act
on a virtual target. Success still requires a declared outcome and an
independent comparison, rather than agreement between simulated components.

The whole snap is a further coordinated showcase, from arm-out/palm-down rest
through elbow and forearm swing and turn, wrist and hand upward during the snap,
then downward return. The cup reach remains a simpler precision example. The
motion, contact and anatomy needed for either task are specific; a human-looking
animation cannot demonstrate a monkey's ability to perform the snap.

\subsection{Separate motion control, neural contribution and physical writing}

The first game milestone uses body motion to control a virtual body. A later
EEG-only milestone would ask whether permitted brain observations contribute to
the specified action. Withheld EMG and measured movement cannot become predictor
inputs through alignment, filtering or cue leakage. Eye activity, covert muscle
activity, residual movement and shared task information need their own controls.
Keep the original brain-only, body-only and paired comparisons, target-persistence
baseline, partition rules and success criteria.

Aim to make the declared intent-to-action mapping as close to one-to-one as
possible by measuring its intermediary. State the participant, session, condition,
allowed context, tolerated error and independent reference. Reduced ambiguity
would be a task-specific result; it would not identify a unique hidden intention
or establish the best achievable mapping. Reverse reconstruction of recorded
EEG features remains conditional and cannot deliver sensation.

The proposed later physical writer is a separate experiment. EEG and surface
EMG record, whereas neuromuscular electrical stimulation (NMES) and functional electrical stimulation (FES) denote stimulation approaches that can produce
contractions or functional movement. Sensation delivery needs its own writer or
sensory transducer. Device, clinical, consent and physical admission remain
separate, and no programme spinal intervention or qualified spinal-writing interface is reported. These writers and
actuators are unqualified and outside the existing acquisition/virtual-work
budget scope; no recording or intervention is reported here.

\subsection{Ask what remains after the help ends}

A proposed blindfolded condition changes visual feedback. It cannot independently
establish neural contribution, causal control or retained skill. Those claims
need matched comparisons, including a trainer-off test that preserves any movement
bridge needed to perform the task. Measure immediate performance, delayed
retention and transfer to specified unfamiliar conditions separately. Report
remaining support, uncertainty and alternative explanations.

The model's training and the trainee's learning also need separate evidence.
Qualified paired data could fit predictions; optional control-policy reinforcement
learning needs states, actions, transitions and rewards. A simulated reward is
not evidence of biological teaching. Final held-out outcomes remain outside
fitting, fine-tuning, rewards and model selection; frozen inference is distinct
from a separately reviewed adaptive update. The retained LINK comparison failed
both model-admission gates and all five benefit claims. These prospective tasks
do not revise that result or establish retained learning, tissue repair or a cure.

% Prospective apparatus and protocol; scientific and publication acceptance remain open.
\section{Build a brain--body system we can test}

The proposed architecture becomes a first experiment only when its measurements
are specified. Begin with one hand and one movement that can be measured well.
Record the brain and muscle signals, capture the movement, and compare the
virtual hand's prediction from one declared input side with that independent
measurement. The error identifies what the next experiment must resolve.

The apparatus would use non-invasive EEG, surface EMG, a calibrated
motion-capture relay-composite, an open musculoskeletal-simulation physics-engine
with a lifelike view, and an open multimodal model trained on eligible neural
and muscular structure and activity data. The EEG/EMG observation pair and the
movement/physics pair would meet inside one conditional learner. Intent-related
brain signals and muscle activity would supply observations; measured movement
would check what the body actually did. Predictions from either end must answer
to observations outside the model. Form constrains what could happen; measured
function tests the prediction.

\subsection{One movement. Two directions.}

The forward test uses EEG and declared task context to predict supported muscle
activity and movement on withheld recordings. The reverse test uses muscle
activity and movement to predict recorded EEG features within that task. Compare
both predictions with independent measurements. Agreement between the models
cannot substitute for that test.

Movement measures what the body did. It narrows the explanations of brain
signals that remain plausible without uniquely identifying intent or subjective
experience. Retain stimulus, task, behaviour and separately reported context.
Show uncertainty instead of displaying a generated neural pattern as the
observed mind.

Surface EMG records muscle electrical activity. A claimed incoming sensory or
feedback pathway needs evidence of its own. The virtual body cannot supply that
missing observation. Keep recordings, predictions and simulation identifiable
throughout the test.

\subsection{A bridge must account for the spine}

The spinal cord helps shape the response. Muscle output depends on sensory
input, movement context, local circuitry and descending input; some spinal
pathways can retain learned changes. Reproducing a specified spinal function
would therefore require more than forwarding its signal: a brain--body bridge
must account for the processing it aims to reproduce.

A changed response can reflect immediate reflex modulation, lasting spinal
plasticity or peripheral adaptation. Those explanations need separate tests.
A change during the task does not establish retained memory, and a retained
change does not by itself locate that memory. Human reflex conditioning and
constrained mouse adaptation supply evidence for different claims.~\cite{thompson2009retained,thompson2013retained,lavaud2024}
The human studies do not establish brain-independent learning; the mouse
preparation does not demonstrate unrestricted walking or reaching. Our human
comparisons use abstracts and selected results, discussion or methods passages.
Relevant mouse-study passages were also examined in the author manuscript;
equivalence to the final published edition and the complete methods and
statistics have not been established here.

The periphery adds a different question: did the sensory signal itself change?
Passive proprioceptive training supplies a task-specific peripheral comparison;
its psychophysical and spindle-afferent experiments were separate.~\cite{ackerley2022retained}
Our passive-training comparison uses abstracts and selected results, discussion
or methods passages. None of these studies demonstrates universal skill
transfer or complex memory stored in a peripheral axon.

The next comparison must choose a function we can measure. Does the controller
respond appropriately to the supported context? Does an adaptation persist
through the declared test? Could peripheral or mechanical changes explain the
outcome? Keep recordings, estimates and simulated state in distinct visible
layers. Surface EMG cannot supply a missing sensory-nerve or spinal channel;
an explanatory cutaway cannot register recorded neural sites.

Externalization remains the research aim: can a specified spinal function be
reproduced outside its usual pathway? Replacing a computation, stimulating a
pathway, controlling a simulated limb and restoring biological function would
answer different questions. Each needs its own task- and site-specific evidence;
no programme spinal intervention or qualified spinal-writing interface is reported. Anatomical access does not establish
neural function, increased excitability does not establish useful plasticity,
and no fixed percentage of a complex skill can be allocated to brain, spine
or periphery from these proposals.

\subsection{Reading a signal and writing an output}

The first prototype would test the recording and prediction route with EEG,
EMG and independent movement measurements. To deliver feedback or act on a
body, it would also need the separate sensory or action route specified in
Section~\ref{b2-output-interfaces}. That method distinguishes indirect sensory
feedback from direct neural writing and virtual control from robotic or later
biological actuation, with the device, measurement, safety and clinical evidence
each would require.

Physical sensory writers, transducers and body actuators remain excluded from
this acquisition and virtual-work plan, unqualified and unpriced. Their absence
defines what this first prototype could test. The consolidated interface method
also separates those later physical procedures from the current observations: no stimulation, participant acquisition, biological actuation or qualified physical-interface training is reported.

\subsection{The first prototype has to earn the next one}

The lean tier tests whether the selected signals and calibrated movement
reference can be synchronized and used for a bounded task. The serious-prototype tier
improves acquisition, independent references and source-specific bidirectional
evaluation. Expand cohorts, tasks, instrumentation and compute in the research-lab
tier only when earlier evidence justifies the next step.

Equipment, calibration, accessories, storage, integration, independent references, safety and consent work, and contingency would all affect the cost. No complete, dated cost estimate has been established for these tiers. The lean, serious-prototype and research-lab scopes therefore specify different experiments, not demonstrated price points or a complete neural reconstruction. The unfulfilled currency, vendor and quotation requirements remain in the budget-preparation provenance record.

\subsection{Let the model earn the claim}

Which open recordings can support the next model comparison? Their rights,
version, coverage, clock, units and missing modalities determine the answer.
Training would proceed only under an independently admitted protocol, with test
participants, sessions and tasks kept outside model selection. Simple baselines
would give the added model a comparison to beat. Examining artifact leakage,
absent inputs, timing errors, physically invalid solutions and uncertainty would
keep that comparison answerable to its actual conditions. Earlier failures and
reserved recordings remain constraints on it. The new ambition must answer to
the evidence already kept.

Begin with constrained virtual hand or foot control. Robotics, remote operation,
unfamiliar effectors and extra limbs remain later research directions, each
requiring separate physical, latency and safety evidence. A successful movement
prediction supports a movement claim. Biological and perceptual claims require
the corresponding broader evidence.

Rights-cleared images and lifelike interactive models would make the apparatus and its predictions inspectable. The comparison should identify the observation, the prediction and the error, with the conditions that limit each claim. A convincing image cannot supply a missing measurement or demonstrate biological function. The next experiment must make those distinctions visible enough to test.

\section{Testing the Proposal}

A prediction becomes testable when an independent observation can disagree with it. The new same-stream forecast makes that requirement concrete: its 990-update state improves on earlier weights but still loses to persistence in both pooled TRAIN streams. The next measurement must therefore check the proposed source-to-target relation against its appropriate simpler alternative, rather than reward model size or effort. Synthetic examples can support controlled development only with real reference support; they cannot manufacture missing measured biology. The original thresholds and failed trials remain part of that comparison.\par

Interoceptomics, the proposed study joining internal-body sensing with connection maps, may extend the programme to heartbeat, breathing and other internal conditions. It needs distinct measurements and claims; the movement-first rationale does not directly reveal all sensation or private experience. Useful limited interfaces remain worth investigating even if a universal translator never emerges.\par

Return to the guided hand. A correct tune would establish a performance under
specified help; a delayed, controlled trainer-off comparison would ask what the
trainee retained. Better measurements and body-grounded tests might connect
those two achievements. Current results establish neither arbitrary skill
teaching nor complete biological virtualization, and the current development
models fail their benefit gates. The next generation must account for what
moves, what help causes the improvement, what remains without that help and
where the explanation breaks. That is the experiment the ambition still owes
us.\par

\section{Definitions and Lexicon}

The \hyperref[programme-definitions]{Scientific Article definition and lexicon supplement} is derived from the shared Brain2Body registry identified by version and fingerprint at the start of that supplement. The fingerprint identifies the exact definition source; it is not a scientific score. Preferred names, historical lookup labels, proposed scope and implemented software retain their distinct meanings. The Simple English Edition has its own supplement using the same registry. Necessary specialist terms are also explained where used in this article.\par

\section{References and Evidence Status}

Primary scientific sources are linked beside the claims they support. The accompanying bibliography preserves verified earlier identities, including the organoid paper's correction, and adds the ground-truth case-study sources used here. A published experiment supplies different evidence from a dated developer disclosure or an observation of pinned source code; those distinctions govern the claims throughout this article.\par

The accompanying \href{research-evidence.json}{computational results and limits}
describe numerical and software studies that have not undergone peer review.
The companion NeuroTranslation manuscript is also a working draft; neither
source is peer-reviewed experimental support. The historical conceptual figure
explains a relationship rather than supplying new biological evidence. The
edition note records the remaining scientific and publication limits.\par


\begin{figure}[htbp]
\section*{Historical Conceptual Framework Illustration}
\centering
\makebox[\linewidth][c]{\includegraphics[width=160mm]{figures/historical-brain2body-framework.pdf}}
\caption{The solid route follows the earlier recorded human finger-attempt computation: 41/60 intended contacts and 54/54 executions of issued commands. Predicting the intended contact and executing an issued command answer different questions. The biological branch remains proposed. This historical schematic supplies no current evidence of anatomical or functional correspondence, lifelike-3D fidelity or biological validity. The separate macaque development study saved 18 runs and 720 epochs; both model gates and all five benefit claims fail. Neither study establishes biological teaching, full-body mapping or retained mastery.}
\end{figure}
\clearpage
% R11 Source candidate derived from R9; historical R6 Source-only status remains explicit.
\section*{Scope of this working draft}
This working draft retains the computational studies and develops the expanded
training proposal, one recorded numerical hand diagnostic and the separate single-acquisition self-forecast comparison. The hand diagnostic
compares two simulated trajectories; the self-forecast comparison concerns recorded muscle/glove features. Neither supplies a new neural fit,
reserved-data evaluation, biological teaching or validated whole-body model.
The definition supplement retains the shared registry's 200 entries from revision 34 and their distinct proposed scopes; its base and proposed editorial-wording fingerprints identify the two records.

The preceding working draft is publicly available as paper PDFs and manuscript
sources.\footnote{Revision 6 was source-only. The subsequent rendered working drafts include the public revision 9 paper/manuscript-source snapshot, Git commit
\nolinkurl{a86bf81a308a71ccfa98bf2e0d5ec7ac329e9ba7}. Public paper sources do not
establish a public release of validated model or GUI code.}
Public availability does not establish peer review or biological validation.
Whole-manuscript scientific, editorial, citation, rights, font, layout, native
and export assessment remains unfinished. Current vocabulary does not change
the original numerical results or establish new model capabilities.

\section*{Evidence and availability}
\addcontentsline{toc}{section}{Evidence and availability}
The retained studies and the new same-stream forecast constrain the proposed paired experiment. LINK retains 18 fits and 720 saved epochs with both model gates and all five benefit claims failed. The separate 990-update TRAIN forecast improves on its earlier states but fails persistence in both streams; it remains unpromoted. The hand diagnostic is a two-trajectory numerical response, not a neural fit or biological result. No reserved-data evaluation or biological teaching is reported.

Complete independent LINK checkpoint/update calculations, reserved confirmation, whole online/offline GUI checks and applicable anatomical and functional reference comparisons remain unfinished. The original tolerances and reserved recordings are unchanged. Scientific, clinical, physical, reference-series, rights and runtime claims need their own evidence. Public availability does not establish peer review or complete manuscript, citation, rights, font, layout, native or export assessment.

The definition supplement retains all 200 entries from canonical registry revision 34 and the full JSON history. Its reader wording is projected through proposed editorial revision 35; the registry base and candidate fingerprints are printed with the definitions. The proposal changes no preferred name, alias, formal definition, measurement or demonstrated capability. Current vocabulary does not reopen protected recordings.

The preceding selected paper PDFs and manuscript sources are publicly available. CC0 covers owned contributions; third-party material retains its licence and attribution requirements. The public R9 source bundle omitted its linked reading aid and was non-standalone. A local reading-aid file alone does not repair that public link; this source bundle includes a separately versioned aid whose delivery status is recorded with its exact dependency manifest. They do not establish a validated model or GUI code release. The online package, original-source static offline export, four byte counts and unavailable browser comparison are documented in Appendix~\ref{app:reproducibility-availability}. A smaller entry file supplies no measured interaction, loading, memory, crash or GPU improvement.

The requested byline is \texttt{@eliminatenoise}. Author identity, applicable affiliation and formal repository/submission responsibilities remain unresolved; no peer review, journal acceptance or new rights clearance is implied. Retained bounded Logothetis, Tang, quantum-measurement and NIST passages support only their stated uses. Complete claim-by-claim citation assessment remains unfinished.

\clearpage
\phantomsection\label{programme-definitions}
% Dedicated current SOURCE export; historical G26/G29 exports remain outside this cohort.
\section*{Brain2Body Definitions and Lookup Index}
\addcontentsline{toc}{section}{Brain2Body Definitions and Lookup Index}
\noindent Registry snapshot revision 35 (PROPOSED EDITORIAL ONLY UNADOPTED); 200 entries. SHA-256 \texttt{fb80b84f\allowbreak{}00beeaeb\allowbreak{}11e69f13\allowbreak{}be30c8d7\allowbreak{}b33db892\allowbreak{}3f105e4d\allowbreak{}3d2408c5\allowbreak{}03087b7f}. Canonical base revision 34; base SHA-256 \texttt{3fffa79b\allowbreak{}c9f404ca\allowbreak{}a7d4910e\allowbreak{}9e37de5a\allowbreak{}97371fdc\allowbreak{}1ad11bc5\allowbreak{}4155d309\allowbreak{}ec5f5394}.\par
\noindent A measurement, a computer state and a claim about biology need different words. These entries explain the current programme meanings. A proposed name does not establish a working interface, a trained model or a biological result.\par
\noindent Trainer and trainee name the current guidance roles. A machine-learning model can still be a learner; its fitted parameters do not measure a person's retained learning. Historical names and original source quotations remain historical.\par
\noindent Lookup labels help locate an entry. Some name related or contrasting concepts, so they are not interchangeable. The complete JSON projection preserves every base-registry field, alias, source and definition history alongside this snapshot's public editorial and reference profile.\par
\noindent Physical scale and units, computational detail and statistical aggregation remain separate. Missing ranges or measurements stay unspecified. Both paper editions retain the same definitions and scientific scope; explanatory wording may differ.\par

\noindent Where no task-specific physical range, computational detail or statistical grouping is supplied, it remains unspecified. A software role supplies no anatomical scale, and the size of a model supplies no averaging rule.\par

% Stable term ID: a-checked-message-between-software-parts
\begin{samepage}
\subsection*{A checked message between software parts}
\noindent A structured message and interface specify what the values mean and how software parts exchange them.\par
\end{samepage}
\begin{samepage}
\noindent The checks include recording, ordered channels, quantity, timing where supplied, units and body axes. Passing a check does not prove that the caller’s claimed physical source is true.\par\nopagebreak[4]

\end{samepage}


\noindent\textit{Current registry lookup labels:} Translation packet, typed port and API; packet; translation port; typed interface; API; application programming interface\par

% Stable term ID: a-group-of-recorded-neural-features
\begin{samepage}
\subsection*{A group of recorded neural features}
\noindent Several recorded neural features considered together when predicting an outcome.\par
\end{samepage}
\begin{samepage}
\noindent Their joint pattern can help predict movement. That does not establish that the same pattern has the same meaning in another animal, task or day.\par\nopagebreak[4]

\end{samepage}


\noindent\textit{Current registry lookup labels:} Neural population; population decoding; population model\par

% Stable term ID: a-map-of-physical-connections
\begin{samepage}
\subsection*{A map of physical connections}
\noindent A structural map records physical connections. A functional recording measures activity over time, such as which recorded features produced events during a task.\par
\end{samepage}
\begin{samepage}
\noindent The structure can guide a model, but it cannot supply missing time-recorded activity. A drawing, connection map and synchronized functional recording have different evidential roles.\par\nopagebreak[4]

\end{samepage}


\noindent\textit{Current registry lookup labels:} Structural connectome and anatomical prior; connectome; connectomic connectivity; anatomical prior; structural connectivity; structural reference; recorded function; structural versus functional measurement\par

% Stable term ID: a-model-spanning-the-body
\begin{samepage}
\subsection*{A model spanning the body}
\noindent A future model connecting relevant observations and actions across the brain, spinal cord, nerves, muscles and movement.\par
\end{samepage}
\begin{samepage}
\noindent The project does not yet have a matched voluntary recording across that whole chain. Separate datasets cannot fill the missing measurements.\par\nopagebreak[4]

\end{samepage}


\noindent\textit{Current registry lookup labels:} Body-wide functional model; body-wide neural fusion; whole-body neural translator\par

% Stable term ID: a-specified-human-related-feature
\begin{samepage}
\subsection*{A specified human-related feature}
\noindent Studying a defined human-related biological component or feature in a model organism.\par
\end{samepage}
\begin{samepage}
\noindent This can overlap with adding donor cells. One selected immune or neural feature does not make an animal a miniature human or establish human-like cognition.\par\nopagebreak[4]

\end{samepage}


\noindent\textit{Current registry lookup labels:} Humanization; humanized; human-related feature\par

% Stable term ID: ability-that-remains-after-guidance-is-removed
\subsection*{Ability that remains after guidance is removed}
\noindent A learned ability that remains without the training guidance and can be tested again later.\par
\noindent Remove trainer commands while keeping any bridge needed to express the trainee’s own control. Switching off that required route can prevent movement regardless of what the trainee learned.\par

\noindent\textit{Current registry lookup labels:} Retained learning; retention; independent pupil control; retained skill; pupil-alone; necessary bridge held constant; pupil-alone with necessary bridge constant; independent learner control; learner-alone; learner-alone with necessary bridge constant; independent trainee control; trainee-alone; trainee-alone with necessary bridge constant\par

% Stable term ID: action2world
\begin{samepage}
\subsection*{Action2World}
\noindent The hand reaches; the cup moves or stays still. Contact, load and the measured outcome tell us whether the action did what we intended, and that feedback shapes the next attempt.\par
\end{samepage}
\noindent The relationship between an action and its consequences for the body and surrounding environment, including the sensory information those consequences produce.\par
\noindent A predicted consequence is not an observed outcome. Specify the object, reference frame, task, units and uncertainty. Footage supplies movement examples without supplying the performer's neural, muscle or force measurements.\par
\begin{samepage}
\noindent\textit{Organizational level:} whole system; network; signal/feature\par\nopagebreak[4]

\end{samepage}


\noindent\textit{Current registry lookup labels:} action2world; action to world; action-to-world\par

% Stable term ID: activity-state-inside-the-muscle-model
\begin{samepage}
\subsection*{Activity state inside the muscle model}
\noindent The model’s internal state that responds to a command and affects its force.\par
\end{samepage}
\begin{samepage}
\noindent This state can change more slowly than the command. The displayed activation comes from the computer model, not recovered physiological activation from the macaque recording.\par\nopagebreak[4]

\end{samepage}


\noindent\textit{Current registry lookup labels:} Muscle activation; activation; activation dynamics\par

% Stable term ID: adjustment-using-evidence-from-the-new-setting
\begin{samepage}
\subsection*{Adjustment using evidence from the new setting}
\noindent Changing or checking a model using examples from where it will be used.\par
\end{samepage}
\begin{samepage}
\noindent Adjusting a decoder, checking probability scores and measuring a body’s physical properties are different kinds of calibration. The new-day model uses labeled examples.\par\nopagebreak[4]

\end{samepage}


\noindent\textit{Current registry lookup labels:} Calibration; recalibration; calibrated transfer; uncalibrated; physiological calibration; physiological fidelity; subject calibration\par

% Stable term ID: agent2agent
\begin{samepage}
\subsection*{Agent2Agent}
\noindent One software agent passes declared information to another.\par
\end{samepage}
\noindent A proposed link carrying specified information between software agents.\par
\begin{samepage}
\noindent Name the agents, message or representation, direction and tested outcome. Software communication does not by itself supply a biological recording, sensory writer, shared experience or demonstrated skill transfer.\par\nopagebreak[4]

\end{samepage}


\noindent\textit{Current registry lookup labels:} agent2agent; agent-to-agent; agent to agent\par

% Stable term ID: agent2brain
\begin{samepage}
\subsection*{Agent2Brain}
\noindent An agent would give a person feedback during a stated task. Tests would have to distinguish help during that task from learning that remains afterwards.\par
\end{samepage}
\noindent A proposed path from a software agent to a person through specified feedback or a separately qualified neural output interface.\par
\begin{samepage}
\noindent Feedback would reach a person through an ordinary sensory route or a separately qualified direct neural interface. Those routes need different devices and independent tests of what they deliver: feedback on a screen is not direct neural delivery. EEG and EMG record; they do not write. The proposed agent-to-person path leaves the reverse observation and feedback direction intact. Help during a task must be tested separately from learning retained after it ends. No stimulation, biological actuation or qualified spinal-writing interface is demonstrated by this proposed path.\par\nopagebreak[4]

\end{samepage}


\noindent\textit{Current registry lookup labels:} Agent2Brain?; agent2brain; agent2brain?; agent-to-brain; agent to brain\par

% Stable term ID: all-outputs-inside-their-ranges-together
\begin{samepage}
\subsection*{All outputs inside their ranges together}
\noindent How often every output of interest falls within its stated range in the same case.\par
\end{samepage}
\begin{samepage}
\noindent A 90\% range for each separate output does not guarantee 90\% coverage of all outputs at once or of a whole body.\par\nopagebreak[4]

\end{samepage}


\noindent\textit{Current registry lookup labels:} Simultaneous coverage; all-output coverage; whole-body coverage\par

% Stable term ID: an-independent-way-the-model-can-move
\begin{samepage}
\subsection*{An independent way the model can move}
\noindent One independent model movement, such as rotation at a joint. The formal shorthand is DOF.\par
\end{samepage}
\begin{samepage}
\noindent The native arm has seven coordinates and no separate finger movements. The human replay controls two goals while bracing other independent joints.\par\nopagebreak[4]

\end{samepage}


\noindent\textit{Current registry lookup labels:} Degree of freedom; DOF; DOFs; independent coordinate; braced joint\par

% Stable term ID: analysis-timing-and-biological-travel-time
\begin{samepage}
\subsection*{Analysis timing and biological travel time}
\noindent An analysis offset aligns two recorded quantities. A conduction delay concerns how long a signal travels through biology.\par
\end{samepage}
\begin{samepage}
\noindent The 40 ms sensory offset estimates an earlier limb state. The motor target’s 4 ms offset comes from sampling. Neither measures neural conduction time.\par\nopagebreak[4]

\end{samepage}


\noindent\textit{Current registry lookup labels:} S1 alignment offset and conduction delay; 40 ms sensory lag; S1 lag; 4 ms sampling offset; conduction delay; physiological latency\par

% Stable term ID: animal-drawing-and-measured-anatomy
\begin{samepage}
\subsection*{Animal drawing and measured anatomy}
\noindent An illustration explains a proposed animal or body. An empirical anatomical atlas is built from verified observations of anatomy; it cannot be established by a convincing drawing.\par
\end{samepage}
\begin{samepage}
\noindent The antechinus illustration is a proposal, separate from acquired macaque and human sources. The implementation has no matched validated antechinus neural/body atlas and makes no living-animal learning claim from the image.\par\nopagebreak[4]

\end{samepage}


\noindent\textit{Current registry lookup labels:} Antechinus proposal; antechinus; marsupial proposal; hypothetical organism; Proposed marsupial setting; anatomy illustration; illustrative model; empirical atlas; anatomical atlas\par

% Stable term ID: another-brain-area-used-in-the-finger-study
\begin{samepage}
\subsection*{Another brain area used in the finger study}
\noindent A brain area recorded alongside the movement area in the human finger study.\par
\end{samepage}
\begin{samepage}
\noindent The formal name is posterior parietal cortex, or PPC. The study provides attempted-finger labels, not measured trajectories of executed finger movement.\par\nopagebreak[4]

\end{samepage}


\noindent\textit{Current registry lookup labels:} Posterior parietal cortex; PPC; Parietal brain area\par

% Stable term ID: assumed-repair-to-the-native-muscle-model
\begin{samepage}
\subsection*{Assumed repair to the native muscle model}
\noindent A separately named model version replaces unusable zero-rate properties with the software’s default values.\par
\end{samepage}
\begin{samepage}
\noindent The source uses a named Schutte muscle model. Restoring 117 default properties permits numerical tests but does not measure the subject’s muscle response rates.\par\nopagebreak[4]

\end{samepage}


\noindent\textit{Current registry lookup labels:} Schutte muscle properties and default-rate fork; Schutte; zero-rate; assumed-rate fork; class defaults\par

% Stable term ID: attachment-order-and-the-chosen-route-segment
\begin{samepage}
\subsection*{Attachment order and the chosen route segment}
\noindent The ordered connection points and the numbered segment to which a wrapping rule applies.\par
\end{samepage}
\begin{samepage}
\noindent The rule uses the source model’s attachment order. The SDK counts from one and includes both range endpoints; the correction does not fit arbitrary points to an error frame.\par\nopagebreak[4]

\end{samepage}


\noindent\textit{Current registry lookup labels:} Attachment topology and SDK range; attachment topology; one-based inclusive range; humerus-to-ulna segment\par

% Stable term ID: axes-pushes-and-turns
\begin{samepage}
\subsection*{Axes, pushes and turns}
\noindent A coordinate frame states the axes used for movement. A handle wrench combines forces and turning moments measured in that frame.\par
\end{samepage}
\begin{samepage}
\noindent The handle record has three force and three moment components. An unknown mapping into the computer-body axes prevents treating those forces as verified simulator inputs.\par\nopagebreak[4]

\end{samepage}


\noindent\textit{Current registry lookup labels:} Coordinate frame and handle wrench; body frame; XY; physical frame; handle wrench; forces and moments\par

% Stable term ID: biological-tissue-repair
\begin{samepage}
\subsection*{Biological tissue repair}
\noindent Restoration or replacement of biological tissue or its function, supported by direct biological evidence.\par
\end{samepage}
\begin{samepage}
\noindent Better simulated tracking, assisted movement and retained behavior are different outcomes. The prototype does not demonstrate neural tissue repair.\par\nopagebreak[4]

\end{samepage}


\noindent\textit{Current registry lookup labels:} Regeneration; tissue regeneration; repair\par

% Stable term ID: biological-virtualization
\begin{samepage}
\subsection*{Biological virtualization}
\noindent For a movement model, matching the path is one question; matching the living system’s response to an intervention is another.\par
\end{samepage}
\noindent A claim that a digital model reproduces specified biological behaviour and causal dynamics within an independently tested domain and stated error bounds.\par
\noindent Independent ground truth must establish measured error and repeatability within declared conditions and tolerances. For causal dynamics, the comparison also needs independently measured responses to appropriate causal interventions; a similar-looking trajectory alone cannot establish them. Approximation can make a model plausible and may contribute to a validated model when its assumptions, errors and domain are tested and disclosed. A musculoskeletal simulation is itself a model to test, not independent empirical ground truth. Numerical detail and repeated verifier results guarantee neither biological accuracy nor complete physiology, subjective continuity, consciousness transfer or implemented Brain2Body capability. This biological claim is distinct from operating-system virtualization.\par
\noindent\textit{Scope:} validation: Specified observed behaviours and causal dynamical responses under declared conditions.\par
\noindent\textit{Organizational level:} whole system; network; organ; local circuit; cell; signal/feature\par
\noindent\textit{Computational granularity:} Declared model and brain-body interfaces\par
\begin{samepage}
\noindent\textit{Statistical aggregation:} Declared validation aggregation; separate from physical scale\par\nopagebreak[4]

\end{samepage}


\noindent\textit{Current registry lookup labels:} virtualization; virtualisation; biological virtualisation\par

% Stable term ID: body-dimensions-and-resistance-to-motion
\begin{samepage}
\subsection*{Body dimensions and resistance to motion}
\noindent Mass is the amount of matter modeled. Inertia describes resistance to changes in motion; anthropometry means body dimensions.\par
\end{samepage}
\begin{samepage}
\noindent The scaled model keeps assumed source masses and applies geometry-related inertia updates. Its length fit does not identify the animal’s actual masses, strengths or individual bone dimensions.\par\nopagebreak[4]

\end{samepage}


\noindent\textit{Current registry lookup labels:} Mass, inertia and anthropometry; inertia; mass; anthropometry; subject scaling; specimen scale\par

% Stable term ID: body-position-stored-by-the-simulator
\begin{samepage}
\subsection*{Body position stored by the simulator}
\noindent The model’s stored configuration, such as its joint positions.\par
\end{samepage}
\begin{samepage}
\noindent The formal variable name is qpos. Its entries depend on the model and are not automatically measured joint angles. The matching velocity variable, qvel, is explained separately.\par\nopagebreak[4]

\end{samepage}


\noindent\textit{Current registry lookup labels:} Generalized position and velocity; qpos; joint state\par

% Stable term ID: body-model-software-and-its-programming-tools
\begin{samepage}
\subsection*{Body-model software and its programming tools}
\noindent OpenSim provides tools for body models, geometry, scaling and simulation. Its software development kit, or SDK, lets code work with those models.\par
\end{samepage}
\begin{samepage}
\noindent Checks on these tools establish the stated model’s numerical behavior. They do not certify the exact upstream model version or the recorded animal’s anatomy.\par\nopagebreak[4]

\end{samepage}


\noindent\textit{Current registry lookup labels:} OpenSim and its SDK; OpenSim; SDK; software development kit; native SDK; OpenSim 4.5.2\par

% Stable term ID: body-position-feedback
\begin{samepage}
\subsection*{Body-position feedback}
\noindent Proprioception means sensing body position and movement. For example, it helps a body monitor where a limb is. Biofeedback uses bodily information to guide action or learning.\par
\end{samepage}
\begin{samepage}
\noindent The implemented correction uses simulated body position. It does not add measured peripheral nerve signals or show that a biological feedback process taught a retained skill.\par\nopagebreak[4]

\end{samepage}


\noindent\textit{Current registry lookup labels:} Proprioception and biofeedback; proprioception; simulated proprioception; biofeedback; Body-state sensing and feedback; proprioceptive feedback\par

% Stable term ID: body2action
\begin{samepage}
\subsection*{Body2Action}
\noindent The proposed link from body mechanics and sensory interaction to a specified, independently measured task action or outcome.\par
\end{samepage}
\noindent A task outcome makes this proposed link testable: the body’s forces, contacts and timing have to be measured in the relevant context. Body2Action connects body mechanics and sensory interaction to that specified action or outcome; it establishes no working controller. World labels are historical conceptual lookup references, not a strict synonym for every environmental interaction.\par
\noindent\textit{Organizational level:} Declared system or interface\par
\noindent\textit{Computational granularity:} Declared process/interface/model\par
\begin{samepage}
\noindent\textit{Statistical aggregation:} Validation domain stated independently of physical scale\par\nopagebreak[4]

\end{samepage}


\noindent\textit{Current registry lookup labels:} Body to action; Body2World; Body to world\par

% Stable term ID: body2body
% Canonical status: "user_requested_proposed_programme_role"
\subsection*{Body2Body}
\noindent A measured movement can constrain a digital body; a digital body can suggest a movement to test. Body2Body asks which quantities can travel between those settings and whether the resulting movement matches an independent reference.\par
\noindent A proposed research direction relating specified physical-body measurements to digital body models and testing mappings or feedback from those models to declared target bodies.\par
\begin{samepage}
\noindent Physical-to-digital and digital-to-physical mappings name a proposed research direction, not a proven symmetric inverse. Observation-to-model fitting, retargeting and a separately qualified physical output have different requirements. Homologous and mechanically analogous features can support unpaired pretraining or transfer when their mechanics, scale, contacts, timing and missing features are declared. That learned relationship is not an actually measured neural pairing. Keep inferred and simulated quantities distinct from measurements. Biological virtualization still requires independent biological outcomes in a stated domain; no real body actuation, retained skill or whole-organism equivalence follows here.\par\nopagebreak[4]

\end{samepage}


\noindent\textit{Current registry lookup labels:} Body to body\par

% Stable term ID: brain-imaging-related-to-blood-oxygen
\begin{samepage}
\subsection*{Brain imaging related to blood oxygen}
\noindent An imaging method, called fMRI, that measures signals associated with changes in blood oxygenation.\par
\end{samepage}
\begin{samepage}
\noindent It does not record individual electrical spikes. The paper discusses a published model for this kind of image, not a project result showing that brain scans and spikes are interchangeable.\par\nopagebreak[4]

\end{samepage}


\noindent\textit{Current registry lookup labels:} Functional magnetic resonance imaging; fMRI\par

% Stable term ID: brain-pad
\subsection*{Brain pad}
\noindent A pad shows the proposed interface at the brain.\par
\noindent The proposed brain-side (cortical) interface, drawn as a pad.\par
\noindent A label for the proposed cortical interface in the programme's illustrations, not a general clinical synonym. The name establishes no acquired neural data, registration, tissue access, function, safety, clinical use, species or physical dimensions.\par

\noindent\textit{Current registry lookup labels:} brain pad; cortical prosthesis; cortical interface; cortical interface prosthesis\par

% Stable term ID: brain-stimulation
\begin{samepage}
\subsection*{Brain stimulation}
\noindent An intervention that applies a specified physical stimulus to alter brain activity.\par
\end{samepage}
\noindent The method must be named, such as transcranial magnetic/electrical stimulation or intracortical microstimulation. These differ in access, risks and capabilities. EEG records; a scalp EEG headset does not provide general sensory writing or safe skill upload. Broad neuromodulation is not identical to precise task-content encoding.\par



\noindent\textit{Source:} A brain-computer interface that evokes tactile sensations improves robotic arm control\par
\noindent\textit{Source identifier:} flesher2021\par
\noindent\textit{DOI:} 10.1126/science.abd0380\par
\noindent\url{https://pubmed.ncbi.nlm.nih.gov/34016775/}\par
\noindent Research precedent for motor decoding plus sensory feedback through intracortical microstimulation; invasive and not a consumer EEG capability\par


\noindent\textit{Current registry lookup labels:} Neural stimulation; Neurostimulation\par

% Stable term ID: brain2agent
\begin{samepage}
\subsection*{Brain2Agent?}
\noindent A recording supplies observations that an agent may use for a stated task.\par
\end{samepage}
\noindent A proposed path from specified brain observations to a software agent through declared recording and interpretation interfaces.\par
\begin{samepage}
\noindent Recorded brain observations would reach an agent through a declared recording and interpretation interface. Where applicable, feedback toward the agent remains distinct from Agent2Brain’s proposed agent-to-person path. Neither direction identifies unique private intent or experience. Recording supplies observations; inference and simulation interpret or model them, while actual sensory delivery requires its own interface. Brain2Agent? remains an open proposal.\par\nopagebreak[4]

\end{samepage}


\noindent\textit{Current registry lookup labels:} Brain2Agent; brain2agent; brain2agent?; brain-to-agent; brain to agent\par

% Stable term ID: brain2body
\begin{samepage}
\subsection*{Brain2Body}
\noindent Connecting brain and body function; the 2 means to.\par
\end{samepage}
\noindent The numeral 2 is the preferred spelling. The programme includes body-to-brain feedback as well as outward pathways; the name does not establish complete biological translation or novelty of body–brain research.\par
\begin{samepage}
\noindent\textit{Organizational level:} whole system; network; local circuit; signal/feature\par\nopagebreak[4]

\end{samepage}


\noindent\textit{Current registry lookup labels:} brain2body; Brain2body; brain 2 body; brain-to-body; brain to body\par

% Stable term ID: brain2brain
\subsection*{Brain2Brain}
\noindent Information links between two brain or neural systems.\par
\noindent The project’s label for studying specified information links between brain or neural systems, including trainer–trainee comparisons.\par
\noindent Brain2Brain is the numeral-style public label for the project’s existing brain-to-brain research direction. Each experiment must specify its signal, task and controls; the name itself establishes no shared thoughts, general skill transfer or successful biological teaching.\par
\noindent\textit{Organizational level:} whole system; network; local circuit; signal/feature\par

\noindent\textit{Current registry lookup labels:} brain2brain; brain 2 brain; brain-to-brain; brain to brain\par

% Stable term ID: brain2motion
\begin{samepage}
\subsection*{Brain2Motion}
\noindent The proposed movement pathway connecting Brain2Body with Body2Action and the resulting measured motion, including the returning sensory feedback.\par
\end{samepage}
\noindent The label concerns the movement part of the wider programme, not every sensory/cognitive function. Brain/spinal/peripheral observations, muscle activity, body mechanics and external outcomes stay distinct. A decoded motion signal, simulation or lifelike animation is not a complete physiological explanation or validated skill transfer.\par
\noindent\textit{Organizational level:} Declared system or interface\par
\noindent\textit{Computational granularity:} Declared process/interface/model\par
\begin{samepage}
\noindent\textit{Statistical aggregation:} Validation domain stated independently of physical scale\par\nopagebreak[4]

\end{samepage}


\noindent\textit{Current registry lookup labels:} Brain to motion\par

% Stable term ID: brain2spine
\begin{samepage}
\subsection*{Brain2Spine}
\noindent A movement command meets a spinal cord that already regulates reflexes, coordinates muscles and adapts through experience. Understanding that middle step helps explain how the same intention can produce different actions.\par
\end{samepage}
\noindent The pathway from brain activity into spinal processing: descending signals influence circuits that combine them with sensory input and the body's current state.\par
\noindent A scalp EEG recording does not identify every descending signal or measure the spinal circuit's hidden state. No measured hidden spinal state or independently qualified spinal-stimulation interface follows from this pathway name.\par
\begin{samepage}
\noindent\textit{Organizational level:} whole system; network; signal/feature\par\nopagebreak[4]

\end{samepage}


\noindent\textit{Current registry lookup labels:} brain2spine; brain to spine; brain-to-spine\par

% Stable term ID: brainnet
\begin{samepage}
\subsection*{BrainNet}
\noindent A network of specified links between brain or neural systems.\par
\end{samepage}
\noindent A working network metaphor for connected brain or neural systems within the programme.\par
\noindent BrainNet pictures specified links between brain or neural systems as a network. It is a programme metaphor; it establishes no validated universal neural network, uniquely coined field or available trademark.\par
\begin{samepage}
\noindent\textit{Organizational level:} whole system; network; local circuit; signal/feature\par\nopagebreak[4]

\end{samepage}


\noindent\textit{Current registry lookup labels:} brainnet; Brain Net\par

% Stable term ID: brain-body-interface
% Canonical status: "user_requested_proposed_programme_role"
\subsection*{Brain–Body Interface}
\noindent A brain recording and a measured hand movement answer different questions about the same reach. A Brain–Body Interface would connect the stated observations and any separately qualified outputs so that the proposed relationship can be tested.\par
\noindent A proposed interface linking specified brain observations with body observations or control outputs in a declared task, with each direction and physical route identified.\par
\begin{samepage}
\noindent This names an interface role, rather than the Brain2Body programme or a model spanning the body. A Digital Bridge is the existing selected-information route around an interruption or between specified recording and control channels; it is related, not an automatic synonym for every Brain–Body Interface. Recording, inference, simulated control and physical output require distinct evidence. EEG and EMG record; a feedback or stimulation output needs its own qualified interface. The label establishes neither complete coverage nor a bidirectional biological loop. No new acronym is introduced.\par\nopagebreak[4]

\end{samepage}


\noindent\textit{Current registry lookup labels:} Brain-Body Interface\par

% Stable term ID: choosing-a-model-and-deciding-if-it-passes
\begin{samepage}
\subsection*{Choosing a model and deciding if it passes}
\noindent Choosing settings and applying a stated pass rule using the validation data.\par
\end{samepage}
\begin{samepage}
\noindent The frozen-transfer rule requires a prediction score at least 0.01 above every specified control. A failing candidate stays rejected rather than being given an easier success rule.\par\nopagebreak[4]

\end{samepage}


\noindent\textit{Current registry lookup labels:} Validation selection and admission; validation-selected; acceptance criterion; admission; failure criterion\par

% Stable term ID: command-sent-to-a-modeled-muscle
\begin{samepage}
\subsection*{Command sent to a modeled muscle}
\noindent The input command that starts the modeled muscle’s activation process.\par
\end{samepage}
\begin{samepage}
\noindent The command is bounded and nonnegative here. It differs from the resulting activity state, recorded muscle electricity and instructions for a biological intervention.\par\nopagebreak[4]

\end{samepage}


\noindent\textit{Current registry lookup labels:} Muscle excitation; excitation; actuator command; muscle control\par

% Stable term ID: communication-route-around-an-interruption
\begin{samepage}
\subsection*{Communication route around an interruption}
\noindent An interface carrying selected information around a modeled break or between stated recording and control channels.\par
\end{samepage}
\begin{samepage}
\noindent A bridge can compensate without repairing tissue or supplying muscle energy. The virtual interruption does not establish a biological injury model or treatment.\par\nopagebreak[4]

\end{samepage}


\noindent\textit{Current registry lookup labels:} Digital bridge; compensatory bridge; brain–spinal bridge; BCI; brain-computer interface; virtual pathway interruption\par

% Stable term ID: comparable-roles-in-different-organisms
\begin{samepage}
\subsection*{Comparable roles in different organisms}
\noindent Comparing roles such as reaching or support while keeping each organism’s anatomy explicit.\par
\end{samepage}
\begin{samepage}
\noindent Related roles do not imply identical neurons or muscles. A combined fictional body cannot replace the measured sources and checks needed for each real organism.\par\nopagebreak[4]

\end{samepage}


\noindent\textit{Current registry lookup labels:} Organism role and homology; organism adapter; limb role; functional coordinate; homology; biped/quadruped\par

% Stable term ID: comparisons-that-test-whether-guidance-matters
\subsection*{Comparisons that test whether guidance matters}
\noindent Matched replay uses the relevant recorded activity under matched task conditions. Stale or unmatched replay uses activity from a different time or condition. No-link and inactive or sham comparisons test whether the proposed information route adds benefit.\par
\noindent These are proposed comparison types, not new animal results. Shuffling disrupts timing; a yoked comparison gives matched help from another case. Guidance amount, timing and later trainee control need their own stated comparisons.\par

\noindent\textit{Current registry lookup labels:} Shuffled and yoked controls; shuffled-input control; time-shuffled; yoked signal comparison; Matched comparisons with altered information; matched replay; matched own-session activity; stale replay; unmatched replay; no-link; sham; inactive comparison\par

% Stable term ID: completed-time-interval-for-counting-events
\begin{samepage}
\subsection*{Completed time interval for counting events}
\noindent A stated time interval groups events into one count. A half-open interval includes its start and excludes its end.\par
\end{samepage}
\begin{samepage}
\noindent Inputs use completed intervals. The exact boundary rule, interval duration and timing must remain explicit rather than being guessed from the display.\par\nopagebreak[4]

\end{samepage}


\noindent\textit{Current registry lookup labels:} Completed time bin and half-open binning; bin; 20 ms bin; 50 ms bin; half-open; bin end; cadence\par

% Stable term ID: computer-processors-memory-and-file-size
\begin{samepage}
\subsection*{Computer processors, memory and file size}
\noindent A CPU is a general processor, a GPU is another kind of parallel processor, and RAM is working memory. Bytes and GiB measure digital data size.\par
\end{samepage}
\begin{samepage}
\noindent The training described in this source edition used four CPU threads; GPU acceleration was neither verified nor used. AI-assisted coding does not establish independent biological validation or an authorship credential.\par\nopagebreak[4]

\end{samepage}


\noindent\textit{Current registry lookup labels:} Computing hardware and data size; CPU; GPU; RAM; AMD; AI; Ryzen AI Max+395; GiB; byte; bytes\par

% Stable term ID: correct-label-rate-and-chance-comparison
\begin{samepage}
\subsection*{Correct-label rate and chance comparison}
\noindent Accuracy is the proportion of labels chosen correctly. A confusion table shows which labels were mistaken for others.\par
\end{samepage}
\begin{samepage}
\noindent A uniform guess among ten balanced classes has a 10\% reference rate. Accuracy and label scores do not prove executed movement or confidence in biological truth.\par\nopagebreak[4]

\end{samepage}


\noindent\textit{Current registry lookup labels:} Classification accuracy and chance; accuracy; uniform chance; ten-class confusion; confusion matrix\par

% Stable term ID: count-prediction-improvement-per-recorded-event
\begin{samepage}
\subsection*{Count-prediction improvement per recorded event}
\noindent Bits per spike describes how much a count-prediction score improves over a defined baseline, divided by the number of recorded spikes.\par
\end{samepage}
\begin{samepage}
\noindent The baseline uses the training mean rate. A positive result concerns estimated average counts, not newly recorded activity or verified stimulation instructions.\par\nopagebreak[4]

\end{samepage}


\noindent\textit{Current registry lookup labels:} Bits per spike; BPS; ΔBPS; bits/spike; likelihood gain\par

% Stable term ID: data-for-learning-choosing-and-checking
\begin{samepage}
\subsection*{Data for learning, choosing and checking}
\noindent Separate portions used to fit the model, choose its settings, and check the selected result.\par
\end{samepage}
\begin{samepage}
\noindent The paper uses local time-ordered splits. If test results have already been inspected, later reuse cannot be described as another independent test.\par\nopagebreak[4]

\end{samepage}


\noindent\textit{Current registry lookup labels:} Training, validation and test partitions; train/validation/test; train/val/test; split; chronological 70/15/15\par

% Stable term ID: data-kept-out-of-a-stated-step
\begin{samepage}
\subsection*{Data kept out of a stated step}
\noindent Data reserved from training or from another specified development step.\par
\end{samepage}
\begin{samepage}
\noindent Independence depends on what was actually excluded. A previously inspected case does not become a fresh test just because it is called held-out.\par\nopagebreak[4]

\end{samepage}


\noindent\textit{Current registry lookup labels:} Holdout; held-out; reserved window; independent evaluation\par

% Stable term ID: describing-motion-and-calculating-what-causes-it
\begin{samepage}
\subsection*{Describing motion and calculating what causes it}
\noindent Describing a body’s position and movement is one task; calculating how modeled forces change that movement is another.\par
\end{samepage}
\begin{samepage}
\noindent The paper calls these kinematics and dynamics. A replayed pose and a muscle-driven trajectory are different results; neither alone proves physiological accuracy.\par\nopagebreak[4]

\end{samepage}


\noindent\textit{Current registry lookup labels:} Kinematics and dynamics; kinematics; dynamics; forward dynamics; numerical plant\par

% Stable term ID: ecotype
% Canonical status: "established_scientific_reference"
\subsection*{Ecotype}
\noindent A population adapted to a coastal habitat is an example of the ecological distinction at issue. Ecotype names the adapted population or form; the surroundings themselves are its environment.\par
\noindent A locally adapted population or form associated with a particular ecological setting.\par
\noindent Ecotype is an established population/ecological concept, not an individual’s surroundings or an obligatory third stage after genotype and phenotype. State the organism and evidence for the claimed local differentiation or adaptation. The retained Turesson source is a historical plant example~\citep{turesson1922ecotype}, not a complete modern evolutionary theory or evidence about the proposed nervous-system experiments.\par

% Stable term ID: electrical-recording-in-outer-brain-tissue
\begin{samepage}
\subsection*{Electrical recording in outer brain tissue}
\noindent Electrodes in the brain’s outer layer, called the cortex, record electrical signals outside the cells.\par
\end{samepage}
\begin{samepage}
\noindent This differs from a scalp recording. It also does not observe every cell in a circuit; the source file defines which events and channels were recorded.\par\nopagebreak[4]

\end{samepage}


\noindent\textit{Current registry lookup labels:} Intracortical extracellular recording; intracortical; extracellular cortical events; Electrical recording inside the cortex; cortex\par

% Stable term ID: enactome
% Canonical status: "user_requested_proposed_data_organization"
\subsection*{Enactome}
\noindent A reach has a path, a time and a result. An Enactome would organize observations like these so that the action and what followed can be compared across stated conditions.\par
\noindent A proposed organization of specified action observations and the recorded relationships between actions and their outcomes in declared tasks.\par
\noindent Declare which actions, body parts, tasks, units, frames and times the data cover. Keep commands, measured movement, sensory consequences and simulated outcomes separate. An observed association does not establish causation, a complete repertoire or the person’s unique intent. This proposed data organization differs from the tentative Enactotype individual-profile label.\par

% Stable term ID: endohumanisation
% Canonical status: "user_requested_proposed_programme_role"
\subsection*{EndoHumanisation}
\noindent Choose a particular human-related feature, then ask what introducing or studying that feature would change in the model organism. EndoHumanisation names that proposed line of work.\par
\noindent A proposed research stage studying specified human-related genetic features within a model organism and testing their relevance to a declared biological question.\par
\noindent The roadmap places XenoTransgenesis at step 1.1. The working interpretation concerns defined genetic features; the actual manipulation and its established scientific name remain to be specified. A selected feature does not make an animal human or establish human cognition, disease fidelity or an appropriate intervention age. Biological relevance, function and animal/participant, ethical, safety and rights requirements need their own evidence and admission. This entry supplies no experimental procedure.\par

% Stable term ID: error-range-based-on-validation-examples
\begin{samepage}
\subsection*{Error range based on validation examples}
\noindent A range placed around a prediction using the errors seen on the data used for model selection.\par
\end{samepage}
\begin{samepage}
\noindent The ranges have constant width for each output. They do not guarantee future coverage when examples overlap in time or their errors are related.\par\nopagebreak[4]

\end{samepage}


\noindent\textit{Current registry lookup labels:} Residual band; empirical residual bounds; validation residual quantile; 90\% per-output residual band\par

% Stable term ID: events-grouped-as-one-neural-feature
\begin{samepage}
\subsection*{Events grouped as one neural feature}
\noindent The source recording process groups events into features that the model can count separately.\par
\end{samepage}
\begin{samepage}
\noindent A group is not a verified lifelong identity for one neuron. It must stay separate from unassigned events and from the channel labels used in other recordings.\par\nopagebreak[4]

\end{samepage}


\noindent\textit{Current registry lookup labels:} Sorted neural unit; sorted unit; sorted cell; sorted feature\par

% Stable term ID: exohumanisation
% Canonical status: "user_requested_proposed_programme_role"
\subsection*{ExoHumanisation}
\noindent Added tissue can survive without doing the job expected of it. ExoHumanisation names a proposed study of a specified human-related graft: what survives, what connects, and what function can actually be measured?\par
\noindent A proposed research stage studying specified human-related donor cells or tissue in a model organism and testing their survival, integration and function.\par
\noindent The roadmap associates step 2.1 with XenoTransplantation. A cross-species graft retains that established meaning; cell survival, anatomical connection, activity and useful functional integration are distinct outcomes. A graft does not establish a human nervous system, a shared mind or whole-organism human equivalence. Species, tissue identity, biological purpose, ethics, safety and rights remain explicit dependencies. This label provides no transplantation procedure or intervention-age window.\par

% Stable term ID: expected-event-count-and-its-fit-score
\begin{samepage}
\subsection*{Expected event count and its fit score}
\noindent A model estimates the average event count for given inputs, then a score checks how that estimate fits the recorded counts.\par
\end{samepage}
\begin{samepage}
\noindent The paper uses Poisson negative log likelihood for that check. A predicted average may be fractional; a good count-fit score does not establish a complete biological spike model.\par\nopagebreak[4]

\end{samepage}


\noindent\textit{Current registry lookup labels:} Conditional count mean and Poisson likelihood; expected count mean; conditional count; Poisson; negative log likelihood; NLL\par

% Stable term ID: external-movement-model
\begin{samepage}
\subsection*{External movement model}
\noindent The external part of the body model, constrained by independently measured surface movement and pose in a declared task and coordinate frame.\par
\end{samepage}
\noindent The motion-capture relay-composite retains the source views, clocks, calibration and coordinate frame for the declared task. The earlier upper/lower hand concept and the later synchronized left/right tripod arm proposal are different arrangements; each needs its own visibility, synchronization and error checks. The composite is derived. Hidden internal states, forces and intent are not uniquely measured. LiDAR and optical infrared tracking require distinct hardware, calibration and uncertainty qualification.\par

% Stable term ID: extra-correction-and-its-scoring-rule
\begin{samepage}
\subsection*{Extra correction and its scoring rule}
\noindent A model can add a correction to an existing controller. A reward scores the result; command-change rate describes how quickly the input changes.\par
\end{samepage}
\begin{samepage}
\noindent The formal names include residual policy, reinforcement learning and actuator slew. A simulator can score tracking or task completion without proving biological learning or correct hidden neural activity.\par\nopagebreak[4]

\end{samepage}


\noindent\textit{Current registry lookup labels:} Residual policy, reward and actuator slew; residual policy; reinforcement learning; RL; verifiable reward; actuator slew; task completion\par

% Stable term ID: file-fingerprints
\begin{samepage}
\subsection*{File fingerprints}
\noindent A checksum or hash is calculated from file bytes to detect changes or match a stated reference.\par
\end{samepage}
\begin{samepage}
\noindent The named methods here include MD5 and SHA-256. A match checks bytes against that reference; it does not prove scientific truth or verify an entire file when only a range was acquired.\par\nopagebreak[4]

\end{samepage}


\noindent\textit{Current registry lookup labels:} Checksums and hashes; checksum; hash; MD5; SHA-256; source checksum; archive hash\par

% Stable term ID: force-in-the-modeled-muscle-and-tendon
\begin{samepage}
\subsection*{Force in the modeled muscle and tendon}
\noindent The force produced by the computer model for a muscle-tendon unit, or a force measured by a stated physical method.\par
\end{samepage}
\begin{samepage}
\noindent The native test reads forces from the running simulator rather than an old lookup table. Those values are still simulated rather than isolated forces recorded in the source animal.\par\nopagebreak[4]

\end{samepage}


\noindent\textit{Current registry lookup labels:} Tendon force; muscle force; isolated muscle force\par

% Stable term ID: form-function-action
\begin{samepage}
\subsection*{Form \(\rightarrow\) function \(\rightarrow\) action}
\noindent The programme narrative follows structural constraints, the context-dependent processes they support, and independently measured actions or outcomes.\par
\end{samepage}
\noindent The three words organize the explanation; every molecular and neural level has its own structure and dynamics. Form informs and constrains function but does not determine the full process alone. Feedback returns outcomes to the system; this is not a one-way causal chain or a guaranteed one-to-one mapping.\par
\noindent\textit{Organizational level:} Declared system or interface\par
\noindent\textit{Computational granularity:} Declared process/interface/model\par
\begin{samepage}
\noindent\textit{Statistical aggregation:} Validation domain stated independently of physical scale\par\nopagebreak[4]

\end{samepage}


\noindent\textit{Current registry lookup labels:} Form to function to action\par

% Stable term ID: functional-electrical-stimulation
\begin{samepage}
\subsection*{Functional electrical stimulation}
\noindent Electrical stimulation coordinated to produce a specified functional movement or task.\par
\end{samepage}
\noindent FES is the appropriate functional-actuation label for the proposed later assisted-body study. It requires separately qualified devices, control/safety limits and specialist-supervised research approval. This registry contains no stimulation settings or self-experiment instructions.\par



\noindent\textit{Source:} Restoring cortical control of functional movement in a human with quadriplegia\par
\noindent\textit{Source identifier:} bouton2016\par
\noindent\textit{DOI:} 10.1038/nature17435\par
\noindent\url{https://www.nature.com/articles/nature17435}\par
\noindent Research precedent linking implanted cortical recording to a custom neuromuscular stimulation sleeve; not an EEG/OpenBCI or elite-skill result\par


\noindent\textit{Current registry lookup labels:} FES\par

% Stable term ID: further-training-for-the-target-task
\begin{samepage}
\subsection*{Further training for the target task}
\noindent Continuing model training so it predicts the desired target, using the relevant inputs and labels.\par
\end{samepage}
\begin{samepage}
\noindent The paper calls this fine-tuning. Changing software weights is not evidence that an animal or person learned a biological skill.\par\nopagebreak[4]

\end{samepage}


\noindent\textit{Current registry lookup labels:} Fine-tuning; paired-output fine-tuning\par

% Stable term ID: future-testing-and-analysis-of-existing-data
\begin{samepage}
\subsection*{Future testing and analysis of existing data}
\noindent A future test checks later unseen outcomes under a fixed plan; an existing-data test analyzes outcomes already recorded.\par
\end{samepage}
\begin{samepage}
\noindent The second source file was chosen before inspection, but its recording day is earlier than the original model’s day. Choosing a file in advance does not make it future biological deployment.\par\nopagebreak[4]

\end{samepage}


\noindent\textit{Current registry lookup labels:} Prospective and retrospective evaluation; prospective; retrospective; prospective biological deployment\par

% Stable term ID: gap-between-data-portions
\begin{samepage}
\subsection*{Gap between data portions}
\noindent Leaving out an interval around a split boundary so overlapping windows do not cross into another portion.\par
\end{samepage}
\begin{samepage}
\noindent The continuous pipeline uses a one-second gap and checks the input and target support. This protects the software split, not every assumption about the original measurement.\par\nopagebreak[4]

\end{samepage}


\noindent\textit{Current registry lookup labels:} Partition embargo; embargo; boundary embargo\par

% Stable term ID: genotype
% Canonical status: "established_scientific_reference"
\subsection*{Genotype}
\noindent A genetic description can record which versions of a DNA sequence an organism carries. Genotype names that genetic make-up, whether the description covers a selected part of the DNA or a wider set. The scope must be stated; the genetic description alone does not tell us every observable trait.\par
\noindent The genetic constitution being described for an organism, with the genetic scope of the description stated.\par
\noindent Genotype is an established scientific term. Its description is not a complete prediction of the organism’s observed traits; phenotype also depends on environment and development. The three labels in the user’s organism-level example are not three obligatory sequential stages or three fixed individual types.\par

% Stable term ID: hallucination
\begin{samepage}
\subsection*{Hallucination}
\noindent A model can predict a movement correctly while inventing the biological mechanism it claims produced it.\par
\end{samepage}
\noindent Invented or ungrounded model content presented as if it corresponded to reality. In programme claims, this includes fabricated biological structure, dynamics or ability presented as established biological fidelity.\par
\noindent Research on AlphaFold3 uses the term for generative-model errors; this programme also applies it to unsupported biological attributions. A disclosed approximation, hypothesis, counterfactual or unvalidated candidate is a different kind of statement, and missing evidence alone does not prove it false. Numerical error concerns the calculation; objective misspecification concerns what is being optimized; causal attribution error concerns the mechanism being credited. Those differences matter when deciding what a result supports. An ungrounded generated answer can happen to be correct without validating its claimed biological mechanism or reliability elsewhere. Informed predictions therefore still require independent performance evaluation. The rolling-dice analogy describes reliance on unverified claims, not literal randomness or equal odds. Calling every approximation a hallucination, or every connectome or simulation fictitious, would erase the distinction.\par
\noindent\textit{Scope:} usage: Model-generated biological content and claims of biological fidelity.\par
\begin{samepage}
\noindent\textit{Computational granularity:} Model output and its interpretation\par\nopagebreak[4]

\end{samepage}


\noindent\textit{Current registry lookup labels:} hallucination; model hallucination\par

% Stable term ID: help-during-current-performance
\subsection*{Help during current performance}
\noindent Help from a controller, replay, trainer or interface that improves performance while it is present.\par
\noindent For example, a trainee can move successfully while following trainer commands. This does not show independent learning or tissue regeneration.\par

\noindent\textit{Current registry lookup labels:} Assistance; compensation; assisted performance\par

% Stable term ID: how-a-muscle-s-pull-turns-a-joint
\begin{samepage}
\subsection*{How a muscle’s pull turns a joint}
\noindent To relate a muscle’s pulling force to joint turning, the model uses its geometry. Consistent work calculations also relate those forces to small coordinate changes.\par
\end{samepage}
\begin{samepage}
\noindent The formal terms are moment arm and virtual work. Their numerical checks do not measure biological muscle strength; the human controller uses them to account for dependent shoulder movements.\par\nopagebreak[4]

\end{samepage}


\noindent\textit{Current registry lookup labels:} Moment arm and virtual work; moment arm; virtual-work projection; dependent shoulder-girdle load\par

% Stable term ID: how-many-labeled-examples-were-supplied
\begin{samepage}
\subsection*{How many labeled examples were supplied}
\noindent The labeled examples given to train or choose a model, including examples from prior recordings.\par
\end{samepage}
\begin{samepage}
\noindent The first 5\% counts new training labels only. Validation labels, earlier-source labels and extra neural samples are additional information.\par\nopagebreak[4]

\end{samepage}


\noindent\textit{Current registry lookup labels:} Label budget and supervision; label budget; paired supervision; adaptation labels; validation labels; source-task supervision\par

% Stable term ID: how-scores-from-several-outputs-are-combined
\begin{samepage}
\subsection*{How scores from several outputs are combined}
\noindent Scores can be averaged across individual outputs or first combined within defined groups. In this entry, macro describes the averaging and weighting rule, not the size of a body part, the scope of a nervous system, or the level of detail in a computer model.\par
\end{samepage}
\begin{samepage}
\noindent The sensory result gives equal weight to five physical output groups. This avoids letting a larger group dominate but does not make their different physical units interchangeable.\par\nopagebreak[4]

\end{samepage}


\noindent\textit{Current registry lookup labels:} Per-output and group-macro scoring; per-output; group-macro; macro R²\par

% Stable term ID: how-uncertain-the-result-remains
\begin{samepage}
\subsection*{How uncertain the result remains}
\noindent The remaining doubt about a prediction, including why it could be wrong or change.\par
\end{samepage}
\begin{samepage}
\noindent Noise in the instrument, errors in the model, changing conditions and missing observations are different causes. The current ranges do not separate all of them.\par\nopagebreak[4]

\end{samepage}


\noindent\textit{Current registry lookup labels:} Prediction uncertainty; uncertainty; model uncertainty; measurement noise; drift\par

% Stable term ID: human-arm-model-used-for-playback
\begin{samepage}
\subsection*{Human arm model used for playback}
\noindent MyoArm is the human arm model used for the displayed muscle-driven replay.\par
\end{samepage}
\begin{samepage}
\noindent The replay uses 63 modeled muscles and two mapped joint goals. Those quantities are simulated, and the macaque velocity mapping does not prove neural or skill transfer into a human body.\par\nopagebreak[4]

\end{samepage}


\noindent\textit{Current registry lookup labels:} MyoArm; human MyoArm; MyoArm right\par

% Stable term ID: internal-musculoskeletal-model
\begin{samepage}
\subsection*{Internal musculoskeletal model}
\noindent The internal anatomy part of the body model: specified bones, joints and soft tissues, assembled from appropriate registered anatomical measurements.\par
\end{samepage}
\noindent Proposed complementary X-ray/CT and MRI inform different structures. Anatomy does not directly measure activation, force, material properties or neural function. Wētā FX provides visual/simulation inspiration, not a verified identical acquisition pipeline.\par

% Stable term ID: interoception
\begin{samepage}
\subsection*{Interoception}
\noindent Information about the body’s internal condition, such as signals related to heartbeat or breathing.\par
\end{samepage}
\begin{samepage}
\noindent\textit{Scope:} Sensing the body’s internal condition differs from sensing a limb’s position and movement. Internal sensing remains a research objective; this prototype does not provide a fully measured body-wide internal signal.\par\nopagebreak[4]

\end{samepage}


\noindent\textit{Current registry lookup labels:} Sensing the body from inside; interoception; interoceptive signal; internal sensations\par

% Stable term ID: interoceptome
% Canonical status: "user_requested_proposed_data_organization"
\subsection*{Interoceptome}
\noindent A heartbeat signal and a breathing signal describe different parts of the body’s internal condition. An Interoceptome would organize the selected observations and their relationships, retaining how, where and when each was recorded.\par
\noindent A proposed organization of specified observations of internal-body sensing and the relationships being studied among them.\par
\noindent This is a data-organization proposal within Interoceptomics. Its coverage must name the organism, body systems, measurements, timing and task. Measured signals, calculated quantities and model estimates remain distinguishable. Touch and limb-position sensing are not relabelled as interoception, and the suffix does not establish a complete internal-body map or an individual biological type.\par

% Stable term ID: interoceptomics
\begin{samepage}
\subsection*{Interoceptomics}
\noindent How does the body sense a heartbeat or its breathing, and how are the nerve connections involved arranged? Interoceptomics brings those questions together in a proposed study, relating observations of internal-body sensing to maps of the connections carrying or processing that information.\par
\end{samepage}
\noindent A proposed field combining interoception and connectomics to study internal-body sensing and its relationships with neural connectivity across scales.\par
\noindent The programme uses Interoceptomics as an umbrella for this proposed research. Giving the work a name does not establish a validated comprehensive map or a recognized external standard.\par
\noindent\textit{Scope:} macro: Whole-body, organ and network relationships relevant to internal-body sensing; micro: Local signals, cells, circuits and interfaces relevant to those relationships.\par
\begin{samepage}
\noindent\textit{Organizational level:} whole system; network; organ; local circuit; cell; signal/feature\par\nopagebreak[4]
\nopagebreak[4]
\end{samepage}


\noindent\textit{Current registry lookup labels:} interoceptomics\par

% Stable term ID: keeping-each-source-separate
\begin{samepage}
\subsection*{Keeping each source separate}
\noindent Keeping track of which recording, channel, time and physical setting a value belongs to.\par
\end{samepage}
\begin{samepage}
\noindent A channel numbered 12 in one day is not automatically the same cell as channel 12 on another day. A software label also cannot prove where supplied data physically came from.\par\nopagebreak[4]

\end{samepage}


\noindent\textit{Current registry lookup labels:} Source conditioning; source-conditioned; source-specific; recording-specific\par

% Stable term ID: keeping-each-trial-together
\begin{samepage}
\subsection*{Keeping each trial together}
\noindent Assigning an entire trial to one data portion instead of dividing it among portions.\par
\end{samepage}
\begin{samepage}
\noindent This prevents pieces of the same trial from being shared between the relevant splits. It does not remove prior inspection or other source limitations.\par\nopagebreak[4]

\end{samepage}


\noindent\textit{Current registry lookup labels:} Whole-trial separation; whole-trial-separated; successful trial; trial partition\par

% Stable term ID: keeping-up-now-and-computing-for-later
\begin{samepage}
\subsection*{Keeping up now and computing for later}
\noindent Real-time computation keeps pace with incoming data; offline computation can take longer and be played back later.\par
\end{samepage}
\begin{samepage}
\noindent Native planning takes about 4.52 seconds for a 20 ms simulated step and is offline. Smooth cached playback does not prove live planning or biological control.\par\nopagebreak[4]

\end{samepage}


\noindent\textit{Current registry lookup labels:} Real-time and offline execution; real-time; offline; paced playback; cached reference\par

% Stable term ID: kind-of-measurement
\begin{samepage}
\subsection*{Kind of measurement}
\noindent What a method records, such as electrical spikes, scalp voltage, muscle electrical activity or a brain-scan signal.\par
\end{samepage}
\begin{samepage}
\noindent These signals measure different things. They cannot be treated as interchangeable numbers merely because all of them relate to the nervous system.\par\nopagebreak[4]

\end{samepage}


\noindent\textit{Current registry lookup labels:} Signal modality; modality; observation process\par

% Stable term ID: large-movement-model
\subsection*{Large Movement Model}
\noindent The proposed movement-focused learned model that brings neural structure and functional observations together with internal musculoskeletal anatomy and independently measured external movement in a declared task.\par
\noindent The shorthand “NeuroTranslate with MoCap MyoSim = Large Movement Model” names the proposed combination of NeuroTranslate, MoCap ReComp and MyoSim Replay: a learned mapping, measured movement and virtual mechanics. The equals sign is architectural shorthand, not a mathematical identity or evidence of a trained model. EEG, EMG, measured movement, virtual mechanics and the trainee retain distinct roles.\par
\noindent This model-family name is analogous in ambition to large language model. It does not replace an accurate architecture description or make every deep neural network a movement model. Scans and capture alone do not train a model. Parameter count, scale, generalization and nervous-system fidelity need evidence. System 3 thinking remains the human–AI partnership framing.\par

\noindent\textit{Current registry lookup labels:} large movement models; LMM; Large movement model\par

% Stable term ID: lasting-change-in-the-nervous-system
\begin{samepage}
\subsection*{Lasting change in the nervous system}
\noindent Changes in the nervous system that can contribute to learning or adaptation.\par
\end{samepage}
\begin{samepage}
\noindent Updating software weights or controller gains is a different event. The practice interface in the computer model does not measure synaptic change in a living organism.\par\nopagebreak[4]

\end{samepage}


\noindent\textit{Current registry lookup labels:} Biological plasticity; synaptic plasticity; biological learning\par

% Stable term ID: learning-before-the-final-task
\begin{samepage}
\subsection*{Learning before the final task}
\noindent Training first on another objective, such as predicting a later neural count, before training for a paired movement target.\par
\end{samepage}
\begin{samepage}
\noindent The paper calls this pretraining. Earlier source-task labels remain part of the information the model received, even when only new-setting labels are counted in a small budget.\par\nopagebreak[4]

\end{samepage}


\noindent\textit{Current registry lookup labels:} Pretraining; neural next-bin pretraining\par

% Stable term ID: mapping-movement-onto-another-body-model
\begin{samepage}
\subsection*{Mapping movement onto another body model}
\noindent Turn a specified movement quantity, such as fingertip velocity components, into movement goals for a particular computer body. The mapping depends on that body and the chosen task.\par
\end{samepage}
\noindent Mapping specified movement quantities into goals for a particular body model, with the source quantities, target coordinates and transformation declared. This does not by itself reconstruct full posture or demonstrate biological transfer.\par
\begin{samepage}
\noindent The retained source-specific adapter takes two macaque fingertip velocity components as its source and produces human-model shoulder and elbow goals. Their signs carry direction; scalar speed would discard it. These are goals for that model and task, so the mapping establishes neither cross-species skill transfer nor verified biological function.\par\nopagebreak[4]

\end{samepage}


\noindent\textit{Current registry lookup labels:} Retargeting; body adapter; organism-specific body adapter; joint goal\par

% Stable term ID: meme-ification
\begin{samepage}
\subsection*{Meme-jacking}
\noindent Borrow the recognition; earn the explanation. A familiar image should help the audience see the scientific point before the script develops it.\par
\end{samepage}
\noindent Use a recognizable meme or cultural reference to make a scientific point immediately legible, drawing on familiarity already earned by others.\par
\begin{samepage}
\noindent Memification hijacking is the user’s expansion. References may come from any online culture, not only films. Apply one purposeful reference per teaching concept without crowding formal methods or labels. Recognition does not grant reuse rights, endorsement, proven comprehension or scientific validity. Preserve the original Meme-ification history.\par\nopagebreak[4]

\end{samepage}
\noindent Meme-ification/memeification are spelling forms; mimification is retained only as the corrected historical transcription.\par


\noindent\textit{Current registry lookup labels:} Meme-ification; meme-ification; memeification; mimification; meme-jacking; meme jacking; memification hijacking\par

% Stable term ID: mind-meld
\subsection*{Mind-Meld}
\noindent What information could one system give another, and would it help with the task? Mind-Meld names the proposed first Brain2Brain encounter, whose signal and effects would have to be measured.\par
\noindent A proposed first Brain2Brain study of a specified information link between neural systems during a declared task.\par
\noindent The roadmap places this at step 3.1 before later NeuroTransfer studies. The name does not describe fused minds, shared private experience, arbitrary thought reading or demonstrated skill transfer. Specify the source, recipient, recording and feedback interfaces, task, controls and independent outcome. Assistance during a link and learning retained without trainer information need separate tests, with any necessary compensatory bridge held constant. A reported natural connection is a motivating case, not evidence that this engineered link exists.\par

% Stable term ID: mind-merge
\subsection*{Mind-Merge}
\noindent An agent could guide practice by giving feedback about a stated task. Mind-Merge names the proposed study of that interaction: what help does the person receive, and what can they do later without it?\par
\noindent A proposed Agent2Brain study of task-specific feedback from a software agent to a person and its measurable effects.\par
\noindent The roadmap places this at step 4.1. Ordinary sensory feedback and direct neural delivery remain different routes with different interfaces and evidence about what is delivered. EEG and EMG record; neither supplies a neural writer. No merged identity, shared mind, retained skill or physical intervention follows from the name. Preserve the reverse observation path, matched practice and delayed/new-task retention tests, removing trainer information while keeping a necessary bridge constant. No implemented or qualified spinal-writing interface is demonstrated here.\par

% Stable term ID: mixed-source-animal-and-added-tissue
\begin{samepage}
\subsection*{Mixed-source animal and added tissue}
\noindent A chimera has biological contributions from different sources. A graft is added tissue; when it crosses species, it is a xenotransplant. An organoid is an organized tissue model.\par
\end{samepage}
\begin{samepage}
\noindent Survival, physical connection, activity and useful function are separate questions. The proposed comparisons do not establish a humanized whole brain or an implemented treatment.\par\nopagebreak[4]

\end{samepage}


\noindent\textit{Current registry lookup labels:} Chimera, graft and organoid; chimera; chimerism; graft; organoid; host–graft; xenotransplant; xenograft; XenoTransplantation\par

% Stable term ID: motion-capture-composite-relay
\subsection*{MoCap ReComp}
\noindent MoCap ReComp is Motion-Capture Relay-Composite. It proposes two dot maps (3D point clouds) of the same moving body, one from above and one from below, aligned in time and a common physical coordinate frame and joined while retaining each point’s source.\par
\noindent Bidirectional denotes those opposing measurements, not neural communication or causal direction. LiDAR ranging and infrared optical motion capture are proposed alternatives; marker triangulation and LiDAR point clouds are different measurements, so the actual modality and camera/ranging geometry must be qualified. The maze geometry calls for measurements perpendicular to the floor, with viewing paths from both sides through material intended to be clear. Transmission and material effects need testing; concentrated and distributed arrays remain alternatives. Hardware, coverage, timing, calibration, spatial alignment, occlusion, motion error and latency require validation against real references. The joined map is a derived record, not a direct observation of hidden joints, tissue mechanics, muscle commands or neural function. Neither successful implementation, complete body reconstruction nor material transparency has been established by the name. The single vertical camera illustration remains schematic.\par
\noindent\textit{Geometry:} Declare the top/bottom sensor positions, measured object, coordinate frame, viewing/illumination paths and coverage. Preserve the requested perpendicular-to-maze-floor geometry; document deviations rather than inferring them from an illustration.\par
\noindent\textit{Modality:} Declare LiDAR ranging or infrared optical motion capture, including marker-based/markerless method, camera count, triangulation geometry, material path and actual calibration. No complete reconstruction or physiological writing follows from opposing views.\par
\noindent\textit{Timing:} Record source timestamps, clocks, acquisition cadence, synchronization error and any interpolation, offline processing or end-to-end latency. No real-time capability is assumed.\par
\noindent\textit{Calibration and registration:} Preserve calibration data and transformations placing both streams into a common physical frame; measure registration error and keep original observations and reconstruction steps traceable.\par
\noindent\textit{Limits:} Record overlap, occlusion, missing observations, material effects, uncertainty and failure cases. Opposing measurements alone do not prove complete 3D pose, forces or biological function.\par
\noindent\textit{Physical scale:} Declare measured distance ranges and units for an actual acquisition; no supported numeric range or physical unit assignment is presently available for this proposal.\par
\noindent\textit{Statistical aggregation:} Declare any chosen evaluation grouping/averaging rule separately from physical scale and computational granularity.\par
\noindent\textit{References and acceptance:} Use \hyperref[app:reference-comparison-protocol]{the Reference Comparison Protocol and unfilled record} to identify the real-reference series needed for each acquisition component and the assembled map. Its provenance, rights, organism-stage, physical-unit, modality, clock, calibration, occlusion, contact, independent held-out, error, uncertainty and acceptance fields remain unfilled and HOLD. It supplies no measured series or apparatus validation. The unchanged programme reference criteria and every stronger original requirement govern their respective claims; no acceptance threshold is introduced or relaxed.\par
\noindent\textit{Outputs:} Retain the top dot map, bottom dot map and joined map as separate traceable outputs. Their timestamps and common-frame transformations must be explicit; do not fill unobserved regions as though they were measured.\par
\noindent\textit{Scope:} proposed application: Complementary top/bottom motion measurements of a body moving through the maze; concentrated coordinated platforms and distributed arrays remain alternatives; measurement boundary: Observed distances or geometry remain distinct from registered/stitched composites and any model-estimated pose or musculoskeletal quantity.\par
\noindent\textit{Organizational level:} whole system; signal/feature\par
\begin{samepage}
\noindent\textit{Computational granularity:} top dot map (3D point cloud); bottom dot map (3D point cloud); time alignment and spatial registration; joined composite motion sequence\par\nopagebreak[4]

\end{samepage}
\noindent Historical spellings name this same proposed capture method. Relay here means passing captured measurements onward; the separate agent-coordination Relay protocol does not describe motion capture.\par


\noindent\textit{Current registry lookup labels:} motion capture composite relay; motion-capture composite relay; motion capture composite; bidirectional composite motion capturing using LiDAR; Motion capture composite relay; motion capture relay composite; motion-capture relay composite; bidirectional LiDAR; opposed-view infrared motion capture; bidirectional infrared; motion-capture relay-composite; Motion-Capture Relay-Composite\par

% Stable term ID: model-that-chooses-a-label
\begin{samepage}
\subsection*{Model that chooses a label}
\noindent A statistical model that assigns scores to categories, such as the ten attempted-finger labels.\par
\end{samepage}
\begin{samepage}
\noindent The selected human model uses logistic classification, but its scores do not establish validated confidence in biological truth. Its reported accuracy comes from development data that have already been seen.\par\nopagebreak[4]

\end{samepage}


\noindent\textit{Current registry lookup labels:} Logistic classifier; logistic regression; classifier; classification; logistic\par

% Stable term ID: model-weights-kept-fixed
\begin{samepage}
\subsection*{Model weights kept fixed}
\noindent Keeping the fitted numbers in a model unchanged during a defined comparison.\par
\end{samepage}
\begin{samepage}
\noindent The tested transfer trains a connector while retaining the older main model. Fixed weights make the comparison clear; they do not prove biological sameness across days.\par\nopagebreak[4]

\end{samepage}


\noindent\textit{Current registry lookup labels:} Frozen parameters; frozen trunk; frozen readout; frozen checkpoint\par

% Stable term ID: modeled-route-from-one-muscle-attachment-to-another
\begin{samepage}
\subsection*{Modeled route from one muscle attachment to another}
\noindent The route and length assigned to a muscle and tendon as the model’s joints move.\par
\end{samepage}
\begin{samepage}
\noindent The native arm retains 39 named routes. Their modeled lengths differ from direct measurements of living tissue, muscle commands and force.\par\nopagebreak[4]

\end{samepage}


\noindent\textit{Current registry lookup labels:} Muscle-tendon path; musculotendon length; modeled muscle length; geometric path\par

% Stable term ID: models-that-convert-inputs-or-predict-outputs
\begin{samepage}
\subsection*{Models that convert inputs or predict outputs}
\noindent One model can turn an input into a representation; another can use available inputs to predict a named result.\par
\end{samepage}
\begin{samepage}
\noindent The paper calls these encoders and decoders. Predicting counts from mechanics or mechanics from counts does not by itself show the biological cause of either.\par\nopagebreak[4]

\end{samepage}


\noindent\textit{Current registry lookup labels:} Encoder and decoder; encoding; decoding; readout; decode; encoder; decoder\par

% Stable term ID: movement-brain-area
\begin{samepage}
\subsection*{Movement brain area}
\noindent A brain area involved in movement. The paper uses the names motor cortex, MC and M1 for relevant source recordings.\par
\end{samepage}
\begin{samepage}
\noindent The signals help predict recorded movement-related targets. Knowing the brain area does not identify matching cells across recordings or reveal a whole-body command pathway.\par\nopagebreak[4]

\end{samepage}


\noindent\textit{Current registry lookup labels:} Motor cortex; M1; MC; primary motor cortex\par

% Stable term ID: muscle-electrical-recording
\begin{samepage}
\subsection*{Muscle electrical recording}
\noindent Electrical activity recorded from muscles; the method is called EMG.\par
\end{samepage}
\begin{samepage}
\noindent Muscle electrical activity differs from force and from the computer model’s command and activity states. The trained targets here are rescaled signal summaries.\par\nopagebreak[4]

\end{samepage}


\noindent\textit{Current registry lookup labels:} Electromyography; EMG\par

% Stable term ID: myoception
\subsection*{Myoception}
\noindent Sensing muscle-driven movement and the feedback that comes with it.\par
\noindent The programme’s term for the sensory side of musculoskeletal contraction and movement, including proprioception, tactile feedback and other task-relevant sensations linked to action.\par
\noindent Myoception names a proposed programme scope; it establishes neither complete recording coverage nor a new external scientific standard. Muscle electrical activity, proprioceptive or tactile measurements, neural recordings and simulated feedback remain distinct observations. EMG does not directly record sensory afferents or the full sensory consequences of contraction. Declare each source’s organism, anatomy, modality, task, units, timing, uncertainty and measured, derived or simulated status. The wider neural-data ambition includes brain, spinal and peripheral activity and relevant visceral signals; the current focus is the motor–sensory system. No body-wide coverage, unique hidden neural state or subjective experience is established by this term.\par

\noindent\textit{Current registry lookup labels:} myoception; M-Y-O-ception\par

% Stable term ID: musculoskeletal-simulation-physics-engine
\subsection*{MyoSim Replay}
\noindent MyoSim Replay is the proposed MusculoSkeletal Simulation Physics-Engine. It models specified muscles, bones, joints and their mechanical interactions using declared anatomy, parameters, forces, contact rules and control inputs. Predictions of movement and force would be tested against independent movement and other applicable physical measurements. The external MyoSim body-model collection is a separate resource; simulation replay is not a measured biological replay.\par
\noindent Within the five-component proposal, the motion-capture relay-composite would provide traceable measured movement to compare with those predictions. That comparison needs matching timing, coordinates, units and calibration, and tests that preserve anatomical and physical constraints, uncertainty, held-out cases, baseline controls, a series of real references and independent acceptance. Movement geometry alone does not measure muscle activation, internal forces, hidden neural state or subjective experience; force, activation and causal claims need additional independent references. Visual effects (VFX) rendering shows the model but cannot validate its mechanics. This remains a proposed simulation component, not a newly demonstrated implementation or external scientific field. Approximate simulation and biological virtualization verified within declared error bounds remain distinct; this component establishes neither universal precision nor whole-body completeness.\par
\noindent\textit{Scope:} programme role: Virtual-body mechanics component; initially constrained hand and foot tasks; inputs and outputs: Declared anatomical parameters, muscle/control inputs and physical context produce simulated movement and mechanical quantities. Independently acquired motion and other necessary observables remain references, not simulator ground truth.\par
\noindent\textit{Organizational level:} whole system; organ\par
\begin{samepage}
\noindent\textit{Computational granularity:} musculoskeletal model; mechanical solver; measurement-to-prediction comparison\par\nopagebreak[4]

\end{samepage}


\noindent\textit{Current registry lookup labels:} musculoskeletal simulation physics engine; musculoskeletal-simulation physics engine; musculoskeletal simulation; musculoskeletal physics model; musculoskeletal-simulation physics-engine; MusculoSkeletal Simulation Physics-Engine\par

% Stable term ID: names-used-for-source-subjects
\begin{samepage}
\subsection*{Names used for source subjects}
\noindent Indy, Han and Monkey L label macaque sources; N1 labels the human finger source.\par
\end{samepage}
\begin{samepage}
\noindent These labels keep different subjects separate. Similar brain-region or task names do not make them one synchronized recording.\par\nopagebreak[4]

\end{samepage}


\noindent\textit{Current registry lookup labels:} Source subject labels; Indy; Han; Monkey L; N1\par

% Stable term ID: nerves-bringing-body-information-in
\begin{samepage}
\subsection*{Nerves bringing body information in}
\noindent Afferent paths carry body-related information toward the brain and spinal cord, together called the central nervous system.\par
\end{samepage}
\begin{samepage}
\noindent For example, a limb-related sensory signal can travel inward. The source-specific cortical predictions do not measure this whole path or make missing nerve activity observed.\par\nopagebreak[4]

\end{samepage}


\noindent\textit{Current registry lookup labels:} afferent; afferent pathway; body-to-CNS; CNS; central nervous system\par

% Stable term ID: nerves-carrying-commands-out
\begin{samepage}
\subsection*{Nerves carrying commands out}
\noindent Efferent paths carry outputs from the brain and spinal cord toward the body, including muscles and organs.\par
\end{samepage}
\begin{samepage}
\noindent A prediction of movement is not a recording of that outward pathway. The prototype does not supply a synchronized voluntary sequence through spinal cord, nerves and muscles.\par\nopagebreak[4]

\end{samepage}


\noindent\textit{Current registry lookup labels:} efferent; efferent pathway; CNS-to-body\par

% Stable term ID: neural-decoding
\begin{samepage}
\subsection*{Neural decoding}
\noindent Estimating declared information such as intended action from recorded neural observations using a specified model.\par
\end{samepage}
\noindent Decode accuracy, channels, source identity, task and held-out domain must be demonstrated. Brain-only control requires exclusion of covert muscle, eye and motion signals; decoding does not write signals into the brain.\par



\noindent\textit{Source:} Restoring cortical control of functional movement in a human with quadriplegia\par
\noindent\textit{Source identifier:} bouton2016\par
\noindent\textit{DOI:} 10.1038/nature17435\par
\noindent\url{https://www.nature.com/articles/nature17435}\par
\noindent Research precedent linking implanted cortical recording to a custom neuromuscular stimulation sleeve; not an EEG/OpenBCI or elite-skill result\par
\noindent\textit{Source:} A brain-computer interface that evokes tactile sensations improves robotic arm control\par
\noindent\textit{Source identifier:} flesher2021\par
\noindent\textit{DOI:} 10.1126/science.abd0380\par
\noindent\url{https://pubmed.ncbi.nlm.nih.gov/34016775/}\par
\noindent Research precedent for motor decoding plus sensory feedback through intracortical microstimulation; invasive and not a consumer EEG capability\par

% Stable term ID: neural-encoding
\begin{samepage}
\subsection*{Neural encoding}
\noindent Mapping declared information to a representation in neural activity; in a sensory-interface design, constructing candidate stimulation patterns for specified feedback.\par
\end{samepage}
\noindent A computational encoder is distinct from a qualified stimulation device and evidence that a participant receives the intended sensation. The term does not mean arbitrary thought writing or skill upload.\par



\noindent\textit{Source:} A brain-computer interface that evokes tactile sensations improves robotic arm control\par
\noindent\textit{Source identifier:} flesher2021\par
\noindent\textit{DOI:} 10.1126/science.abd0380\par
\noindent\url{https://pubmed.ncbi.nlm.nih.gov/34016775/}\par
\noindent Research precedent for motor decoding plus sensory feedback through intracortical microstimulation; invasive and not a consumer EEG capability\par

% Stable term ID: playing-back-a-recording
\begin{samepage}
\subsection*{Neural Replay}
\noindent Revisiting a stored sequence of neural samples with its stated timing and channel order.\par
\end{samepage}
\begin{samepage}
\noindent This playback is not a live recording or instructions for stimulating a nervous system. A displayed arm movement does not show what the source animal’s muscles actually did.\par\nopagebreak[4]

\end{samepage}


\noindent\textit{Current registry lookup labels:} Neural replay; recorded replay; archived replay; cached playback; Playing back a recording\par

% Stable term ID: neuroagent-tutor
\subsection*{NeuroAgent Trainer}
\noindent A software agent would guide a trainee during a stated task.\par
\noindent A proposed software-agent trainer role in a declared Agent2Brain NeuroTransfer study.\par
\noindent Identify the actual software model and feedback route. Computed activations/events do not make a biological trainee digital. Illustrative connections do not establish direct neural writing, successful transfer or retained learning. NeuroAgent Tutor is a historical label; NeuroAgent Trainer is the current name. It describes a proposed software-agent trainer role, without selecting the agent or qualifying its feedback interface.\par
\noindent\textit{Organizational level:} whole system; signal/feature\par

\noindent\textit{Current registry lookup labels:} NeuroAgent Tutor\par

% Stable term ID: neuroarcher
\begin{samepage}
\subsection*{NeuroArcher}
\noindent A proposed archery game that studies coordinated use of both hands and arms through physical motion control, brain-signal control of a digital body, and a separately admitted later assisted-body experiment.\par
\end{samepage}
\noindent The near-term prop is a toy bow with small suction-cup tips and an over-arm arrow pouch. Target/creature choice is unresolved. Simulated performance does not establish muscle stimulation or retained biological archery skill.\par





\noindent\textit{Current registry lookup labels:} Neuro Archer\par

% Stable term ID: neuroboxer
\begin{samepage}
\subsection*{NeuroBoxer}
\noindent A proposed boxing game that compares arm-motion control, brain-signal control of a digital body, and a separately admitted later assisted-body experiment.\par
\end{samepage}
\noindent Prime-era Mike Tyson is the requested virtual opponent and meme reference. A game, recorded prediction, stimulation-assisted motion and retained biological skill learning need different evidence. No actual opponent, head contact, finished game or stimulation capability is claimed.\par





\noindent\textit{Current registry lookup labels:} Neuro Boxer\par

% Stable term ID: neuroconnectome
\begin{samepage}
\subsection*{NeuroConnectome}
\noindent A map of specified neural connections used to constrain a model of the nervous system.\par
\end{samepage}
\begin{samepage}
\noindent Connectome is the established term. The programme's NeuroConnectome label does not establish complete coverage or identify functional dynamics, state, plasticity, sensory context or unique thought from wiring alone. Declare species, anatomical coverage, resolution and provenance.\par\nopagebreak[4]

\end{samepage}


\noindent\textit{Current registry lookup labels:} Neuroconnectome\par

% Stable term ID: neurohandler
% Canonical status: "user_requested_proposed_programme_role"
\subsection*{NeuroHandler}
\noindent How well can a hand carry out the chosen task, and what information helps it do so? NeuroHandler is the proposed manual-dexterity experiment, with the hand’s measured action providing the result to explain.\par
\noindent A proposed manual-dexterity experiment studying specified neural/body observations, task control and independently measured outcomes.\par
\noindent The earlier proposed manual tasks remain: rapid hand rubbing, the requested knuckle gesture and the interlaced-hand stretch. NeuroHandler names the manual-dexterity experiment; it does not establish a finished apparatus or successful model. Define the action, available inputs, physical or simulated setting, comparison and outcome before evaluation. Keep motion control, neural prediction, assistance and retained learning separate; any physical output requires separate qualification and admission. No dexterity gain, clinical benefit or intervention procedure is established.\par

% Stable term ID: neurojuggler
% Canonical status: "user_requested_proposed_programme_role"
\subsection*{NeuroJuggler}
\noindent Moving each limb successfully does not settle whether the limbs work together. NeuroJuggler is the proposed coordination experiment: define the joint task, then measure the relationship between the limb actions and the result.\par
\noindent A proposed multi-limb coordination experiment studying how specified limb actions are coordinated in a declared task.\par
\noindent The earlier coordination proposal remains a hot-potato transfer between the hands, followed preferably by two-ball juggling. Those are proposed tasks; the final apparatus and qualified measurement setup remain unestablished. The label supplies no successful performance or implemented neural interface. Declare the task, limbs, inputs, timing, body model or physical setting and independent outcomes. Separate coordination from isolated-limb performance, simulated success from biological function, and current assistance from retained learning. Physical output, safety, rights and held-out tests keep their own admission requirements.\par

% Stable term ID: neuromuscular-electrical-stimulation
\begin{samepage}
\subsection*{Neuromuscular electrical stimulation}
\noindent Electrical stimulation intended to elicit muscle contraction through the relevant neuromuscular pathways.\par
\end{samepage}
\noindent NMES is a writing/actuation intervention, whereas EMG records muscle activity. EMG electrodes or an OpenBCI recording board do not by themselves provide qualified stimulation. Neither contraction nor imposed movement demonstrates retained skill.\par



\noindent\textit{Source:} Restoring cortical control of functional movement in a human with quadriplegia\par
\noindent\textit{Source identifier:} bouton2016\par
\noindent\textit{DOI:} 10.1038/nature17435\par
\noindent\url{https://www.nature.com/articles/nature17435}\par
\noindent Research precedent linking implanted cortical recording to a custom neuromuscular stimulation sleeve; not an EEG/OpenBCI or elite-skill result\par


\noindent\textit{Current registry lookup labels:} NMES; Electrical muscle stimulation; EMS\par

% Stable term ID: neuroreplay
\subsection*{NeuroReplay}
\noindent What changes when a trainee receives replay-based guidance? NeuroReplay names that proposed study: compare performance while guidance is available, then test what remains when the guidance ends.\par
\noindent A proposed Stage 0 research direction for studying replay-guided sensorimotor feedback and its effects on current performance and later learning.\par
\noindent This is the new roadmap label for steps 0.2–0.5, with a derived working scope; it does not rename Neural Replay, the existing stored-sample playback entry playing-back-a-recording. A stored sequence can supply a comparison, but playback alone establishes neither delivery to a nervous system nor biological learning. Fix the task, signal, feedback route, timing and equal-practice comparisons; distinguish matched, stale or unmatched replay, no-link and inactive or sham controls. Test delayed retention and new tasks after removing trainer information while keeping any necessary compensatory bridge constant. The questioned antechinus age ranges are not established developmental or intervention windows. No physical writer, procedure or present result is supplied by this label.\par

% Stable term ID: neuroscience-recording-files
\begin{samepage}
\subsection*{Neuroscience recording files}
\noindent NWB stores neuroscience data and descriptions. In the paired source files here, SET holds recording descriptions and FDT holds binary samples.\par
\end{samepage}
\begin{samepage}
\noindent A stale unit in the NWB muscle source must not override verified processing. SET headers match source hashes; FDT previews have checked local byte ranges rather than full verification of each upstream file.\par\nopagebreak[4]

\end{samepage}


\noindent\textit{Current registry lookup labels:} Neurodata Without Borders file; NWB; Recording header and sample file; SET and FDT files; SET; FDT; EEGLAB header; byte-range preview\par

% Stable term ID: neurotransfer
\begin{samepage}
\subsection*{NeuroTransfer}
\noindent Specify what would move between which systems, then test what the recipient can do with and without assistance.\par
\end{samepage}
\noindent The proposed transfer of specified task-relevant information or a mapping between declared systems, with assistance and retained learning evaluated separately.\par
\begin{samepage}
\noindent The label does not establish a successful biological transfer or retained skill. A weight copy, recorded replay, guided movement and independently retained learning answer different questions. Each proposed interface needs its own source, target, task, controls, rights and qualified outcome.\par\nopagebreak[4]

\end{samepage}


\noindent\textit{Current registry lookup labels:} neurotransfer\par

% Stable term ID: neurotranslate
\begin{samepage}
\subsection*{NeuroTranslate}
\noindent Brain2Body Neural Activity is the long-form component label paired with NeuroTranslate in this proposed architecture. Recorded activity is an observation. NeuroTranslate names the mapping operation that could relate specified observations to movement, or reconstruct specified recorded features in the reverse direction. The observation, learned relationship and predicted value remain different things.\par
\noindent State the observations, the target and the task, then test how well the mapping works.\par
\end{samepage}
\noindent The proposed operation of mapping specified neural and body observations between declared inputs and targets in a stated task.\par
\begin{samepage}
\noindent A model could predict movement from specified observations, or reconstruct specified recorded features in the reverse direction. Each direction needs its own test in the stated task. Neither establishes a unique hidden state, sensory delivery or biological skill transfer. NeuroTranslate names this proposed mapping operation; NeuroTranslation names the field and NeuroTranslator the implemented software.\par\nopagebreak[4]

\end{samepage}


\noindent\textit{Current registry lookup labels:} neurotranslate; Brain2Body Neural Activity\par

% Stable term ID: neurotranslation
\begin{samepage}
\subsection*{NeuroTranslation}
\noindent How could a particular neural recording be related to a defined movement task? NeuroTranslation is the proposed study of mapping nerve signals and the patterns that describe them between clearly defined situations like these. Every mapping must state where its inputs come from, what its outputs mean, and how it is tested.\par
\end{samepage}
\noindent A proposed field for translating or mapping neural signals and representations between explicitly defined source and target domains, with stated context, scale, tests and limitations.\par
\noindent Each connection must be assessed against its named source and stated test. The project has not demonstrated a universal brain language or an uploaded skill.\par
\noindent\textit{Scope:} macro: Relationships between neural systems, body systems and defined tasks or organisms; micro: Mappings involving identified local signals, channels, features and circuit-level representations.\par
\begin{samepage}
\noindent\textit{Organizational level:} whole system; network; organ; local circuit; signal/feature\par\nopagebreak[4]
\nopagebreak[4]
\end{samepage}


\noindent\textit{Current registry lookup labels:} neurotranslation; neural translation; Connecting neural recordings to defined tasks; Neurotranslation\par

% Stable term ID: source-conditioned-neurotranslation-architecture
\begin{samepage}
\subsection*{NeuroTranslation Network}
\noindent A plan for connecting more neural and body components while keeping the origin and meaning of each value clear.\par
\end{samepage}
\begin{samepage}
\noindent The proposed network would connect neural and body components while carrying the origin and meaning of each value through the mapping. Each link could then be examined in its own recording context. Body-wide translation and biological teaching remain unproved.\par\nopagebreak[4]

\end{samepage}


\noindent\textit{Current registry lookup labels:} Proposed larger design; Source-Conditioned Neurotranslation Architecture; SCNA; Neurotranslation Network\par

% Stable term ID: neurotranslation-engine
\begin{samepage}
\subsection*{NeuroTranslator}
\noindent The software that joins particular recordings, prediction models, playback and computer-body checks.\par
\end{samepage}
\begin{samepage}
\noindent Predictions and playback can be checked against the particular recordings they use. That is the scope of the implemented NeuroTranslator; neither a universal brain language nor skill upload has been demonstrated.\par\nopagebreak[4]

\end{samepage}


\noindent\textit{Current registry lookup labels:} Working translation software; Neurotranslation Engine; NTE; Neurotranslator\par

% Stable term ID: neurotwin-tutor
\subsection*{NeuroTwin Trainer}
\noindent A biological trainer would guide a trainee during a stated task.\par
\noindent A proposed biological trainer role in a declared Brain2Brain NeuroTransfer study.\par
\noindent The requested name alone establishes neither genetic twins, a validated digital twin, communication nor retained learning. Biological electrical activity is electrochemical and can include graded voltages and spikes; identify any actual measurement, otherwise keep the activity illustrative. NeuroTwin Tutor is a historical label; NeuroTwin Trainer is the current name. These names belong to a proposed programme view. They establish neither an implemented interface nor a demonstrated learning result.\par
\noindent\textit{Organizational level:} whole system; signal/feature\par

\noindent\textit{Current registry lookup labels:} NeuroTwin Tutor\par

% Stable term ID: numerical-arrays-saved-arrays-and-prediction
\begin{samepage}
\subsection*{Numerical arrays, saved arrays and prediction}
\noindent NumPy provides tools for numerical arrays. NPZ saves arrays; inference means running a trained model to produce an output.\par
\end{samepage}
\begin{samepage}
\noindent Portable execution still needs the correct processing and source identity. Matching saved-model outputs does not establish success on new organisms or biological skills.\par\nopagebreak[4]

\end{samepage}


\noindent\textit{Current registry lookup labels:} NumPy inference and NPZ archive; NumPy; NPZ; portable inference; inference\par

% Stable term ID: numerical-checks-on-allowed-movement
\begin{samepage}
\subsection*{Numerical checks on allowed movement}
\noindent Checks that the computer model stays defined and within its stated joint limits.\par
\end{samepage}
\begin{samepage}
\noindent Joint-limit penetration means going past an angular bound, not entering tissue. The 0.001 rad tolerance is a numerical rule, not a certificate of correct neural commands.\par\nopagebreak[4]

\end{samepage}


\noindent\textit{Current registry lookup labels:} Mechanical feasibility and anatomical tolerance; anatomical limit; limit penetration; finite state; joint range; numerical tolerance\par

% Stable term ID: numerical-software-and-its-data-file
\begin{samepage}
\subsection*{Numerical software and its data file}
\noindent MATLAB is numerical-computing software, and MAT is a format for its structured data.\par
\end{samepage}
\begin{samepage}
\noindent The second Indy recording uses the original MAT source. The format does not prove matching cells, past-only source processing or complete observation.\par\nopagebreak[4]

\end{samepage}


\noindent\textit{Current registry lookup labels:} MATLAB and MAT file; MATLAB; MAT\par

% Stable term ID: observe-it-connect-it-keep-it-useful
\begin{samepage}
\subsection*{Observe it, connect it, keep it useful}
\noindent Observation burden, prediction on excluded cases and lasting task benefit need separate measurements. A clear picture does not answer those experimental questions.\par\nopagebreak[4]

\end{samepage}


\noindent\textit{Current registry lookup labels:} Access it, Ground it, Keep it; access; grounding; stability; signal stability\par

% Stable term ID: one-multiplier-for-modeled-lengths
\begin{samepage}
\subsection*{One multiplier for modeled lengths}
\noindent A single number chosen to reduce differences between the model’s muscle lengths and the source-derived length references.\par
\end{samepage}
\begin{samepage}
\noindent The number is fitted on training poses and checked on distinct observed poses. It is not complete body calibration; no movement-control experiment is run on the scaled candidate.\par\nopagebreak[4]

\end{samepage}


\noindent\textit{Current registry lookup labels:} Scalar length fit; global length factor; uniform factor; a*; scalar bridge; recording length-fit candidate; SDK-scaled candidate\par

% Stable term ID: one-recording-period
\begin{samepage}
\subsection*{One recording period}
\noindent One experiment or recording period with its own subject, channels, timing and task.\par
\end{samepage}
\begin{samepage}
\noindent Two recording days, or recordings from different animals, do not form one observed nervous system simply because the files are loaded together.\par\nopagebreak[4]

\end{samepage}


\noindent\textit{Current registry lookup labels:} Recording and session; source recording; session; recording family\par

% Stable term ID: open-body-model-collection
\begin{samepage}
\subsection*{Open body-model collection}
\noindent MyoSim supplies computer-body models with muscles and bones, including models used in this project.\par
\end{samepage}
\begin{samepage}
\noindent The collection is earlier work, not a new project invention. Availability of an arm model does not establish matching macaque and human muscles or neural commands.\par\nopagebreak[4]

\end{samepage}


\noindent\textit{Current registry lookup labels:} MyoSim, MyoSuite and MyoArm; MyoSim\par

% Stable term ID: open-muscle-control-task-tools
\begin{samepage}
\subsection*{Open muscle-control task tools}
\noindent MyoSuite provides tasks and tools for experiments with muscle-driven computer-body control.\par
\end{samepage}
\begin{samepage}
\noindent It is earlier software reused here. Success on a simulated task is not proof of biological learning or a recovered neural command.\par\nopagebreak[4]

\end{samepage}


\noindent\textit{Current registry lookup labels:} MyoSuite\par

% Stable term ID: other-published-models-of-neural-recordings
\begin{samepage}
\subsection*{Other published models of neural recordings}
\noindent NDT2 learns patterns across recordings made inside the cortex, the brain’s outer layer. POYO handles differing sets of recorded neural features by representing events and learning relationships between them.\cite{ndt2_2023,poyo2023}\par
\end{samepage}
\begin{samepage}
\noindent The paper cites these as earlier work. It does not claim to have implemented them or use their existence as proof of universal transfer or matching cells.\par\nopagebreak[4]

\end{samepage}


\noindent\textit{Current registry lookup labels:} NDT2; large-scale neural pretraining; Published model for varying recorded populations; POYO; neural-event tokenization; cross-attention bottleneck\par

% Stable term ID: passing-information-and-gaining-a-lasting-skill
\begin{samepage}
\subsection*{Passing information and gaining a lasting skill}
\noindent Receiving useful task information differs from acquiring a durable ability that works independently in new conditions.\par
\end{samepage}
\begin{samepage}
\noindent A short software label budget is not measured biological practice time. Sewing, cube solving and breakdancing remain future examples without matched skill recordings here.\par\nopagebreak[4]

\end{samepage}


\noindent\textit{Current registry lookup labels:} Task information transfer and skill acquisition; task transfer; skill transfer; skill upload; compressed acquisition; learning speed; learning; physical skill learning; skill acquisition\par

% Stable term ID: past-only-inputs-and-processing-that-uses-later-data
\begin{samepage}
\subsection*{Past-only inputs and processing that uses later data}
\noindent The local model uses completed past inputs. Some earlier source processing may still have used later samples.\par
\end{samepage}
\begin{samepage}
\noindent Past-only indexing cannot reverse earlier filtering or rescaling that used the whole recording. It does not certify the full measurement-to-output delay in a living system.\par\nopagebreak[4]

\end{samepage}


\noindent\textit{Current registry lookup labels:} Causal indexing and acausal preprocessing; causal; algorithmic support; acausal filtering; offline interpolation; upstream filtering delay\par

% Stable term ID: peripheral-nerve-electrical-recording
\begin{samepage}
\subsection*{Peripheral nerve electrical recording}
\noindent A method for recording electrical signals associated with peripheral nerves, called ENG.\par
\end{samepage}
\begin{samepage}
\noindent The human preview contains aggregate electrical observations. It does not provide separately identified nerve-cell spike trains or a complete voluntary command sequence.\par\nopagebreak[4]

\end{samepage}


\noindent\textit{Current registry lookup labels:} Electroneurography; ENG; surface ENG; peripheral recording\par

% Stable term ID: phenotype
% Canonical status: "established_scientific_reference"
\subsection*{Phenotype}
\noindent The trait is what must be observed: for example, a specified body characteristic or response in the setting being studied. Phenotype names these observable traits, rather than the genetic description alone.\par
\noindent The observable traits of an organism arising through genotype, environment and development.\par
\noindent State the actual trait, measurement method, developmental and environmental context. A phenotype is not a universal function that can be read uniquely from structure or sequence. Retain the distinction between a measured trait and a model’s estimate.\par

% Stable term ID: placing-body-parts-from-joint-positions
\begin{samepage}
\subsection*{Placing body parts from joint positions}
\noindent Calculating where body parts and muscle routes lie when joint positions are given.\par
\end{samepage}
\begin{samepage}
\noindent The formal name is forward kinematics, or FK. It places geometry; it does not calculate muscle-driven motion through time or recover muscle commands.\par\nopagebreak[4]

\end{samepage}


\noindent\textit{Current registry lookup labels:} Forward kinematics; FK; kinematic replay\par

% Stable term ID: planned-commands-corrections-and-constant-input
\begin{samepage}
\subsection*{Planned commands, corrections and constant input}
\noindent A planned command is feedforward; adjusting it from current model information is feedback; a constant muscle input is tonic control.\par
\end{samepage}
\begin{samepage}
\noindent The human correction can change the bounded command goal. Recorded movement stays a scoring reference, and the native correction candidates remain rejected against constant-input comparisons.\par\nopagebreak[4]

\end{samepage}


\noindent\textit{Current registry lookup labels:} Feedforward, feedback and tonic control; feedforward; feedback; corrected controller; tonic excitation; constant tonic\par

% Stable term ID: prediction-score-compared-with-a-mean-value-guess
\begin{samepage}
\subsection*{Prediction score compared with a mean-value guess}
\noindent R² compares squared prediction errors with a constant mean-value reference. Higher values usually indicate better prediction on the stated data.\par
\end{samepage}
\begin{samepage}
\noindent The score can be negative. Different recordings, target processing and ways of combining outputs prevent treating all R² values as directly comparable.\par\nopagebreak[4]

\end{samepage}


\noindent\textit{Current registry lookup labels:} Coefficient of determination; R²; R2; R\textasciicircum{}2; variance explained\par

% Stable term ID: processing-within-spinal-and-peripheral-pathways
\begin{samepage}
\subsection*{Processing within spinal and peripheral pathways}
\noindent Body-related information can be processed locally in spinal and peripheral structures rather than passing through one simple unchanged wire.\par
\end{samepage}
\begin{samepage}
\noindent A cortex-only prediction does not observe those local processes. The separate recordings and computer-body checks do not establish a complete measured brain-to-body causal chain.\par\nopagebreak[4]

\end{samepage}


\noindent\textit{Current registry lookup labels:} local spinal processing; local peripheral processing; spinal/peripheral processing\par

% Stable term ID: programmable-neurobiology
\begin{samepage}
\subsection*{Programmable neurobiology}
\noindent The proposed ability to specify a bounded nervous-system or sensorimotor outcome, select a qualified intervention, predict its effects, measure what happens and adapt within a validated closed loop.\par
\end{samepage}
\noindent\textit{Scope:} Brain, spinal, peripheral, muscle, mechanical and sensory feedback roles need separate evidence. Bidirectional prediction is different from safe sensory/action writing. An intermediary may improve a mapping without making it one-to-one. Arbitrary skill upload or universal disease/disability repair is not demonstrated.\par
\noindent\textit{Organizational level:} Declared system or interface\par
\noindent\textit{Computational granularity:} Declared process/interface/model\par
\noindent\textit{Statistical aggregation:} Validation domain stated independently of physical scale\par

% Stable term ID: programmable-synthetic-biology
\begin{samepage}
\subsection*{Programmable synthetic biology}
\noindent Engineering specified biological behaviours by combining design, context-dependent molecular prediction, experimental measurement and iterative correction.\par
\end{samepage}
\noindent DNA/regulation, RNA/cell state and protein/interactions are explanatory layers, not a sufficient literal three-model recipe. Metabolism, feedback, environment and measured function matter. AlphaGenome and AlphaFold do not jointly certify whole-cell programmability; AlphaTranscriptome is a desired capability placeholder here.\par
\noindent\textit{Organizational level:} Declared system or interface\par
\noindent\textit{Computational granularity:} Declared process/interface/model\par
\begin{samepage}
\noindent\textit{Statistical aggregation:} Validation domain stated independently of physical scale\par\nopagebreak[4]

\end{samepage}


\noindent\textit{Current registry lookup labels:} Programmable molecular biology\par

% Stable term ID: public-reaching-recordings
\begin{samepage}
\subsection*{Public reaching recordings}
\noindent Public files pairing cortical electrical events with reaching behavior. The sources use names such as RTT and MC\_RTT.\par
\end{samepage}
\begin{samepage}
\noindent The original model’s release and the separately acquired Indy day have different processing. Their dates and target definitions must remain explicit.\par\nopagebreak[4]

\end{samepage}


\noindent\textit{Current registry lookup labels:} Reaching source archive; RTT; MC\_RTT; MC\textbackslash{}\_RTT; O'Doherty reaching archive; Sabes reaching corpus\par

% Stable term ID: public-source-records-and-document-links
\begin{samepage}
\subsection*{Public source records and document links}
\noindent Repositories such as DANDI and OpenNeuro identify source versions. PDF is a document format, and a URL gives a resource location.\par
\end{samepage}
\begin{samepage}
\noindent Source licenses and credit must be retained. A local paper draft or public input dataset does not mean the project’s new code is published, licensed or submitted to arXiv or a journal.\par\nopagebreak[4]

\end{samepage}


\noindent\textit{Current registry lookup labels:} Public records, document files and links; DANDI; OpenNeuro; PDF; URL; arXiv; public deposition; license; attribution\par

% Stable term ID: published-models-for-brain-scans-and-scalp-recordings
\begin{samepage}
\subsection*{Published models for brain scans and scalp recordings}
\noindent BrainLM learns from fMRI brain-scan data. Neuro-GPT combines an EEG input model with a model of changes over time and is discussed for motor imagery.\cite{brainlm2023,neurogpt2024}\par
\noindent BrainLM version 1 is supported here at the abstract-level model and modality description. Neuro-GPT uses author version 4 for the cited EEG mechanism; exact equivalence to the IEEE final publication remains unverified. No performance result is transferred.\par
\end{samepage}
\begin{samepage}
\noindent These are earlier methods for different measurements. The project does not claim to have trained them, make scans and electrical spikes interchangeable, or upload skills.\par\nopagebreak[4]

\end{samepage}


\noindent\textit{Current registry lookup labels:} BrainLM; Published model of scalp signals; Neuro-GPT; NeuroGPT\par

% Stable term ID: published-neural-model-comparison-tasks
\begin{samepage}
\subsection*{Published neural-model comparison tasks}
\noindent The Neural Latents Benchmark, or NLB, provides shared tasks for comparing models of neural recordings.\cite{nlb2021}\par
\end{samepage}
\begin{samepage}
\noindent It informs the choice of comparisons. The paper’s own local scores are not official NLB benchmark scores.\par\nopagebreak[4]

\end{samepage}


\noindent\textit{Current registry lookup labels:} Neural Latents Benchmark; NLB\par

% Stable term ID: published-tests-of-changing-recordings
\begin{samepage}
\subsection*{Published tests of changing recordings}
\noindent FALCON is a published benchmark concerned with neural decoding when recordings change and recalibration is limited.\par
\end{samepage}
\begin{samepage}
\noindent The paper uses related source recordings and its own data splits. Its local muscle-signal scores are not official FALCON scores.\par\nopagebreak[4]

\end{samepage}


\noindent\textit{Current registry lookup labels:} FALCON benchmark; FALCON; FALCON M1; FALCON M1-A\par

% Stable term ID: ranges-that-need-statistical-assumptions
\begin{samepage}
\subsection*{Ranges that need statistical assumptions}
\noindent A confidence interval gives a range for a quantity being estimated. A conformal prediction range tries to contain a future actual result at a stated frequency when its required assumptions hold.\par
\end{samepage}
\begin{samepage}
\noindent The transfer figure shows one fixed run. The paper’s ranges are based on earlier prediction errors, with related errors across time; they do not guarantee how often future true values will fall inside.\par\nopagebreak[4]

\end{samepage}


\noindent\textit{Current registry lookup labels:} Confidence interval and conformal coverage; confidence interval; repeated-seed interval; Prediction range with tested coverage; conformal; conformal coverage; prospective-conformal; distribution-free coverage\par

% Stable term ID: recent-signal-model
\begin{samepage}
\subsection*{Recent-signal model}
\noindent A model component that keeps a learned memory of recent inputs as it predicts the next result.\par
\end{samepage}
\begin{samepage}
\noindent Its formal name is gated recurrent unit, or GRU. This established method uses completed input history; its memory does not establish a universal neural language.\par\nopagebreak[4]

\end{samepage}


\noindent\textit{Current registry lookup labels:} Gated recurrent unit; GRU; compact recurrent model; recurrent network\par

% Stable term ID: record-of-what-was-observed
\begin{samepage}
\subsection*{Record of what was observed}
\noindent A record showing which channels or intervals were actually measured and usable.\par
\end{samepage}
\begin{samepage}
\noindent Missing observations stay missing; they do not become zeros. The second Indy source does not supply separate observation intervals for every recorded unit.\par\nopagebreak[4]

\end{samepage}


\noindent\textit{Current registry lookup labels:} Observation mask; observation interval; invalid support; missing observation\par

% Stable term ID: recorded-neural-events
\begin{samepage}
\subsection*{Recorded neural events}
\noindent A spike is an electrical event identified in a recording. A spike train is the sequence of those events over time.\par
\end{samepage}
\begin{samepage}
\noindent A recorded event count is a whole number tied to one channel and time interval. A model’s average expected count can be a fraction and is a different kind of value.\par\nopagebreak[4]

\end{samepage}


\noindent\textit{Current registry lookup labels:} Spike and spike train; neural event; spike count; spike train\par

% Stable term ID: records-linking-files-checks-and-claims
\begin{samepage}
\subsection*{Records linking files, checks and claims}
\noindent A manifest lists specified files and identities; a receipt records checks; an evidence matrix links a claim to its supporting record.\par
\end{samepage}
\begin{samepage}
\noindent These records also keep failures and withdrawn results visible. A catalog listing does not show that every dataset was acquired or used to train a model.\par\nopagebreak[4]

\end{samepage}


\noindent\textit{Current registry lookup labels:} Manifest, receipt and evidence matrix; manifest; source manifest; audit receipt; processing receipt; evidence matrix; withdrawal record\par

% Stable term ID: relay
\begin{samepage}
\subsection*{Relay}
\noindent Pass focused work between agents, then check their results before accepting it.\par
\end{samepage}
\noindent A manual coordination protocol that passes bounded work and verified evidence between specialist sessions, with one integration owner and an independent acceptance reviewer.\par
\noindent A specialist can work on a bounded question while an independent reviewer checks its evidence before the integration owner accepts it. This is a project coordination workflow, distinct from an established external standard or native Ultra mode. Efficiency and quality claims require matched evidence from the tested task; no universal performance equivalence or measured saving follows from the protocol.\par
\begin{samepage}
\noindent\textit{Computational granularity:} finite delegated tasks and evidence artifacts\par\nopagebreak[4]

\end{samepage}


\noindent\textit{Current registry lookup labels:} relay; manual Relay protocol; Relay coordination\par

% Stable term ID: rescaled-muscle-signal
\begin{samepage}
\subsection*{Rescaled muscle signal}
\noindent A processed summary of muscle electrical activity put onto a relative scale.\par
\end{samepage}
\begin{samepage}
\noindent The 16 targets have no physical voltage unit. A stale millivolt label in the source does not make them measured voltage, force or muscle activation.\par\nopagebreak[4]

\end{samepage}


\noindent\textit{Current registry lookup labels:} Normalized EMG envelope; EMG envelope; quantile normalization; dimensionless EMG\par

% Stable term ID: rescaling-using-training-data-only
\begin{samepage}
\subsection*{Rescaling using training data only}
\noindent Learning a rescaling rule from the training portion, then applying the same rule to later values.\par
\end{samepage}
\begin{samepage}
\noindent This avoids using local test values to fit that rule. It cannot undo source processing that already used future samples or the whole recording.\par\nopagebreak[4]

\end{samepage}


\noindent\textit{Current registry lookup labels:} Training-only normalization; scaler; target statistics; normalization\par

% Stable term ID: response-after-an-electrical-stimulus
\begin{samepage}
\subsection*{Response after an electrical stimulus}
\noindent An electrical recording made while observing the response to a stimulus.\par
\end{samepage}
\begin{samepage}
\noindent The human preview is a stimulus-response experiment, not voluntary movement. The stimulus itself can also leave an unwanted signal in the recording; that is an artifact.\par\nopagebreak[4]

\end{samepage}


\noindent\textit{Current registry lookup labels:} Evoked potential and stimulus artifact; evoked stimulation; evoked recording\par

% Stable term ID: results-already-seen-during-development
\begin{samepage}
\subsection*{Results already seen during development}
\noindent Results that have been inspected and can therefore influence later decisions.\par
\end{samepage}
\begin{samepage}
\noindent The human 90\% result and local motor test windows have this limitation. They can be reported, but their later reuse is not a new independent test.\par\nopagebreak[4]

\end{samepage}


\noindent\textit{Current registry lookup labels:} Development evidence; development holdout; pilot; descriptive local test; inspected test\par

% Stable term ID: reusing-learning-in-another-recording
\begin{samepage}
\subsection*{Reusing learning in another recording}
\noindent Trying to use information learned in one setting to improve performance in another.\par
\end{samepage}
\begin{samepage}
\noindent The comparison must state the new-setting labels and the controls. The reported frozen motor transfer and unadjusted human cross-day result do not establish successful biological or cross-species transfer.\par\nopagebreak[4]

\end{samepage}


\noindent\textit{Current registry lookup labels:} Cross-recording transfer; transfer; cross-day transfer; cross-species transfer\par

% Stable term ID: routing-a-muscle-around-a-modeled-obstacle
\begin{samepage}
\subsection*{Routing a muscle around a modeled obstacle}
\noindent A rule bends a modeled muscle route around an obstacle. One obstacle shape, a torus, is shaped like a doughnut.\par
\end{samepage}
\begin{samepage}
\noindent The correction restricts three elbow wrapping rules to a stated attachment segment. It removes reproduced route errors without deleting a muscle or claiming a new anatomical measurement.\par\nopagebreak[4]

\end{samepage}


\noindent\textit{Current registry lookup labels:} Muscle wrapping and torus; wrap; wrapping route; torus; ELBOW\_ULNA; ELBOW\textbackslash{}\_ULNA; off-torus contact\par

% Stable term ID: saved-model-and-matching-output-check
\begin{samepage}
\subsection*{Saved model and matching-output check}
\noindent A checkpoint saves a model’s state. A parity check asks whether another execution produces the specified matching output.\par
\end{samepage}
\begin{samepage}
\noindent Repeating a result needs the same support, ordering, conversions, data splits, settings and body-model version. Passing software checks does not certify biological validity.\par\nopagebreak[4]

\end{samepage}


\noindent\textit{Current registry lookup labels:} Checkpoint and parity; checkpoint; replay/checkpoint parity; model parity; reproduction\par

% Stable term ID: scalp-electrical-recording
\begin{samepage}
\subsection*{Scalp electrical recording}
\noindent Electrodes on the scalp record electrical potential; the method is called EEG.\par
\end{samepage}
\begin{samepage}
\noindent This is an aggregate electrical signal, not a list of individual cortical spikes. The project’s paired preview follows electrical stimulation rather than voluntary motor training.\par\nopagebreak[4]

\end{samepage}


\noindent\textit{Current registry lookup labels:} Electroencephalography; EEG; scalp potentials\par

% Stable term ID: settings-needed-to-repeat-a-run
\begin{samepage}
\subsection*{Settings needed to repeat a run}
\noindent The numbers and choices fixing one training run’s starting point and procedure.\par
\end{samepage}
\begin{samepage}
\noindent One fixed run can be reproduced. It does not show how results vary across repeated starts, animals or future recording days.\par\nopagebreak[4]

\end{samepage}


\noindent\textit{Current registry lookup labels:} Random seed and configuration; seed; single-seed; fixed protocol\par

% Stable term ID: shared-internal-model
\begin{samepage}
\subsection*{Shared internal model}
\noindent A model can learn an internal pattern and reuse that pattern when making several predictions.\par
\end{samepage}
\begin{samepage}
\noindent The paper calls the pattern a latent representation and the reusable model a backbone or trunk. The tested shared model and frozen-transfer candidates were rejected.\par\nopagebreak[4]

\end{samepage}


\noindent\textit{Current registry lookup labels:} Latent representation and backbone; latent model; shared latent; backbone; trunk; shared parameters\par

% Stable term ID: simple-weighted-prediction-model
\begin{samepage}
\subsection*{Simple weighted prediction model}
\noindent A linear model that combines inputs using fitted weights while discouraging excessively large weights.\par
\end{samepage}
\begin{samepage}
\noindent The formal method is ridge regression. It provides a meaningful comparison: a more complex model must add value rather than merely produce a plausible display.\par\nopagebreak[4]

\end{samepage}


\noindent\textit{Current registry lookup labels:} Ridge regression; ridge; linear control; linear model\par

% Stable term ID: simplicity-principle
\begin{samepage}
\subsection*{Simplicity principle}
\noindent One concept becomes understandable before the next element changes the picture. Separate views or successive reveals can make their relationship easier to follow.\par
\end{samepage}
\noindent The relationship between elements is part of the explanation, rather than something a crowded picture can leave the reader to infer. Simplicity cannot remove scientific qualifications, evidence, source rights or required methods. Actual comprehension and exported readability still need review.\par

% Stable term ID: skin-prick-suit
\subsection*{Skin-prick suit}
\noindent A suit shows proposed interface contacts across the body.\par
\noindent The proposed body-side (peripheral) interface, drawn as a suit.\par
\noindent A label for the proposed peripheral interface in the programme's illustrations, not a general clinical synonym. The name establishes no acquired neural data, registration, tissue access, function, safety, clinical use, species or physical dimensions.\par

\noindent\textit{Current registry lookup labels:} skin prick suit; peripheral suit; peripheral prosthesis; peripheral interface\par

% Stable term ID: small-connector-model
\begin{samepage}
\subsection*{Small connector model}
\noindent A smaller fitted mapping that connects one recording to an existing model or interface.\par
\end{samepage}
\begin{samepage}
\noindent The connector can change while the older model stays fixed. That does not establish matching neurons or useful transfer without new-setting calibration.\par\nopagebreak[4]

\end{samepage}


\noindent\textit{Current registry lookup labels:} Adapter; linear adapter; recording adapter\par

% Stable term ID: software-that-calculates-a-computer-body-s-motion
\begin{samepage}
\subsection*{Software that calculates a computer body’s motion}
\noindent MuJoCo calculates how a modeled body moves under forces and control inputs.\par
\end{samepage}
\begin{samepage}
\noindent It is established software reused here. Its availability does not prove that the chosen body model or mapping matches a particular animal’s physiology.\par\nopagebreak[4]

\end{samepage}


\noindent\textit{Current registry lookup labels:} MuJoCo; dynamics engine\par

% Stable term ID: spinal-strip
\subsection*{Spinal strip}
\noindent A strip shows the proposed interface along the modeled spine.\par
\noindent The proposed spinal interface, drawn as a strip along the modeled spine.\par
\noindent A label for the proposed spinal interface in the programme's illustrations, not a general clinical synonym. The name establishes no acquired neural data, registration, tissue access, function, safety, clinical use, species or physical dimensions. It does not validate a modeled dorsal-surface or tailbud overlay as biological spinal-cord or conus anatomy.\par

\noindent\textit{Current registry lookup labels:} spinal strip; spinal strap; spinal prosthesis; spinal interface\par

% Stable term ID: spine2body
\begin{samepage}
\subsection*{Spine2Body}
\noindent A muscle signal becomes a movement through tendons, joints, contact and load. Measure those steps together, and a failed reach becomes a specific question about activation, coordination or mechanics.\par
\end{samepage}
\noindent The pathway through which spinal output reaches peripheral motor nerves and muscles, producing action through the body's mechanics, while sensory signals return from the body.\par
\noindent Surface EMG measures muscle electrical activity, not peripheral nerve firing or force. Spinal plasticity, momentary reflex modulation and peripheral adaptation remain distinct. Simulated movement is a prediction to compare with independent measurements.\par
\begin{samepage}
\noindent\textit{Organizational level:} whole system; network; signal/feature\par\nopagebreak[4]

\end{samepage}


\noindent\textit{Current registry lookup labels:} spine2body; spine to body; spine-to-body\par

% Stable term ID: stopping-training-under-a-stated-rule
\begin{samepage}
\subsection*{Stopping training under a stated rule}
\noindent Ending training when the chosen development rule says to stop, rather than always using every allowed round.\par
\end{samepage}
\begin{samepage}
\noindent Compared methods have equal round limits but can perform different numbers of actual updates. This is not measured biological learning time.\par\nopagebreak[4]

\end{samepage}


\noindent\textit{Current registry lookup labels:} Early stopping; epoch cap; training update\par

% Stable term ID: structured-task-text-and-supported-instructions
\begin{samepage}
\subsection*{Structured task text and supported instructions}
\noindent JSON stores structured text values. The task compiler converts a supported instruction into a defined program.\par
\end{samepage}
\begin{samepage}
\noindent The compiler chooses available recorded building blocks and limited computer-controller practice. It does not program arbitrary biological skills or supply missing task evidence.\par\nopagebreak[4]

\end{samepage}


\noindent\textit{Current registry lookup labels:} JSON and task compiler; JSON; task JSON; natural-language compiler; supported primitive\par

% Stable term ID: surface-electrical-recording-over-the-spine
\begin{samepage}
\subsection*{Surface electrical recording over the spine}
\noindent Surface electrodes placed over spinal regions record electrical potentials; the method is called ESG.\par
\end{samepage}
\begin{samepage}
\noindent These measurements combine electrical contributions. They are not separately identified spinal-neuron spike trains.\par\nopagebreak[4]

\end{samepage}


\noindent\textit{Current registry lookup labels:} Electrospinography; ESG; surface spinal recording; surface ESG\par

% Stable term ID: system-3-thinking
\begin{samepage}
\subsection*{System 3 thinking}
\noindent System 3 thinking is the partnership; a large movement model is the movement-focused model we aim to build within it.\par
\end{samepage}
\noindent The programme’s human–AI partnership framing in which an external supervisory model helps interpret, train or calibrate task-specific neural/body mappings; the movement-focused model being built is called a large movement model.\par
\noindent Expert footage provides movement examples, but it does not supply the depicted experts’ neural or muscle recordings. An external supervisor would help interpret, train or calibrate task-specific neural/body mappings while the large movement model remains the movement-focused model being built within that partnership. Shaw and Nave use System3 for artificial cognition outside the brain in a different research account. This programme’s task-specific framing establishes no new anatomical brain layer, biological neo-neocortex, third anatomical thinking system or achieved universal nervous-system model.\par



\noindent\textit{Source:} Thinking—Fast, Slow, and Artificial\par
\noindent\textit{Source identifier:} shaw-nave2026\par
\noindent\textit{DOI:} 10.2139/ssrn.6097646\par
\noindent\url{https://papers.ssrn.com/sol3/papers.cfm?abstract_id=6097646}\par
\noindent Prior proposed System3 usage for artificial cognition outside the brain; not an established third anatomical brain system\par
\noindent\textit{Source:} Thinking Fast, Slow, Artificially: AI and Your Brain\par
\noindent\textit{Source identifier:} wharton-authors2026\par
\noindent\href{https://executiveeducation.wharton.upenn.edu/thought-leadership/wharton-at-work/2026/05/thinking-fast-slow-and-artificially/}{\nolinkurl{https://executiveeducation.wharton.upenn.edu/}}\\
\href{https://executiveeducation.wharton.upenn.edu/thought-leadership/wharton-at-work/2026/05/thinking-fast-slow-and-artificially/}{\nolinkurl{thought-leadership/wharton-at-work/2026/05/}}\\
\href{https://executiveeducation.wharton.upenn.edu/thought-leadership/wharton-at-work/2026/05/thinking-fast-slow-and-artificially/}{\nolinkurl{thinking-fast-slow-and-artificially/}}\par
\noindent Primary institutional account of Shaw and Nave System3 research\par


\noindent\textit{Current registry lookup labels:} System three thinking; System3 thinking\par

% Stable term ID: testing-on-an-excluded-day-animal-or-species
\begin{samepage}
\subsection*{Testing on an excluded day, animal or species}
\noindent Checking a model on complete identities that were kept out of the relevant development steps.\par
\end{samepage}
\begin{samepage}
\noindent The paper uses formal names such as leave-session-out or leave-species-out. A good score within one session does not answer these broader tests.\par\nopagebreak[4]

\end{samepage}


\noindent\textit{Current registry lookup labels:} Leave-out generalization; leave-session-out; leave-animal-out; leave-species-out; unseen-organism support\par

% Stable term ID: testing-possible-next-actions-in-the-model
\begin{samepage}
\subsection*{Testing possible next actions in the model}
\noindent Simulating near-future actions before choosing the next command.\par
\end{samepage}
\begin{samepage}
\noindent The copied test states must not change the live model. The native result remains rejected and too slow for real-time use; cached playback is a separate function.\par\nopagebreak[4]

\end{samepage}


\noindent\textit{Current registry lookup labels:} Short-horizon planning and lookahead; horizon controller; lookahead; copied-state rollout; model predictive planning\par

% Stable term ID: the-comparison-needed-to-judge-improvement
\begin{samepage}
\subsection*{The comparison needed to judge improvement}
\noindent A defined comparison that shows whether the proposed method adds useful information or performance.\par
\end{samepage}
\begin{samepage}
\noindent A simple predictor, random fixed model and constant muscle command answer different questions. Producing a plausible animation is not a substitute for beating the relevant comparison.\par\nopagebreak[4]

\end{samepage}


\noindent\textit{Current registry lookup labels:} Baseline and control; simple baseline; tonic baseline; random frozen control; behavior-only predictor\par

% Stable term ID: thought2action
\begin{samepage}
\subsection*{Thought2Action}
\noindent Begin with one intention and one action: reach for the cup. Measure the neural, muscular and mechanical steps between them. The aim is to make that intent-to-action relationship as close to one-to-one as possible within the declared task, person, state, context and tolerated error.\par
\end{samepage}
\noindent A task-conditioned mapping from a participant's stated or inferred intent to a specified action, evaluated through independently measured outcomes.\par
\noindent Intent inferred from EEG is not direct access to private thought. Movement constrains competing explanations; it does not certify an achieved universal bijection or the best possible mapping.\par
\begin{samepage}
\noindent\textit{Organizational level:} whole system; network; signal/feature\par\nopagebreak[4]

\end{samepage}


\noindent\textit{Current registry lookup labels:} thought2action; thought to action; thought-to-action\par

% Stable term ID: thought2world
\begin{samepage}
\subsection*{Thought2World}
\noindent Follow the intention all the way to its consequence, then follow sensation and perception back to the next intention. Measuring the middle makes errors easier to locate and the next experiment easier to define.\par
\end{samepage}
\noindent The composed pathway Thought2Action followed by Action2World: intent is translated into an action, and that action produces an outcome that can be independently checked.\par
\noindent This is a proposed programme pathway, not a claim of unrestricted thought control. It is not silently equivalent to the existing Body2World lookup label, Body2Action or Brain2Motion. External touch and movement tasks alone do not solve internal-body sensing or the full Interoceptomics programme.\par
\begin{samepage}
\noindent\textit{Organizational level:} whole system; network; signal/feature\par\nopagebreak[4]

\end{samepage}


\noindent\textit{Current registry lookup labels:} thought2world; thought to world; thought-to-world\par

% Stable term ID: timing-relationships-between-spikes
\begin{samepage}
\subsection*{Timing relationships between spikes}
\noindent How closely spikes can repeat and how spikes from different recorded features occur together.\par
\end{samepage}
\begin{samepage}
\noindent The project checks predicted average counts. It does not establish these finer timing relationships or a complete account of biological spike noise.\par\nopagebreak[4]

\end{samepage}


\noindent\textit{Current registry lookup labels:} Refractoriness and spike synchrony; refractoriness; joint spike synchrony; biological count-noise model\par

% Stable term ID: touch-brain-area
\begin{samepage}
\subsection*{Touch brain area}
\noindent A brain area that processes body-related sensory information. The source study records from somatosensory cortex, including Area 2.\par
\end{samepage}
\begin{samepage}
\noindent The formal shorthand is S1. The selected model estimates an earlier limb state; it does not recover voluntary motor commands or spinal and nerve recordings.\par\nopagebreak[4]

\end{samepage}


\noindent\textit{Current registry lookup labels:} Somatosensory cortex and Area 2; S1; Area 2; Area2\_Bump; somatosensory\par

% Stable term ID: guide-and-learner
\subsection*{Trainer and trainee}
\noindent The trainer supplies guidance; the trainee’s later control is tested. An experienced trainer helping a novice asks a different question from two novices practicing together.\par
\noindent The proposed comparisons include experienced-to-novice, novice-to-novice, individual practice and matched no-link conditions. Later trainee-alone testing removes trainer commands while retaining any bridge required to express movement.\par

\noindent\textit{Current registry lookup labels:} Tutor and pupil; biological tutor; artificial tutor; live tutoring; tutor withdrawal; experienced-to-novice tutor; novice-to-novice; individual practice; Trainer and learner; biological trainer; artificial trainer; live training; trainer withdrawal; experienced-to-novice trainer; Guide and learner; Guide and trainee\par

% Stable term ID: training-a-fresh-model
\begin{samepage}
\subsection*{Training a fresh model}
\noindent Training a new model for the current recording without importing an older model’s learned main weights.\par
\end{samepage}
\begin{samepage}
\noindent This is the scratch specialist. It performs better than the frozen-transfer candidate in both reported label budgets.\par\nopagebreak[4]

\end{samepage}


\noindent\textit{Current registry lookup labels:} Scratch training; scratch specialist; new-recording scratch specialist\par

% Stable term ID: training-with-neural-samples-without-movement-labels
\begin{samepage}
\subsection*{Training with neural samples without movement labels}
\noindent Using neural samples for a training objective that does not require paired movement labels.\par
\end{samepage}
\begin{samepage}
\noindent These samples still provide information and require computation. The first 5\% label condition permits full training-neural support, so it is not only 28.5 seconds of total observation.\par\nopagebreak[4]

\end{samepage}


\noindent\textit{Current registry lookup labels:} Neural-only training support; neural-only pretraining; unlabeled neural support\par

% Stable term ID: trying-to-move-and-actually-moving
\begin{samepage}
\subsection*{Trying to move and actually moving}
\noindent Trying or being cued to move is different from measuring the movement that occurred.\par
\end{samepage}
\begin{samepage}
\noindent The human labels describe attempted finger movement. They do not independently separate cue from intent or provide executed finger trajectories.\par\nopagebreak[4]

\end{samepage}


\noindent\textit{Current registry lookup labels:} Attempted versus executed movement; attempted finger identity; cued attempted movement; executed kinematics\par

% Stable term ID: typical-error-on-the-original-scale
\begin{samepage}
\subsection*{Typical error on the original scale}
\noindent RMS or RMSE squares the deviations, averages them, and takes a square root so the result has the original unit.\par
\end{samepage}
\begin{samepage}
\noindent An error per joint differs from a combined seven-joint error. An error against a changed command goal also differs from an error against the same recorded reference.\par\nopagebreak[4]

\end{samepage}


\noindent\textit{Current registry lookup labels:} Root mean square error; RMSE; RMS; root-mean-square; tracking error\par

% Stable term ID: unassigned-event-channel
\begin{samepage}
\subsection*{Unassigned event channel}
\noindent A channel collecting events not assigned to the sorted features used by the model.\par
\end{samepage}
\begin{samepage}
\noindent These events are excluded from the 171 sorted features in the second Indy recording. They cannot silently become another identified cell.\par\nopagebreak[4]

\end{samepage}


\noindent\textit{Current registry lookup labels:} Residual or hash cell; residual/hash; hash slot; threshold crossing\par

% Stable term ID: units-for-distance-speed-and-angle
\begin{samepage}
\subsection*{Units for distance, speed and angle}
\noindent Metres, centimetres, millimetres and micrometres measure distance. Metres per second measures speed; degrees and radians measure angle. Radians per second (rad/s) measures the rate of change of angle.\par
\end{samepage}
\begin{samepage}
\noindent One metre is 100 cm or 1000 mm; one mm is 1000 µm. Right-angle and varied-angle maze layouts differ in geometry, not in a validated difficulty multiplier. Unit conversion or a small model error does not prove anatomical or decoded-pose accuracy.\par\nopagebreak[4]

\end{samepage}


\noindent\textit{Current registry lookup labels:} Length, velocity and angle units; m; metre; cm; centimetre; mm; millimetre; µm; micrometre; um; m/s; rad; radian; degrees; °; rad/s; orthogonal; nonorthogonal; varied-angle maze; maze layout\par

% Stable term ID: units-for-pushes-turns-and-electrical-potential
\begin{samepage}
\subsection*{Units for pushes, turns and electrical potential}
\noindent Newtons measure force; newton metres measure turning moment; millivolts measure electrical potential. A dimensionless value has no physical unit.\par
\end{samepage}
\begin{samepage}
\noindent The handle record mixes force and moment units. The rescaled muscle targets are dimensionless despite a stale mV label and cannot be read as force or activation.\par\nopagebreak[4]

\end{samepage}


\noindent\textit{Current registry lookup labels:} Force, moment and voltage units; N; newton; N m; N·m; newton metre; mV; millivolt; dimensionless\par

% Stable term ID: units-for-time-and-sample-rate
\begin{samepage}
\subsection*{Units for time and sample rate}
\noindent Seconds measure time; milliseconds are thousandths of a second. Hertz counts repetitions or samples per second.\par
\end{samepage}
\begin{samepage}
\noindent A 250 Hz source sample interval is 4 ms; a 500 Hz simulation interval is 2 ms. These rates differ from playback pace, biological travel time and time spent computing.\par\nopagebreak[4]

\end{samepage}


\noindent\textit{Current registry lookup labels:} Time and frequency units; s; second; seconds; ms; millisecond; milliseconds; Hz; hertz; sample rate; 250 Hz; 500 Hz\par

% Stable term ID: unwanted-stimulus-signal-in-a-recording
\begin{samepage}
\subsection*{Unwanted stimulus signal in a recording}
\noindent Part of a recording caused by the stimulus or measuring setup rather than the intended response alone.\par
\end{samepage}
\begin{samepage}
\noindent Stimulus artifacts remain in the human previews. Their presence and the evoked task limit interpretation as voluntary motor evidence.\par\nopagebreak[4]

\end{samepage}


\noindent\textit{Current registry lookup labels:} stimulus artifact; artifact\par

% Stable term ID: use-in-a-new-setting-without-labeled-adjustment
\begin{samepage}
\subsection*{Use in a new setting without labeled adjustment}
\noindent Applying a learned model to a new setting without using its labels to adjust the model.\par
\end{samepage}
\begin{samepage}
\noindent The second M1 model uses explicit new-recording labels. Its test score is a result after calibration, not zero-shot success.\par\nopagebreak[4]

\end{samepage}


\noindent\textit{Current registry lookup labels:} Zero-shot transfer; zero-shot\par

% Stable term ID: speed-estimated-from-earlier-positions
\begin{samepage}
\subsection*{Velocity estimated from earlier positions}
\noindent Two recorded fingertip positions can give an estimate of how fast the fingertip moved and in which direction between the samples. The signed components describe velocity; speed is its non-negative magnitude.\par
\end{samepage}
\noindent A position-derived velocity estimate formed from the difference between current and earlier recorded positions divided by their elapsed time; individual components retain sign. Speed is the magnitude of the velocity vector.\par
\begin{samepage}
\noindent In the separate calibrated motor recording, each interval uses the last available fingertip sample and the actual elapsed time to form a velocity target in m/s. This makes the calculation causal with respect to available position samples, not proof that the tracker supplied them without filtering or delay. Upstream filtering and latency remain unestablished; the derived target is not an original, directly supplied velocity trace.\par\nopagebreak[4]

\end{samepage}


\noindent\textit{Current registry lookup labels:} Backward finite difference; backward difference; finite difference; causally derived velocity; 250 Hz position support; Speed estimated from earlier positions\par

% Stable term ID: speed-of-each-model-movement
\begin{samepage}
\subsection*{Velocity of each model movement}
\noindent The model’s stored velocity for its movement degrees of freedom, including direction.\par
\end{samepage}
\begin{samepage}
\noindent The formal variable name is qvel. Its signed generalized velocities depend on the model’s degrees of freedom; speed is a magnitude. Component units and frames depend on the declared coordinate and model. In MuJoCo, quaternion orientations can make qpos and qvel different lengths, so qvel is not an entry-by-entry derivative of qpos.\par\nopagebreak[4]

\end{samepage}


\noindent\textit{Current registry lookup labels:} qvel; generalized velocity; joint speed; Speed of each model movement\par

% Stable term ID: what-a-model-estimated
\begin{samepage}
\subsection*{What a model estimated}
\noindent A model’s estimate of a defined result, such as speed, from the inputs it was given.\par
\end{samepage}
\begin{samepage}
\noindent An estimate remains an estimate even when it matches the test data well. It does not become a new measurement of the spinal cord, nerves or muscles.\par\nopagebreak[4]

\end{samepage}


\noindent\textit{Current registry lookup labels:} Decoded or inferred prediction; decoded; inferred; reconstructed; prediction\par

% Stable term ID: what-an-instrument-recorded
\begin{samepage}
\subsection*{What an instrument recorded}
\noindent A value produced by a stated measuring method, such as an electrode or a movement tracker.\par
\end{samepage}
\begin{samepage}
\noindent Recorded data can contain noise, filtering or artifacts. Calculating muscle length from tracked markers does not turn that length into a direct tissue measurement.\par\nopagebreak[4]

\end{samepage}


\noindent\textit{Current registry lookup labels:} Recorded measurement; measured; observed; recorded signal\par

% Stable term ID: what-the-computer-body-produced
\begin{samepage}
\subsection*{What the computer body produced}
\noindent A value produced while running a computer model, such as a modeled muscle’s activity.\par
\end{samepage}
\begin{samepage}
\noindent Computer-body movement and forces can be checked for numerical errors. They remain separate from measurements of a living body and from evidence that an organism learned.\par\nopagebreak[4]

\end{samepage}


\noindent\textit{Current registry lookup labels:} Simulated quantity; simulated; numerical model state\par

% Stable term ID: what-was-calculated-from-tracking
\begin{samepage}
\subsection*{What was calculated from tracking}
\noindent A value calculated from measurements together with a body model, such as an estimated joint angle.\par
\end{samepage}
\begin{samepage}
\noindent The study’s joint angles and muscle lengths use tracked movement and a model. They are not direct recordings of muscle commands, tissue length or peripheral nerve spikes.\par\nopagebreak[4]

\end{samepage}


\noindent\textit{Current registry lookup labels:} Model-derived quantity; tracking-derived; source-derived; biomechanical covariate\par

% Stable term ID: where-a-value-came-from
\begin{samepage}
\subsection*{Where a value came from}
\noindent The record of a value’s source and the steps used to produce it.\par
\end{samepage}
\begin{samepage}
\noindent For example, a measured hand position, an angle calculated from that position and a computer-generated muscle state have different origins. The software must preserve those differences.\par\nopagebreak[4]

\end{samepage}


\noindent\textit{Current registry lookup labels:} Provenance; lineage; source integrity\par

% Stable term ID: which-recorded-channel-is-which
\begin{samepage}
\subsection*{Which recorded channel is which}
\noindent The labels and order used to identify neural features within one recording.\par
\end{samepage}
\begin{samepage}
\noindent The labels do not identify the same neurons across days. The project did not accept unverified matches between cells in different recordings.\par\nopagebreak[4]

\end{samepage}


\noindent\textit{Current registry lookup labels:} Source-local unit identity; ordered channel identity; cell correspondence; unit order; neuron index\par

% Stable term ID: which-samples-a-value-depends-on
\begin{samepage}
\subsection*{Which samples a value depends on}
\noindent The input history and exact time interval needed to produce a value or target.\par
\end{samepage}
\begin{samepage}
\noindent A single timestamp does not show all of this support. Overlapping histories and invalid samples near data splits need explicit checks.\par\nopagebreak[4]

\end{samepage}

\noindent\textit{Current registry lookup labels:} Temporal support and context; support; history window; context; next-target support; overlapping histories\par

% Stable term ID: whole-body-interface
% Canonical status: "user_requested_proposed_programme_role"
\subsection*{Whole Body Interface}
\noindent The hand is one part of the question. Whole Body Interface names the wider ambition to connect declared body systems, while showing which pathways have actually been observed and which remain missing.\par
\noindent A proposed extension of brain–body interfacing toward specified observations and outputs across the brain, spinal, peripheral and body systems.\par
\noindent Whole body states an ambition, not achieved coverage. An interface exchanges specified observations or outputs; the existing model spanning the body is a related modelling ambition, not the same object. State the recorded systems, missing pathways, units, clocks, frames and output routes. Separate datasets can support qualified learning and transfer, but they do not become simultaneous measurements of one body. Physiological input/output direction, software input/output and ordinary sensory feedback remain distinct; no qualified spinal writer or biological actuation is demonstrated.\par

% Stable term ID: xenotransgenicism
% Canonical status: "user_requested_working_label_scope_unsettled"
\subsection*{XenoTransgenesis}
\noindent Which genetic feature would be studied, in which organism, and what would count as the relevant change? XenoTransgenesis names that proposed question within EndoHumanisation.\par
\noindent A working roadmap name for proposed research involving specified genetic components from another species in a model organism.\par
\begin{samepage}
\noindent XenoTransgenesis is the preferred spelling; XenoTransgenicism remains a historical lookup label. This is not an established replacement name for every form of transgenesis, genome editing or humanisation; use the established term appropriate to any actual manipulation. Neither a genetic change nor its name establishes the intended functional effect or human equivalence. No implemented procedure, qualified organism or developmental schedule is established by this working label.\par\nopagebreak[4]

\end{samepage}


\noindent\textit{Current registry lookup labels:} XenoTransgenicism\par

\begin{samepage}
\subsection*{Composed pathways}
\noindent The proposed intermediary has a positive purpose: understand and measure the processes between intent and action so the task-specific mapping can be made as close to one-to-one as possible. The task, participant, state, context, target and tolerated error must be specified and independently tested. This is an aim for identifiability and control, not an achieved universal bijection or a demonstrated optimum.\par
\end{samepage}
\noindent Form \texorpdfstring{\ensuremath{\rightarrow}}{->} function \texorpdfstring{\ensuremath{\rightarrow}}{->} action and intent \texorpdfstring{\ensuremath{\rightarrow}}{->} action describe connected questions. The following proposed pathways separate the links that would need independent tests before their combined prediction could be trusted.\par
\noindent\textit{Biological pathway:} Brain2Spine \texorpdfstring{\ensuremath{\rightarrow}}{->} Spine2Body; resulting in Brain2Body\par
\noindent Relate brain activity to spinal processing, and spinal/peripheral processing to the body; compose the tested links into an intent-to-body-action pathway.\par
\noindent\textit{Functional pathway:} Thought2Action \texorpdfstring{\ensuremath{\rightarrow}}{->} Action2World; resulting in Thought2World\par
\noindent Relate a declared intention to a measured action, and that action to its measured consequence in the task environment.\par
\noindent\textit{Molecular analogy:} DNA \texorpdfstring{\ensuremath{\rightarrow}}{->} RNA; RNA \texorpdfstring{\ensuremath{\rightarrow}}{->} protein; resulting in DNA \texorpdfstring{\ensuremath{\rightarrow}}{->} protein\par
\noindent The intermediate processes explain how an initial specification reaches its downstream product. This analogy illustrates the need for a measured middle; it does not assert identical biology or a context-free one-to-one relationship.\par
\noindent\textit{System 3 engineering route:} Neural activity \texorpdfstring{\ensuremath{\rightarrow}}{->} Musculoskeletal simulation \texorpdfstring{\ensuremath{\rightarrow}}{->} Programmable model\par
\noindent The proposed engineering route to NeuroTranslation, coordinated by the System3 human–AI supervisory model.\par
\noindent Paired source-bound brain/spinal/peripheral/muscle and sensory/outcome observations constrain mechanics and a task-specific learner. A musculoskeletal simulation represents the biomechanical middle. Model read/write capabilities must be qualified individually.\par
\noindent Intent \texorpdfstring{\ensuremath{\rightarrow}}{->} action \texorpdfstring{\ensuremath{\rightarrow}}{->} outcome returns through sensory feedback to perception and subsequent intent; actual observed/inferred/simulated links remain distinguishable.\par
\noindent Interoceptomics combines interoception and connectomics. Touch, proprioception and external sensorimotor tasks can support this programme; they do not by themselves demonstrate internal-body sensing. Combining the proposed pathways advances a research aim rather than establishing a comprehensive solution.\par
\noindent Brain2Body names the wider programme. Body2World remains a contextual lookup label for Body2Action; Brain2Motion and Body2Action retain their own scopes. These related names do not make their measurement targets equivalent to those of the five composed paths.\par
\noindent No spinal recording or writing interface is implemented or qualified by this proposed composition. Each claimed recording or writing capability needs its own independent evidence.\par


\subsection*{A cross-species graft}

\noindent A chimera has biological contributions from different sources. A graft is added tissue; when it crosses species, it is a xenotransplant. An organoid is an organized tissue model.\par

\noindent Survival, physical connection, activity and useful function are separate questions. The proposed comparisons do not establish a humanized whole brain or an implemented treatment.\par

\begin{samepage}

\noindent Cross-species-graft component only; not a strict synonym for every chimera, graft or organoid, or achieved ExoHumanisation.\par\nopagebreak[4]

\noindent\textit{Additional lookup label:} XenoTransplantation\par

\end{samepage}

\subsection*{Proposed names and explanatory comparisons}

% Proposal ID: neo-neocortex-inspiration

\subsubsection*{Neo-neocortex}

\noindent A metaphor for an external machine contribution to human cognition; no anatomical addition is claimed.\par

% Proposal ID: alpha-transcriptome-capability-placeholder

\subsubsection*{AlphaTranscriptome}

\noindent AlphaTranscriptome is a proposed name for a desired RNA/cell-state prediction capability. It is not an adopted general-purpose model. AlphaGenome already predicts some expression and splicing~\citep{alphagenomeRetained}; other RNA/cell-state models have bounded domains. These precedents leave whole-cell programmability to be independently demonstrated.\par

% Proposal ID: alpha-cell-supplied-passage-concept

\subsubsection*{AlphaCell}

\noindent AlphaCell names a proposed integration of cellular modelling layers. The idea establishes no implementation, novelty or entitlement to a product.\par

\noindent Connetype, Interotype, Enactotype? These tentative profile names leave a question open. No formal empirical definition or repeatability threshold has been fixed for any of the three.\par

% Proposal ID: connetype-individual-profile-proposal

\subsubsection*{Connetype}

\noindent A connection profile would describe the selected observations for one individual. Its anatomical coverage, resolution, method and uncertainty must remain visible. The profile does not identify that person’s unique function, thoughts or a fixed biological type.\par

\noindent An individual’s specified neural-connection profile under a declared measurement and representation.\par


\noindent The tentative name is Connetype. Connectype remains an unapproved editorial alternative; it is neither a lookup alias nor an adopted spelling.\par

% Proposal ID: interotype-individual-profile-proposal

\subsubsection*{Interotype}

\noindent A heartbeat-related observation can vary with state and measurement conditions. A proposed individual profile must state what was observed and how repeatability would be assessed; it cannot treat touch or proprioception as interchangeable with interoception.\par

\noindent An individual’s profile of specified internal-body sensing observations under declared conditions.\par


\noindent Interotype is a tentative name, rather than an established scientific type.\par

% Proposal ID: enactotype-individual-profile-proposal

\subsubsection*{Enactotype}

\noindent A profile of selected actions is bounded by the tasks, observations and conditions used to construct it. Repeatability, state dependence and generalization remain empirical questions; a stable innate action type or inferred unique intent is not established.\par

\noindent An individual’s profile of specified actions and action–outcome observations in declared tasks.\par


\noindent Enactotype is a tentative name, rather than an established scientific type.\par

% Proposal ID: form-function-feature-explanatory-lens

\subsubsection*{Form, Function, Feature}

\noindent Start with a structure, ask what it does in the stated setting, then name the observable feature being compared. A hand’s form, its grasping function and a measured property of a particular grasp answer related but different questions.\par

\noindent Feature must identify the actual observed or defined trait in its context; it is not a new universal third biological layer or a substitute for function or action. The existing Form \(\rightarrow\) function \(\rightarrow\) action account and intent–action distinction remain unchanged. Structure constrains function; context, history, selection and co-option prevent a guaranteed one-to-one or teleological chain.\par

% Proposal ID: semiosis-sensorimotor-explanatory-analogy

\subsubsection*{Semiosis as a sensorimotor analogy}

\noindent A drawing can resemble a hand, a pointing finger can indicate an object, and an instruction can follow a learned convention. Such signs can guide action through different relations; they do not assign perception, sensation and action to three exclusive boxes.\par

\noindent Icon, index and symbol describe sign–object relations, while sensation, perception and action name processes. A sign can combine relations. No one-to-one biological correspondence, neural code, private-experience identity or strict isomorphism is established. Outside-in and inside-out perspectives must state whether they concern physiology, modelling or interpretation.\par

% Proposal ID: agent-vs-neurology-stage5-proposal

\subsubsection*{Agent vs Neurology}

\noindent The roadmap ends by asking what an agent could contribute to a defined neurological problem. Agent vs Neurology retains that proposed destination; the specific problem, task and comparison still have to be chosen.\par

\noindent An unresolved Stage 5 research question, not a selected disorder, intervention, diagnostic system or clinical trial.\par

\noindent No specific disease, patient population, endpoint, treatment mechanism, universal cure or successful clinical benefit is inferred. Any future study must define those matters and its ordinary feedback or separately qualified physical interface under the original scientific, participant, rights, safety and independent gates. A prior software, metadata or model result cannot supply this evidence.\par

\noindent Agent vs Neurology names Stage 5 of the proposed roadmap.\par

\nocite{*}\bibliographystyle{unsrtnat}\bibliography{references,myohand-numerical-citations-common,myohand-numerical-citations-brain2body-only,training-record-citations,../../common/public-reference-citations}
\clearpage
\appendix
% Source-only shared appendix. No data, apparatus, reference-series or output acceptance.
\section{Reference Comparison Protocol: Unfilled Record}
\label{app:reference-comparison-protocol}

A moving hand gives the instruments something in common to answer to. Scalp
EEG would observe electrical activity at declared scalp sites; surface EMG
would observe muscle electrical activity at declared arm sites. Neither tells
us, by itself, where the hand went. The proposed LEFT and RIGHT tripod views
would provide complementary movement observations. Each acquisition needs its
own identifiable series of real references, followed by an independent
comparison of the assembly. Two streams can agree because they share the same
mistaken calibration.

\textbf{This is an unfilled protocol and record.} No observations, admitted
reference series, measured errors or acceptance decisions have been entered.
The apparatus and its anatomical and functional claims remain unvalidated.
The older above/below maze illustration retains its separate schematic
context, perpendicular-to-floor geometry and material-path questions. Its
directions are not synonyms for the current LEFT/RIGHT tripod arrangement.
Actual positions, viewing paths, coverage and transmission would need evidence.

\subsection*{Identify the body and the observation}

A photograph can constrain a visible proportion; it cannot supply an EEG
channel or a calibrated trajectory. Record each reference's provider, stable
identifier, version or acquisition context, checksum if materialized, license,
access terms and supporting claim. Identify the organism, developmental stage,
participant or specimen, session, body part, task and view or time interval.
Record physical dimensions and units, posture, support and contact, so that the
comparison concerns the same body and movement. Keep those dimensions separate
from computational detail and from the rule used to average errors. Existing
datasets, atlases and photographs retain their separate study identities;
none is silently admitted for this apparatus.

EEG needs sites, reference, channel order and electrical units. Surface EMG
needs named muscles, sleeve and contact identities, reference and electrical
units; it supplies neither incoming sensory-nerve traffic nor direct force.
Each tripod view needs its selected modality and camera or ranging geometry.
Optical triangulation and LiDAR ranging measure different things. Preserve raw
observations, derived features and model estimates separately. These recordings
do not identify hidden spinal activity, individual neurons or a unique brain state.

\subsection*{Make time and position comparable}

A reach can follow the measured path and still arrive late. Record source
timestamps, clock identities, cadence and units, synchronization events,
offsets, drift, dropped samples and interpolation. Retain calibration
observations, frame transformations, their uncertainty and checks against an
independent reference. Report offline processing and end-to-end latency
explicitly. An alignment error must not disappear inside a movement score,
and no real-time capability follows from this protocol.

Compare scalp EEG, each arm's surface EMG, each tripod acquisition, and the
time/frame calibration and any actual material path against their respective
series of real references. Keep overlap, occlusion, dropout and failed views
alongside anatomy, proportion, support and contact. Preserve both motion
observations and the joined map. A missing finger remains missing unless its
replacement is labelled as an estimate. The assembled map then needs its own
independent movement comparison; agreement among its inputs cannot establish
the accuracy of their common reconstruction.

\subsection*{Test what the assembly claims}

Independently measured movement can test a virtual body's path. Contact and
force claims need appropriate measurements of contact and force; a lifelike
image cannot establish valid mechanics. A learner's prediction needs a
withheld measured target, declared input availability, baselines, negative
controls, independent held-out participants, sessions, tasks or species where
applicable, and leakage checks. Keep the reserved recordings closed and the
original tolerances unchanged. This unfilled record contains no new fitting, biological teaching, participant acquisition, stimulation, physical intervention or model-quality result. No programme spinal intervention is reported here.

For each component and the assembled apparatus, retain matched reference
views or time intervals, the target observable, result, error with units,
repeatability, uncertainty and failures. Declare the applicable original
acceptance rule, its source, domain, comparator and tolerances before evaluation.
The programme's unchanged real-reference contract and every stronger original
scientific, anatomical, functional, physical, held-out, rights, visual, runtime
and export criterion still govern their respective claims. No universal
numeric threshold is introduced or relaxed. Biological function and clinical
readiness need evidence beyond a matched motion record.

\subsection*{Unfilled comparison record}

Use one record per reference and comparison, retaining a series for every
component and for the assembled map. All values below are \texttt{null};
reference-series lists are empty and every acceptance state is \texttt{HOLD}.

\begin{description}
\item[Source and rights.] Record ID, component or assembly, provider, source
identifier, version or acquisition context, local checksum if materialized,
license, access and supporting claim.
\item[Body and task.] Organism, developmental stage, participant or specimen,
session, body part, task, view or time interval, physical dimensions and units,
posture, support, contact, actual sensor geometry and material path.
\item[Acquisition.] Modality, measured quantity and units, sites, montage,
reference, channel and contact identities, source frame and clock, timestamps,
cadence, synchronization events, offsets, drift, calibration observations and
transforms, calibration uncertainty, interpolation, missingness, occlusion,
offline processing and end-to-end latency.
\item[Independent comparison.] Component-series and assembly-series identifiers,
independent reference and its acquisition, aligned interval, raw/derived/simulated
provenance, input availability, held-out split and custodian, baseline, negative
controls, leakage checks and target observable.
\item[Result and decision.] Result, error and units, repeatability, uncertainty,
failure cases, applicable original acceptance rule and source, domain,
comparator, tolerances, independent reviewer, evidence and verdict.
\end{description}

No measured series or scientific validation is supplied by this empty record.
Rights, physical, clinical, runtime, native, render, export and publication
assessment remains separate and open. No reference-series dataset, qualified apparatus or validated model release is supplied by this empty record. The public paper/manuscript-source draft is a different object.

\section{Computational reproducibility and software availability}
\label{app:reproducibility-availability}

The studies retain separate recording identities, targets, scalers, software and comparison rules. The original LINK study's completed runs are distinct from the later same-stream TRAIN diagnostic. The latter's literal metrics, denominators, code-version declarations and protocol identities are in \href{training-diagnostic-summary.json}{the scientific diagnostic summary}; they are recorded evidence, not a new model execution. Original rejected results, report corrections and protected data partitions remain part of the record.

\subsection{Software-availability boundary}
The preceding public R9 snapshot omitted \texttt{research-evidence.json}, despite local Source copies and links. Its public Source was therefore non-standalone. The accompanying source candidate stages the current aid and its scientific-summary dependency as regular files with an exact relative-path allowlist; this local preparation is not an uploaded or verified public repair.

The historical online package moved opaque resources from its HTML entry file, reducing that file from 113,298,857 to 34,295,597 bytes. Its complete directory was 93,831,812 bytes and still required the assets. The distinct original-source static offline export was 113,581,899 bytes; API-dependent research views were unavailable offline. The candidate browser comparison was unavailable, so the packaging result supplies no newer browser qualification or complete interface acceptance. File size does not measure loading time, interaction latency, memory, crashes or GPU benefit.

\subsection{Unperformed physical interfaces and missing clinical support}
The recording-and-virtual-modelling proposals report no programme spinal intervention, physical sensory writer or biological actuator. Their device, participant, safety, clinical and cost qualifications remain missing. The earlier source notes refer to an inaccessible clinical-methods source without a resolved bibliographic identity; no assessed recorder-versus-implanted-writer methods comparison is available from it. That exact missing source is an unresolved citation node, not evidence for an interface. All original read-only, admission and physical-intervention holds remain in the unchanged owner/source records.

\subsection{Budget preparation remains incomplete}
The lean, serious-prototype and research-lab tiers ask different experimental questions. None has a complete dated cost estimate. Their preparation record requires equipment, calibration, accessories, storage, integration, safety/consent work, independent references and contingency; USD planning figures, dated vendor evidence and a dated conversion basis; and a distinction between listed prices, engineering allowances and quotations. These are unfulfilled preparation requirements, not reported prices, purchase authority or evidence of a complete neural reconstruction.

\subsection{MyoHand computing-resource record}
The following retained statement records the earlier study and its source-edition scope; it is not a newly executed calculation.

The run used CPU dynamics. Readback showed a one-thread BLAS pool and dynamics pool size zero; transient compiler parallelism was unmeasured. Its 263 observed loaded-file identities provide bounded identity evidence, not complete system closure. Resource and performance qualification remain unresolved. The reported process high-water RSS of 5,187,464 KiB can retain pre-exec history and cannot be assigned to this model or establish a model-only 3 GiB bound. Tool/native timing near 448 ms and in-probe timing near 491.927 ms have unreconciled boundaries. There is no accepted end-to-end speedup, useful complete latency or full CPU/GPU partition.

\end{document}
